% L’industrie — Géographie Tle D — Cours signé OnBuch+
% Programme officiel MINESEC. Compiler avec : tectonic N08-industrie.tex
\newcommand{\DOCMATIERE}{Géographie}
\newcommand{\DOCNIVEAU}{Tle D}
\newcommand{\DOCMODULE}{Module 1 : Le Cameroun}
\newcommand{\DOCLECON}{N08}
\newcommand{\DOCTITRE}{L’industrie}
% OnBuch+ — préambule « cours signés » ENRICHI (compilé avec tectonic / XeLaTeX)
% Système visuel « Pop » d'OnBuch+ : fond crème (jamais blanc), bordures encre
% épaisses, ombres « dures » décalées, titres Archivo Black en capitales.
% Version MATHÉMATIQUES : graphes (pgfplots/tikz), boîtes pédagogiques
% (définition, propriété, méthode, attention, le savais-tu, exercice résolu,
% exercice, TP, bilan), page de garde et signature OnBuch+ ancrée en bas.
%
% Macros attendues dans main.tex AVANT \input{preamble} :
%   \DOCMATIERE  \DOCNIVEAU  \DOCMODULE  \DOCLECON (numéro)  \DOCTITRE
\documentclass[11pt]{article}

\usepackage{fontspec}
\usepackage[french]{babel}
\usepackage[a4paper,margin=2cm,top=3.4cm,bottom=2.8cm,headheight=1.4cm,footskip=1.3cm]{geometry}
\usepackage{amsmath,amssymb,amsfonts}
\usepackage{mathtools} % \xmapsto, \coloneqq... souvent utilisés par les modèles
\usepackage{cancel} % simplifications barrées (bilans de forces)
% \abs{z} (module, valeur absolue) et \norm{u} (norme), utilisés par les modèles
\providecommand{\abs}{}\let\abs\relax\DeclarePairedDelimiter{\abs}{\lvert}{\rvert}
\providecommand{\norm}{}\let\norm\relax\DeclarePairedDelimiter{\norm}{\lVert}{\rVert}
\usepackage[table]{xcolor} % [table] : fournit aussi \rowcolors (alternance de lignes)
\usepackage{pagecolor}
\usepackage{tikz}
\usetikzlibrary{arrows.meta,positioning,calc,shapes.geometric,shapes.misc,shapes.arrows,decorations.pathmorphing,decorations.pathreplacing,decorations.markings,patterns,shadows,mindmap,trees,fit,backgrounds,matrix,intersections,angles,quotes}
\usepackage{pgfplots}
\usepackage[version=4]{mhchem} % équations nucléaires \ce{^{235}_{92}U}
\usepackage[european,siunitx]{circuitikz} % schémas électriques
\pgfplotsset{compat=1.18}
\usepgfplotslibrary{fillbetween,groupplots} % aires entre courbes (calcul intégral)
\usepackage{enumitem}
\usepackage{fancyhdr}
% [most] charge toutes les bibliothèques tcolorbox (listings, minted, raster,
% documentation...) et gonfle inutilement la mémoire TeX sur un document long
% (~300 boîtes) : on ne charge que ce qui est réellement utilisé.
\usepackage[skins,breakable]{tcolorbox}
\usepackage{graphicx}
\graphicspath{{/home/user/t-t/geo-onbuch/images/}}
\usepackage{colortbl}
\usepackage{booktabs}
\usepackage{tabularx}
\usepackage{multirow}
\usepackage{diagbox} % case d'en-tête diagonale des tableaux à double entrée
% Colonnes extensibles centrée / gauche / droite (C, L, R), souvent utilisées par les modèles
\newcolumntype{C}{>{\centering\arraybackslash}X}
\newcolumntype{L}{>{\raggedright\arraybackslash}X}
\newcolumntype{R}{>{\raggedleft\arraybackslash}X}
\usepackage{array}
\usepackage{multicol}
\usepackage{siunitx}
% parse-numbers=false : accepte aussi \SI{2,5 \times 10^{-3}}{...}
\sisetup{locale=FR,output-decimal-marker={,},group-minimum-digits=4,parse-numbers=false}
\DeclareSIUnit\ohm{\ensuremath{\symup{\Omega}}} % U+2126 absent de la police
\DeclareSIUnit\division{div} % réglages d'oscilloscope (ms/div, V/div)
\DeclareSIUnit\tr{tr} % tours (vitesse de rotation, ex. \SI{1500}{\tr\per\minute})
\DeclareSIUnit\molar{M} % concentration molaire (ex. \SI{3}{\molar}, \SI{1}{\micro\molar})
\DeclareSIUnit\an{an} % année (le modèle confond parfois avec \year, un compteur TeX) — ex. \SI{5}{\kilo\meter\per\an}
\DeclareSIUnit\FCFA{FCFA} % franc CFA (ex. \SI{500}{\FCFA\per\kilo\gram})
\DeclareSIUnit\degreeBrix{\textdegree Brix} % teneur en sucre d'un sirop/jus (réfractométrie)
\DeclareSIUnit\month{mois} % le modèle utilise parfois \month, un compteur TeX, comme unité de durée
\DeclareSIUnit\microliter{\micro\liter} % le modèle écrit parfois \microliter en un seul token
\DeclareSIUnit\million{millions} % dénombrement (ex. \SI{15}{\million\per\mL} spermatozoïdes)
\usepackage{qrcode}
% \degree : commande fréquemment utilisée par les modèles pour un angle ou une
% température en dehors de \SI{}{} (non fournie par les paquets chargés).
\providecommand{\degree}{^{\circ}}
% \qed (fin de démonstration), utilisé par les modèles sans amsthm
\providecommand{\qed}{\hfill$\square$}
% \barre : double barre des valeurs interdites dans un tableau de variations (tkz-tab)
\providecommand{\barre}{\ensuremath{\|}}
% environnement proof (amsthm non chargé) utilisé par les modèles
\ifdefined\proof\else\newenvironment{proof}[1][Démonstration]{\par\noindent\textit{#1.}\ }{\qed\par}\fi
\usepackage{lastpage}
\usepackage{hyperref}
\hypersetup{hidelinks}

% ============================================================
% Palette « Pop » OnBuch+ — mêmes hex que la classe P (plus_ui.dart)
% ============================================================
\definecolor{poporange}{HTML}{F59321}
\definecolor{popdark}{HTML}{E07A0C}
\definecolor{popcream}{HTML}{FFF6E8}
\definecolor{popink}{HTML}{1C1714}
\definecolor{poppurple}{HTML}{7A5AE0}
\definecolor{popgreen}{HTML}{2E9E6A}
\definecolor{poppink}{HTML}{E0457B}
\definecolor{popblue}{HTML}{2F7DE1}
\definecolor{popgold}{HTML}{C98A12}
\definecolor{popmuted}{HTML}{8A7F76}
\definecolor{popline}{HTML}{EADFD2}
\definecolor{popgray}{HTML}{8A7F76}
\colorlet{popgrey}{popgray}
% Alias tolérants (le modèle nomme parfois les couleurs différemment)
\colorlet{popwhite}{white}
\colorlet{popblack}{popink}
\colorlet{popred}{poppink}
\colorlet{popyellow}{popgold}
% Teintes claires (fonds de tableaux/encadrés) — le modèle nomme parfois
% les couleurs avec un suffixe L (ex. popblueL) sans les avoir définies.
\colorlet{poporangeL}{poporange!20}
\colorlet{popdarkL}{popdark!20}
\colorlet{poppurpleL}{poppurple!20}
\colorlet{popgreenL}{popgreen!20}
\colorlet{poppinkL}{poppink!20}
\colorlet{popblueL}{popblue!20}
\colorlet{popgoldL}{popgold!20}
\colorlet{popredL}{popred!20}
\colorlet{popgrayL}{popgray!20}
% Teintes claires pour fonds de tableaux / zones
\colorlet{popblueL}{popblue!10!white}
\colorlet{poporangeL}{poporange!14!white}
\colorlet{popgreenL}{popgreen!12!white}
\colorlet{poppurpleL}{poppurple!12!white}
\colorlet{poppinkL}{poppink!12!white}
\colorlet{popgoldL}{popgold!14!white}

\pagecolor{popcream}

% ============================================================
% Typographie
% ============================================================
\defaultfontfeatures{Ligatures=TeX}
\setmainfont{PlusJakartaSans-Medium.ttf}[Path=../../../fonts/,BoldFont=PlusJakartaSans-Bold.ttf]
% Maths en sans-serif (Fira Math) avec chiffres et lettres droites de Plus
% Jakarta, pour que les formules se fondent dans le texte.
\usepackage[math-style=ISO,bold-style=ISO]{unicode-math}
\setmathfont{FiraMath-Regular.otf}[Path=../../../fonts/]
\setmathfont{PlusJakartaSans-Medium.ttf}[Path=../../../fonts/,range={up/num,up/{Latin,latin},\mathup}]
\usepackage{icomma} % virgule décimale sans espace en mode math
% Caractères mathématiques Unicode absents des polices de texte (Plus Jakarta,
% Archivo Black) : rendus par leur équivalent mathématique, en texte comme en
% formule — sinon ils s'affichent en carré vide (ex. « Arithmétique dans ℤ »).
\usepackage{newunicodechar}
\newunicodechar{ℝ}{\ensuremath{\mathbb{R}}}
\newunicodechar{ℤ}{\ensuremath{\mathbb{Z}}}
\newunicodechar{ℂ}{\ensuremath{\mathbb{C}}}
\newunicodechar{ℕ}{\ensuremath{\mathbb{N}}}
\newunicodechar{ℚ}{\ensuremath{\mathbb{Q}}}
\newunicodechar{𝔻}{\ensuremath{\mathbb{D}}}
\newunicodechar{α}{\ensuremath{\alpha}}
\newunicodechar{β}{\ensuremath{\beta}}
\newunicodechar{γ}{\ensuremath{\gamma}}
\newunicodechar{δ}{\ensuremath{\delta}}
\newunicodechar{ε}{\ensuremath{\varepsilon}}
\newunicodechar{θ}{\ensuremath{\theta}}
\newunicodechar{λ}{\ensuremath{\lambda}}
\newunicodechar{μ}{\ensuremath{\mu}}
\newunicodechar{σ}{\ensuremath{\sigma}}
\newunicodechar{τ}{\ensuremath{\tau}}
\newunicodechar{φ}{\ensuremath{\varphi}}
\newunicodechar{ψ}{\ensuremath{\psi}}
\newunicodechar{ω}{\ensuremath{\omega}}
\newunicodechar{Σ}{\ensuremath{\Sigma}}
% majuscules produites par \MakeUppercase dans les titres (α-aminés -> Α)
\newunicodechar{Α}{\ensuremath{\alpha}}
\newunicodechar{Β}{\ensuremath{\beta}}
\newunicodechar{Γ}{\ensuremath{\Gamma}}
\newunicodechar{Δ}{\ensuremath{\Delta}}
\newunicodechar{Ε}{\ensuremath{\varepsilon}}
\newunicodechar{Θ}{\ensuremath{\Theta}}
\newunicodechar{Λ}{\ensuremath{\Lambda}}
\newunicodechar{Μ}{\ensuremath{\mu}}
\newunicodechar{Τ}{\ensuremath{\tau}}
\newunicodechar{Φ}{\ensuremath{\Phi}}
\newunicodechar{Ψ}{\ensuremath{\Psi}}
\newunicodechar{Ω}{\ensuremath{\Omega}}
\newunicodechar{↦}{\ensuremath{\mapsto}}
\newunicodechar{∈}{\ensuremath{\in}}
\newunicodechar{∉}{\ensuremath{\notin}}
\newunicodechar{∧}{\ensuremath{\wedge}}
\newunicodechar{⇔}{\ensuremath{\Leftrightarrow}}
\newunicodechar{⟺}{\ensuremath{\Longleftrightarrow}}
\newunicodechar{⇒}{\ensuremath{\Rightarrow}}
\newunicodechar{⟹}{\ensuremath{\Longrightarrow}}
\newunicodechar{∘}{\ensuremath{\circ}}
\newunicodechar{∪}{\ensuremath{\cup}}
\newunicodechar{∩}{\ensuremath{\cap}}
\newunicodechar{≡}{\ensuremath{\equiv}}
\newunicodechar{⊂}{\ensuremath{\subset}}
\newunicodechar{⊥}{\ensuremath{\perp}}
\newunicodechar{∃}{\ensuremath{\exists}}
\newunicodechar{∀}{\ensuremath{\forall}}
\newunicodechar{∛}{\ensuremath{\sqrt[3]{\,}}}
\newunicodechar{∜}{\ensuremath{\sqrt[4]{\,}}}
\newunicodechar{⃗}{\ensuremath{{}^{\to}}}
\newunicodechar{ⁿ}{\ifmmode^{n}\else\textsuperscript{\ensuremath{n}}\fi}
\newunicodechar{ˣ}{\ifmmode^{x}\else\textsuperscript{\ensuremath{x}}\fi}
\newunicodechar{ᵇ}{\ifmmode^{b}\else\textsuperscript{\ensuremath{b}}\fi}
\newunicodechar{ʳ}{\ifmmode^{r}\else\textsuperscript{\ensuremath{r}}\fi}
\newunicodechar{ᵘ}{\ifmmode^{u}\else\textsuperscript{\ensuremath{u}}\fi}
\newunicodechar{ʸ}{\ifmmode^{y}\else\textsuperscript{\ensuremath{y}}\fi}
\newunicodechar{ᵃ}{\ifmmode^{a}\else\textsuperscript{\ensuremath{a}}\fi}
\newunicodechar{ᵉ}{\ifmmode^{e}\else\textsuperscript{\ensuremath{e}}\fi}
\newunicodechar{ᵏ}{\ifmmode^{k}\else\textsuperscript{\ensuremath{k}}\fi}
\newunicodechar{ᵖ}{\ifmmode^{p}\else\textsuperscript{\ensuremath{p}}\fi}
\newunicodechar{ᴺ}{\ifmmode^{N}\else\textsuperscript{\ensuremath{N}}\fi}
\newunicodechar{ᶜ}{\ifmmode^{c}\else\textsuperscript{\ensuremath{c}}\fi}
\newunicodechar{ʰ}{\ifmmode^{h}\else\textsuperscript{\ensuremath{h}}\fi}
\newunicodechar{ᵠ}{\ifmmode^{q}\else\textsuperscript{\ensuremath{q}}\fi}
\newunicodechar{ₙ}{\ifmmode_{n}\else\textsubscript{\ensuremath{n}}\fi}
\newunicodechar{ₐ}{\ifmmode_{a}\else\textsubscript{\ensuremath{a}}\fi}
\newunicodechar{ᵢ}{\ifmmode_{i}\else\textsubscript{\ensuremath{i}}\fi}
\newunicodechar{ᵣ}{\ifmmode_{r}\else\textsubscript{\ensuremath{r}}\fi}
\newunicodechar{ₖ}{\ifmmode_{k}\else\textsubscript{\ensuremath{k}}\fi}
\newunicodechar{ᵦ}{\ifmmode_{\beta}\else\textsubscript{\ensuremath{\beta}}\fi}
\newunicodechar{ₓ}{\ifmmode_{x}\else\textsubscript{\ensuremath{x}}\fi}
\newfontfamily\popdisplay{ArchivoBlack-Regular.ttf}[Path=../../../fonts/]
\newfontfamily\poplogo{SpaceGrotesk-Bold.ttf}[Path=../../../fonts/]
\newfontfamily\popsemi{PlusJakartaSans-SemiBold.ttf}[Path=../../../fonts/]
\newfontfamily\popxbold{PlusJakartaSans-ExtraBold.ttf}[Path=../../../fonts/]
\newfontfamily\popxboldls{PlusJakartaSans-ExtraBold.ttf}[Path=../../../fonts/,LetterSpace=3.0]

\setlist[enumerate]{leftmargin=1.7em,itemsep=5pt,topsep=4pt}
\setlist[itemize]{leftmargin=1.7em,itemsep=4pt,topsep=4pt,label={\color{poporange}\textbullet}}
\setlength{\parindent}{0pt}
\setlength{\parskip}{5pt}
\renewcommand{\arraystretch}{1.25}

% pgfplots : style Pop par défaut
\pgfplotsset{
  every axis/.append style={axis line style={popink,line width=1pt},tick style={popink},
    grid style={popline,line width=0.5pt},label style={font=\small},tick label style={font=\footnotesize}},
  popcurve/.style={poporange,line width=1.8pt,no markers,smooth},
}
% Même style disponible dans un \draw TikZ simple (hors pgfplots)
\tikzset{popcurve/.style={poporange,line width=1.8pt}}
\tikzset{popgrid/.style={popline,very thin}}
\colorlet{popcurve}{poporange} % le nom du style est parfois utilisé comme couleur

% ============================================================
% Pill / squiggle
% ============================================================
\newcommand{\poppill}[3]{%
  \tikz[baseline=-0.55ex]\node[draw=popink,line width=1.2pt,rounded corners=2.6pt,%
    fill=#2,inner xsep=6.5pt,inner ysep=3pt]{%
    \popxboldls\fontsize{7}{7}\selectfont\color{#3}\MakeUppercase{#1}};%
}
\newcommand{\popsquiggle}[1]{%
  \tikz[baseline=-0.4ex]{
    \draw[#1,line width=2.6pt,line cap=round]
      plot [smooth] coordinates {(0,0.09) (0.34,0.01) (0.68,0.21) (1.02,0.01) (1.36,0.10)};
  }%
}
% Mot-clé surligné (marqueur orange sous le texte)
\newcommand{\cle}[1]{{\popxbold\color{popdark}#1}}

% ============================================================
% En-tête / pied
% ============================================================
\pagestyle{fancy}
\fancyhf{}
\renewcommand{\headrulewidth}{0pt}
\renewcommand{\footrulewidth}{0pt}
\lhead{\raisebox{-0.15ex}{\poplogo\fontsize{13}{13}\selectfont\color{popink}OnBuch\color{poporange}+}}
\rhead{\poppill{\DOCMATIERE\ \textbullet\ \DOCNIVEAU\ \textbullet\ Leçon \DOCLECON}{white}{popink}}
\lfoot{\popxbold\fontsize{7.6}{7.6}\selectfont\color{popmuted}\MakeUppercase{Cours signé OnBuch+}\ \textbullet\ {\popsemi\DOCTITRE}}
\rfoot{\tikz[baseline=-0.6ex]\node[rounded corners=8.5pt,fill=popink,minimum height=17pt,minimum width=17pt,inner xsep=5pt,inner sep=0]{\popxbold\fontsize{8}{8}\selectfont\color{popcream}\thepage\,/\,\pageref*{LastPage}};}

% ============================================================
% Titres
% ============================================================
\newcommand{\chaptitre}[2]{%
  \begin{tcolorbox}[enhanced,colback=poporange,colframe=popink,boxrule=2.6pt,arc=11pt,
      drop shadow=popink,left=15pt,right=15pt,top=12pt,bottom=11pt,before skip=3pt,after skip=18pt]
    \poppill{#2}{popink}{popcream}\par\vspace{9pt}
    {\raggedright\hyphenpenalty=10000\exhyphenpenalty=10000\popdisplay\fontsize{22}{24}\selectfont\color{popcream}\MakeUppercase{#1}\par}
  \end{tcolorbox}
}
% Section numérotée : pastille orange avec le numéro + titre Archivo
\newcounter{popsec}
\newcommand{\coursec}[1]{%
  \stepcounter{popsec}\setcounter{popsub}{0}%
  \par\vspace{16pt}\noindent
  \begin{minipage}{\linewidth}
  \tikz[baseline=-0.7ex]\node[draw=popink,line width=1.6pt,fill=poporange,rounded corners=4pt,minimum size=22pt,inner sep=0]{\popdisplay\fontsize{12}{12}\selectfont\color{popcream}\thepopsec};%
  \hspace{8pt}{\popdisplay\fontsize{15}{17}\selectfont\color{popink}\MakeUppercase{#1}}\par
  \vspace{4pt}\noindent{\color{poporange}\rule{36pt}{3.6pt}}
  \end{minipage}\par\nopagebreak\vspace{8pt}
}
\newcounter{popsub}
\newcommand{\courssub}[1]{%
  \stepcounter{popsub}%
  \par\vspace{9pt}\noindent{\popxbold\fontsize{11.5}{13}\selectfont\color{popink}{\color{poporange}\thepopsec.\thepopsub}\hspace{6pt}#1}\par\nopagebreak\vspace{4pt}
}

% ============================================================
% Boîtes pédagogiques — même recette : fond blanc, bordure épaisse colorée,
% ombre dure encre, badge plein. Titre optionnel [#1].
% ============================================================
\tcbset{popbase/.style={breakable,enhanced,colback=white,boxrule=2.3pt,arc=9pt,
  shadow={3pt}{-3pt}{0pt}{popink},left=14pt,right=14pt,top=11pt,bottom=11pt,before skip=11pt,after skip=15pt,
  fontupper=\popsemi\fontsize{10.6}{14.5}\selectfont\color{popink},
  fontlower=\popsemi\fontsize{10.6}{14.5}\selectfont\color{popink}}}
% Coche dessinée (le glyphe ✓ n'existe pas dans les polices du document)
\newcommand{\popcheck}[1]{\tikz[baseline=-0.5ex]\draw[#1,line width=1.6pt,line cap=round,line join=round](-0.13,0.0)--(-0.04,-0.09)--(0.13,0.1);}
\providecommand{\coche}{\popcheck{popgreen}}
\providecommand{\resultat}[1]{\par\smallskip\noindent\fcolorbox{popgreen}{popgreenL}{\parbox{\dimexpr\linewidth-2\fboxsep-2\fboxrule\relax}{\textbf{Résultat.} #1}}\par\smallskip}
\newcommand{\popboxhead}[3]{\poppill{#1}{#2}{white}\ifx\relax#3\relax\else\hspace{6pt}{\popxbold\fontsize{10.6}{12}\selectfont\color{popink}#3}\fi\par\vspace{7pt}}

\newtcolorbox{aretenir}[1][]{popbase,colframe=popgreen,before upper={\popboxhead{À retenir}{popgreen}{#1}}}
\newtcolorbox{exemplebox}[1][]{popbase,colframe=poporange,before upper={\popboxhead{Exemple}{poporange}{#1}}}
\newtcolorbox{definition}[1][]{popbase,colframe=popblue,before upper={\popboxhead{Définition}{popblue}{#1}}}
\newtcolorbox{propriete}[1][]{popbase,colframe=poppurple,before upper={\popboxhead{Propriété}{poppurple}{#1}}}
\newtcolorbox{methode}[1][]{popbase,colframe=popgold,before upper={\popboxhead{Méthode}{popgold}{#1}}}
\newtcolorbox{attention}[1][]{popbase,colframe=poppink,before upper={\popboxhead{Attention}{poppink}{#1}}}
\newtcolorbox{experience}[1][]{popbase,colframe=popblue,colback=popblueL,before upper={\popboxhead{Expérience}{popblue}{#1}}}
% « Le savais-tu ? » : carte violette pleine (contexte, culture, Cameroun)
\newtcolorbox{savaistu}[1][]{popbase,colback=poppurple,colframe=popink,
  fontupper=\popsemi\fontsize{10.4}{14.2}\selectfont\color{white},
  before upper={\poppill{Le savais-tu\,?}{white}{poppurple}\ifx\relax#1\relax\else\hspace{6pt}{\popxbold\color{white}#1}\fi\par\vspace{7pt}}}
% Exercice résolu : énoncé en haut, \tcblower puis la solution sur fond crème
\newcounter{popexo}\newcounter{popexr}
\newtcolorbox{exoresolu}[1][]{popbase,colframe=poporange,colbacklower=popcream,
  segmentation style={popink,line width=1.2pt,dashed},
  before upper={\stepcounter{popexr}\popboxhead{Exercice résolu \thepopexr}{poporange}{#1}},
  before lower={\poppill{Solution}{popink}{popcream}\par\vspace{6pt}}}
% Exercice (non résolu en place) — la correction est dans la partie Corrigés
\newtcolorbox{exercice}[1][]{popbase,colframe=popink,
  before upper={\stepcounter{popexo}\popboxhead{Exercice \thepopexo}{popink}{#1}}}
% Corrigé
\newtcolorbox{corrige}[1][]{popbase,colframe=popgreen,colback=popgreenL,
  before upper={\popboxhead{Corrigé}{popgreen}{#1}}}
% Objectifs / prérequis / bilan
\newtcolorbox{objectifs}{popbase,colframe=popink,colback=poporangeL,
  before upper={\popboxhead{Objectifs de la leçon}{popink}{}}}
\newtcolorbox{prerequis}{popbase,colframe=popink,colback=popblueL,
  before upper={\popboxhead{Prérequis}{popblue}{}}}
\newtcolorbox{bilan}{popbase,colframe=popink,colback=white,boxrule=2.8pt,
  before upper={\popboxhead{Fiche bilan}{poporange}{L'essentiel à connaître pour l'examen}}}
% Figure encadrée avec légende
\newtcolorbox{popfigure}[1][]{enhanced,breakable=false,colback=white,colframe=popline,boxrule=1.4pt,arc=8pt,
  left=10pt,right=10pt,top=10pt,bottom=8pt,before skip=10pt,after skip=12pt,halign=center,
  lower separated=false,halign lower=center,
  fontlower=\popsemi\fontsize{9}{11.5}\selectfont\color{popmuted},
  #1}
\newcommand{\legende}[1]{\par\vspace{6pt}{\popsemi\fontsize{9}{11.5}\selectfont\color{popmuted}#1}}

% ============================================================
% Signature OnBuch+
% ============================================================
\newcommand{\poplogoinline}[1]{{\poplogo\fontsize{#1}{#1}\selectfont\color{popink}OnBuch\color{poporange}+}}
\newcommand{\signatureonbuch}{%
  \begin{tcolorbox}[enhanced,colback=popink,colframe=popink,boxrule=2.2pt,arc=10pt,
      drop shadow=poporange,left=14pt,right=14pt,top=10pt,bottom=10pt,before skip=0pt,after skip=0pt]
    \tikz[baseline=-0.7ex]\node[circle,fill=popgreen,draw=popcream,line width=1.3pt,minimum size=22pt,inner sep=0]{\popcheck{white}};%
    \hspace{9pt}\begin{minipage}[c]{0.62\linewidth}
      {\popxbold\fontsize{12}{13}\selectfont\color{popcream}Signé\ \textbullet\ {\poplogo OnBuch\color{poporange}+}}\par\vspace{2pt}
      {\raggedright\popsemi\fontsize{8}{10.5}\selectfont\color{popline}\MakeUppercase{Équipe pédagogique OnBuch+}\\\MakeUppercase{Conforme au programme officiel MINESEC}\par}
    \end{minipage}\hfill
    {\poplogo\fontsize{22}{22}\selectfont\color{popcream}OnBuch\color{poporange}+}
  \end{tcolorbox}%
}

% ============================================================
% Page de garde
% #1 titre · #2 matière · #3 niveau · #4 module · #5 n° leçon · #6 durée ·
% #7 description / objectifs (texte ou itemize)
% ============================================================
\newcommand{\pagegarde}[7]{%
  \thispagestyle{empty}
  \newgeometry{a4paper,margin=1.7cm,top=1.6cm,bottom=1.4cm}%
  \begin{tikzpicture}[remember picture,overlay]
    \fill[poporange!18] ([xshift=-3.2cm,yshift=-3.6cm]current page.north east) circle (4.6cm);
    \fill[poppurple!14] ([xshift=2.4cm,yshift=7.6cm]current page.south west) circle (3.8cm);
  \end{tikzpicture}%
  {\poplogo\fontsize{30}{30}\selectfont\color{popink}OnBuch\color{poporange}+}\hfill
  \raisebox{0.6ex}{\poppill{Cours signé \textbullet\ Premium}{popink}{popcream}}\par
  \vspace{1pt}\popsquiggle{poppurple}\par
  \vspace{8pt}
  \begin{tcolorbox}[enhanced,colback=poporange,colframe=popink,boxrule=2.8pt,arc=13pt,
      drop shadow=popink,left=18pt,right=18pt,top=12pt,bottom=13pt,after skip=10pt]
    \poppill{#2 \textbullet\ #3}{popink}{popcream}\hspace{5pt}\poppill{#4}{white}{popink}\par\vspace{8pt}
    {\popdisplay\fontsize{14}{14}\selectfont\color{popink}LEÇON #5}\par\vspace{4pt}
    {\raggedright\hyphenpenalty=10000\exhyphenpenalty=10000\popdisplay\fontsize{25}{27}\selectfont\color{popcream}\MakeUppercase{#1}\par}
    \vspace{8pt}
    {\popxbold\fontsize{9.5}{9.5}\selectfont\color{popink}\MakeUppercase{Durée conseillée : #6}}
  \end{tcolorbox}
  \begin{tcolorbox}[enhanced,colback=white,colframe=popink,boxrule=2.2pt,arc=10pt,
      drop shadow=popink,left=16pt,right=16pt,top=10pt,bottom=10pt,after skip=8pt]
    \poppill{Dans cette leçon}{poppurple}{white}\par\vspace{5pt}
    \popsemi\fontsize{9.3}{12.6}\selectfont\color{popink}#7
  \end{tcolorbox}
  \noindent\begin{minipage}[t]{0.487\linewidth}
    \begin{tcolorbox}[enhanced,colback=popgreen,colframe=popink,boxrule=2.2pt,arc=10pt,
        drop shadow=popink,left=11pt,right=11pt,top=10pt,bottom=10pt,height=3.1cm,valign=center]
      \tcbox[colback=white,colframe=popink,boxrule=1.3pt,arc=5pt,boxsep=0pt,nobeforeafter,
        left=3pt,right=3pt,top=3pt,bottom=3pt,valign=center]{\qrcode[height=1.9cm]{https://play.google.com/store/apps/details?id=com.luvvix.onbuch&hl=fr}}%
      \hspace{8pt}\begin{minipage}[c]{0.5\linewidth}
        {\popdisplay\fontsize{11}{12.5}\selectfont\color{white}\MakeUppercase{Télécharge OnBuch}}\par\vspace{4pt}
        {\raggedright\popsemi\fontsize{8.4}{11}\selectfont\color{white}Gratuit \textbullet\ Google Play\\Tuteur IA, annales, concours, résultats\par}
      \end{minipage}
    \end{tcolorbox}
  \end{minipage}\hfill
  \begin{minipage}[t]{0.487\linewidth}
    \begin{tcolorbox}[enhanced,colback=poppurple,colframe=popink,boxrule=2.2pt,arc=10pt,
        drop shadow=popink,left=11pt,right=11pt,top=10pt,bottom=10pt,height=3.1cm,valign=center]
      \tikz[baseline=-0.7ex]\node[circle,fill=white,draw=popink,line width=1.3pt,minimum size=1.6cm,inner sep=0]{\popdisplay\fontsize{24}{24}\selectfont\color{poppurple}+};%
      \hspace{8pt}\begin{minipage}[c]{0.6\linewidth}
        {\popdisplay\fontsize{11}{12.5}\selectfont\color{white}\MakeUppercase{Passe à OnBuch+}}\par\vspace{4pt}
        {\raggedright\popsemi\fontsize{8.4}{11}\selectfont\color{white}Tous les cours signés, Tuteur IA illimité, fiches PDF — onglet OnBuch+ de l'app\par}
      \end{minipage}
    \end{tcolorbox}
  \end{minipage}
  \vspace{8pt}
  \popcontenu
  \vfill
  \signatureonbuch
  \restoregeometry
  \newpage
}

% Rangée « Ce cours contient » (page de garde)
\newcommand{\poptile}[3]{% #1 couleur #2 gros texte #3 légende
  \begin{tcolorbox}[enhanced,colback=#1,colframe=popink,boxrule=2pt,arc=9pt,shadow={2.5pt}{-2.5pt}{0pt}{popink},
      left=6pt,right=6pt,top=9pt,bottom=9pt,halign=center,valign=center,height=1.75cm,width=\linewidth,nobeforeafter]
    {\popdisplay\fontsize{12.5}{13}\selectfont\color{white}\MakeUppercase{#2}}\par\vspace{3pt}
    {\centering\popsemi\fontsize{7.8}{9.5}\selectfont\color{white}#3\par}
  \end{tcolorbox}}
\newcommand{\popcontenu}{%
  \par\noindent{\popxboldls\fontsize{7.5}{7.5}\selectfont\color{popmuted}CE COURS SIGNÉ CONTIENT}\par\vspace{7pt}
  \noindent\begin{minipage}[t]{0.235\linewidth}\poptile{popblue}{Cours}{complet, illustré, pas à pas}\end{minipage}\hfill
  \begin{minipage}[t]{0.235\linewidth}\poptile{poporange}{Méthodes}{et exercices résolus}\end{minipage}\hfill
  \begin{minipage}[t]{0.235\linewidth}\poptile{poppink}{Exercices}{type examen + corrigés}\end{minipage}\hfill
  \begin{minipage}[t]{0.235\linewidth}\poptile{popgreen}{Fiche bilan}{l'essentiel à retenir}\end{minipage}\par}

% Sommaire
\newtcolorbox{sommaire}{enhanced,colback=white,colframe=popink,boxrule=2.2pt,arc=10pt,
  drop shadow=popink,left=18pt,right=18pt,top=13pt,bottom=13pt,before skip=6pt,after skip=20pt,
  before upper={\poppill{Sommaire}{poppurple}{white}\par\vspace{9pt}\popsemi\fontsize{10.6}{14.5}\selectfont\color{popink}}}

% Signature de fin de cours — poussée tout en bas de la dernière page
\newcommand{\findecours}{%
  \par\vfill\nopagebreak
  \begin{center}\popsquiggle{poporange}\end{center}\vspace{4pt}
  \signatureonbuch
}
\AtBeginDocument{\def\llbracket{\mathopen{[\mkern-3.5mu[}}\def\rrbracket{\mathclose{]\mkern-3.5mu]}}} % intervalles d'entiers (glyphe ⟦ absent de la police)
% Glyphes absents de Fira Math : redéfinis à partir de glyphes disponibles
\AtBeginDocument{%
  \def\setminus{\mathbin{\backslash}}\let\smallsetminus\setminus
  \def\vdots{\mathord{\vbox{\baselineskip3.2pt\lineskiplimit0pt\kern4pt\hbox{.}\hbox{.}\hbox{.}}}}%
  \def\ddots{\mathinner{\mkern1mu\raise6.5pt\hbox{.}\mkern2mu\raise3.6pt\hbox{.}\mkern2mu\raise0.7pt\hbox{.}\mkern1mu}}%
  \def\triangle{\mathord{\scalebox{0.9}{$\symup{\Delta}$}}}%
  \def\nabla{\mathord{\raisebox{\depth}{\scalebox{1}[-1]{$\symup{\Delta}$}}}}%
  \def\checkmark{\popcheck{popdark}}%
  \def\star{\mathbin{\ast}}\def\bigstar{\mathord{\ast}}%
  \def\diamond{\mathbin{\scalebox{0.6}{$\Diamond$}}}%
  \def\bigcup{\mathop{\mathchoice{\raisebox{-0.3em}{\scalebox{1.7}{$\cup$}}}{\scalebox{1.25}{$\cup$}}{\cup}{\cup}}\displaylimits}%
  \def\bigcap{\mathop{\mathchoice{\raisebox{-0.3em}{\scalebox{1.7}{$\cap$}}}{\scalebox{1.25}{$\cap$}}{\cap}{\cap}}\displaylimits}%
  \def\lesseqqgtr{\mathrel{\lessgtr}}\def\bigoplus{\mathop{\oplus}\displaylimits}%
}
\providecommand{\ssi}{si et seulement si} % abréviation fréquente
\AtBeginDocument{\providecommand{\wideparen}{\overparen}} % arc de cercle

%% FIGFIX-BEGIN v5 — correctifs globaux des figures (docs/onbuch-generation/tools/figfix)
\usepackage{adjustbox}\usepackage{etoolbox}\tcbuselibrary{raster}
% toute figure TikZ/pgfplots plus large que la ligne est réduite à la largeur disponible
\BeforeBeginEnvironment{tikzpicture}{\begin{adjustbox}{max width=\linewidth}}
\AfterEndEnvironment{tikzpicture}{\end{adjustbox}}
% légendes pgfplots lisibles par-dessus les courbes
\pgfplotsset{every axis/.append style={legend style={fill=white,fill opacity=0.92,text opacity=1,draw=black!15,inner sep=3pt}}}
% étiquettes d'arêtes (syntaxe edge["..."]) avec fond blanc
\tikzset{every edge quotes/.append style={fill=white,inner sep=1.2pt}}
%% FIGFIX-END
\begin{document}

\pagegarde{L’industrie}{Géographie}{Tle D}{Module 1 : Le Cameroun}{N08}{2 h + TP 3 (2 h)}{Cette leçon étudie l'industrie au Cameroun : ses formidables atouts énergétiques et miniers, l'élan industriel brisé par la crise des années 1980-1990, et la structure actuelle dominée par les agro-industries et les PME. Elle s'appuie sur l'analyse de cartes, de photographies et de données chiffrées, et comprend un TP de construction d'un camembert de la production d'énergie.
\vspace{4pt}
\begin{multicols}{2}\small
\begin{itemize}
\item À la fin de cette leçon, je sais définir l'industrie et distinguer les grands types d'industries (lourdes, agro-industries, PME)
\item À la fin de cette leçon, je sais localiser sur une carte les principaux barrages hydroélectriques, centrales à gaz et sites solaires du Cameroun
\item À la fin de cette leçon, je sais inventorier les principales réserves minières du Cameroun (bauxite, fer, terres rares) et les situer
\item À la fin de cette leçon, je sais expliquer comment la crise économique a brisé l'élan industriel camerounais
\item À la fin de cette leçon, je sais citer des exemples d'industries lourdes (ALUCAM, DANGOTE, CIMENCAM) et montrer la prédominance des agro-industries et des PME
\item À la fin de cette leçon, je sais décrire une zone industrielle et expliquer son rôle dans l'intégration nationale
\end{itemize}\end{multicols}}

\begin{sommaire}
\begin{multicols}{2}
\begin{enumerate}
\item Généralités et définitions
\item Un énorme potentiel énergétique
\item D'importantes réserves minières
\item Un élan industriel brisé par la crise
\item Les types d'industries au Cameroun
\item Les zones industrielles et les perspectives
\item TP 3 : Construction d'un camembert de la production d'énergie
\item Activité d'intégration
\item Méthodes et exercices résolus
\item Exercices
\item Corrigés des exercices
\item Fiche bilan
\end{enumerate}
\end{multicols}
\end{sommaire}

\begin{objectifs}
\begin{itemize}
\item À la fin de cette leçon, je sais définir l'industrie et distinguer les grands types d'industries (lourdes, agro-industries, PME)
\item À la fin de cette leçon, je sais localiser sur une carte les principaux barrages hydroélectriques, centrales à gaz et sites solaires du Cameroun
\item À la fin de cette leçon, je sais inventorier les principales réserves minières du Cameroun (bauxite, fer, terres rares) et les situer
\item À la fin de cette leçon, je sais expliquer comment la crise économique a brisé l'élan industriel camerounais
\item À la fin de cette leçon, je sais citer des exemples d'industries lourdes (ALUCAM, DANGOTE, CIMENCAM) et montrer la prédominance des agro-industries et des PME
\item À la fin de cette leçon, je sais décrire une zone industrielle et expliquer son rôle dans l'intégration nationale
\item À la fin de cette leçon, je sais construire et légender un cercle (camembert) représentant la production d'énergie par source
\item À la fin de cette leçon, je sais interpréter des documents géographiques (cartes, photographies, tableaux statistiques) relatifs à l'industrie camerounaise
\end{itemize}
\end{objectifs}

\begin{prerequis}
\begin{itemize}
\item Les grandes régions naturelles et les fleuves du Cameroun (Sanaga, Wouri, Bénoué) vus en Première
\item La notion de ressources naturelles et de développement (Module 1, leçons précédentes)
\item La crise économique des années 1980-1990 au Cameroun (repères historiques)
\item Les notions de pourcentage et de proportionnalité (mathématiques, utiles pour le TP)
\item La lecture et la légendation d'une carte (savoir-faire de Seconde et Première)
\end{itemize}
\end{prerequis}

\newpage

\coursec{Généralités et définitions}

En juillet 2023, plusieurs quartiers de Douala et de Yaoundé subissent des délestages quotidiens : les boutiques ferment plus tôt, les soudeurs interrompent leur travail, les congélateurs des poissonneries de Mokolo dégèlent. Pourtant, le Cameroun possède le deuxième potentiel hydroélectrique d'Afrique (après la RDC), estimé à environ \num{20000} MW, dont à peine un dixième est réellement exploité. Au même moment, l'usine ALUCAM d'Édéa, fleuron historique de l'industrie nationale, tourne au ralenti, tandis que les cimenteries DANGOTE (usine de broyage de Mimboman) et CIMENCAM (usine de Bonabéri), à Douala, produisent à plein régime pour répondre à la demande du bâtiment. Ce contraste saisissant pose la problématique de toute la leçon : comment un pays doté d'un énorme potentiel énergétique et minier peut-il connaître un élan industriel brisé ? Quels types d'industries dominent aujourd'hui le tissu productif camerounais ? Pour répondre à ces questions, il faut d'abord s'entendre sur ce qu'est l'industrie, distinguer ses grandes catégories, mesurer sa place dans l'économie nationale et comprendre son rôle dans l'intégration du territoire.

\courssub{Qu'est-ce que l'industrie ?}

\textbf{L'intuition d'abord.} Observe un simple paquet de biscuits acheté au marché de Mfoundi, une tige de fer chez le quincailler de Bafoussam ou une bouteille de gaz domestique à Maroua. Aucun de ces objets ne sort directement de la nature : le blé a été moulu, le minerai de fer a été fondu, le pétrole a été raffiné. Entre la matière première et l'objet final, il y a eu \cle{transformation}. C'est précisément cette transformation qui définit l'industrie.

\begin{definition}[L'industrie]
L'\cle{industrie} est l'ensemble des activités économiques qui transforment les matières premières (agricoles, minières, forestières, énergétiques) en biens de consommation (savon, farine, boissons, ciment) ou en biens d'équipement (machines, matériaux, installations). Elle constitue, avec la construction et l'énergie, le cœur du \cle{secteur secondaire} de l'économie.
\end{definition}

Le résultat de la transformation peut être un \cle{produit fini}, directement consommable (un sac de ciment CIMENCAM, une bouteille de jus TOP), ou un \cle{produit semi-fini}, qui servira d'intrant à une autre industrie (l'aluminium d'ALUCAM, la farine livrée aux boulangeries, les plaques de tôle destinées aux ateliers de chaudronnerie).

\begin{attention}[Industrie ou artisanat ?]
Ne confonds pas l'\cle{industrie} et l'\cle{artisanat}. L'artisanat est une production essentiellement manuelle, à petite échelle, souvent familiale (le tailleur de quartier, le forgeron traditionnel, la potière). L'industrie repose sur des machines, une organisation du travail en usine et une production en série. Au Cameroun, la frontière est parfois floue : les \cle{PME} (petites et moyennes entreprises) font le lien entre les deux, comme les minoteries semi-industrielles ou les unités de transformation artisanale du manioc en gari. Au Baccalauréat, précise toujours l'échelle et le niveau de mécanisation.
\end{attention}

\courssub{Les trois grandes catégories d'industries}

Pour classer les industries, on se pose une question simple : d'où vient la matière et que devient-elle ? Trois grandes catégories apparaissent.

\begin{definition}[Extractive, transformation, énergétique]
\begin{itemize}
\item L'\cle{industrie extractive} prélève les richesses du sous-sol ou du sol sans les transformer profondément : extraction du pétrole brut (opérateurs comme PERENCO, sous contrat SNH), du gaz naturel, des minerais (bauxite, fer), exploitation forestière.
\item L'\cle{industrie de transformation} convertit les matières premières en produits nouveaux : raffinage du pétrole (SONARA à Limbé), cimenteries (CIMENCAM, DANGOTE, FIGUIL), brasseurs (SABC), savonneries, agro-industries (SOSUCAM, PHP, CDC).
\item L'\cle{industrie énergétique} produit et distribue l'énergie nécessaire à toutes les autres : barrages hydroélectriques (Songloulou, Édéa, Lagdo, Memve'ele), centrales thermiques et à gaz (Kribi, Dibamba), production d'électricité assurée principalement par ENEO.
\end{itemize}
\end{definition}

Ces trois catégories sont en réalité enchaînées : l'extractive fournit la matière, la transformation lui donne de la valeur, l'énergétique fournit la force motrice. C'est ce qu'illustre la chaîne de production ci-dessous, à savoir reproduire dans un croquis.

\begin{popfigure}
\begin{center}
\begin{tikzpicture}[
box/.style={draw=popblue, fill=popblueL, rounded corners=3pt, minimum height=1.1cm, text width=2.6cm, align=center, font=\small},
arr/.style={-{Stealth[length=3mm]}, very thick, poporange},
lab/.style={font=\scriptsize\itshape, popmuted, align=center}
]
\node[box, align=center] (mp) at (0,0){\textbf{Matière première}\\ coton, cacao, bauxite, pétrole brut};
\node[box, align=center] (tr) at (4.6,0){\textbf{Transformation}\\ usine, machines, main-d'œuvre};
\node[box, align=center] (pf) at (9.2,0){\textbf{Produit fini}\\ tissu, chocolat, aluminium, essence};
\node[box, fill=popgreenL, draw=popgreen, align=center] (ma) at (13.4,0){\textbf{Marché}\\ consommateurs, exportation};
\draw[arr] (mp) -- (tr);
\draw[arr] (tr) -- (pf);
\draw[arr] (pf) -- (ma);
\node[box, fill=popgoldL, draw=popgold, align=center] (en) at (4.6,-2.2){\textbf{Énergie}\\ barrages, centrales à gaz};
\draw[arr, popgold] (en) -- (tr);
\node[lab] at (2.3,0.55) {1. approvisionnement};
\node[lab] at (6.9,0.55) {2. fabrication};
\node[lab] at (11.3,0.55) {3. distribution};
\end{tikzpicture}
\end{center}
\legende{Schéma de la chaîne de production industrielle : 1. la matière première est extraite ou récoltée ; 2. elle est transformée en usine grâce à l'énergie et à la main-d'œuvre ; 3. le produit fini est distribué sur le marché national ou exporté. L'énergie (en jaune) est l'intrant indispensable de toute transformation.}
\end{popfigure}

\begin{exemplebox}[De la fève de cacao à la tablette de chocolat]
Le Cameroun produit environ \num{300000} tonnes de cacao par an (ordre de grandeur), mais n'en transforme localement qu'environ 20 à 30 \%. La filière illustre la chaîne : la fève (matière première, régions du Centre et du Sud-Ouest) est séchée, expédiée à Douala, broyée et transformée en pâte et en beurre de cacao (produits semi-finis), puis une partie devient chocolat (produit fini). Tant que le pays exporte surtout la fève brute, il laisse l'essentiel de la \cle{valeur ajoutée} — et des emplois — aux industriels étrangers.
\end{exemplebox}

\courssub{La place de l'industrie dans l'économie camerounaise}

Pour mesurer le poids de l'industrie, les géographes utilisent deux indicateurs : la contribution au \cle{PIB} (produit intérieur brut, c'est-à-dire la richesse produite en un an) et la part de l'\cle{emploi}.

L'économie camerounaise se répartit en trois secteurs : le \cle{secteur primaire} (agriculture, élevage, pêche, forêt), le \cle{secteur secondaire} (industries, énergie, construction) et le \cle{secteur tertiaire} (commerce, transports, services, administration). Les ordres de grandeur récents (sources INS et Banque mondiale, à retenir comme approximations) sont les suivants :

\begin{center}
\begin{tabularx}{\linewidth}{|l|X|X|}
\hline
\rowcolor{popblueL} \textbf{Secteur} & \textbf{Part approximative du PIB} & \textbf{Contenu principal} \\
\hline
Primaire & environ 15 à 17 \% & cacao, café, coton, bananes, vivrier, bois \\
\hline
Secondaire & environ 25 à 28 \% & agro-industries, ciment, pétrole raffiné, énergie, BTP \\
\hline
Tertiaire & environ 50 à 55 \% & commerce, transports, télécommunications, services publics \\
\hline
\end{tabularx}
\end{center}

\begin{popfigure}
\begin{center}
\begin{tikzpicture}
\begin{axis}[
ybar, bar width=1.4cm,
width=0.85\linewidth, height=6cm,
ymin=0, ymax=60,
ylabel={Part du PIB (en \%)},
symbolic x coords={Primaire, Secondaire, Tertiaire},
xtick=data,
nodes near coords,
every node near coord/.append style={font=\small\bfseries},
ymajorgrids=true, grid style={popline!40},
]
\addplot[fill=popgreen, draw=popdark] coordinates {(Primaire,16)};
\addplot[fill=poporange, draw=popdark] coordinates {(Secondaire,27)};
\addplot[fill=popblue, draw=popdark] coordinates {(Tertiaire,52)};
\end{axis}
\end{tikzpicture}
\end{center}
\legende{Diagramme en barres de la répartition approximative du PIB du Cameroun par secteur (ordres de grandeur, sources INS et Banque mondiale, années récentes). Le secteur secondaire représente environ un quart de la richesse nationale ; la somme des trois barres est inférieure à 100 \% car les impôts nets sur produits (TVA, droits de douane) ne sont pas répartis par secteur ici.}
\end{popfigure}

\begin{exemplebox}[Un quart du PIB, mais une industrie dominée par l'agro-alimentaire]
Le secteur secondaire représente environ un quart du PIB camerounais. Mais attention à la composition : sur ce quart, une part importante revient aux \cle{agro-industries} (brasseries, sucreries comme SOSUCAM, huileries, minoteries) et au BTP, tandis que les industries lourdes (métallurgie, chimie) restent très minoritaires. En termes d'emploi, l'industrie moderne n'occupe qu'une faible part de la population active (environ 10 \% en comptant le secteur informel manufacturier), loin derrière l'agriculture qui emploie encore près d'un actif sur deux.
\end{exemplebox}

\begin{attention}[Deux pièges de rédaction fréquents]
Premier piège : écrire que « le Cameroun est un pays industriel ». Faux : avec environ un quart du PIB et une faible part de l'emploi, c'est un pays à \cle{industrialisation inachevée}. Deuxième piège : confondre « secteur secondaire » et « industrie » : le secondaire inclut aussi le BTP et l'énergie ; l'industrie manufacturière seule pèse moins lourd (de l'ordre de 12 à 15 \% du PIB, ordre de grandeur à manier avec prudence).
\end{attention}

\courssub{Localisation des principales unités industrielles citées}

Avant d'analyser le rôle intégrateur, il est indispensable de situer dans l'espace les exemples mobilisés jusqu'ici. Le croquis schématique ci-dessous localise les sites clés : les pôles industriels côtiers (Douala, Limbé, Kribi), les sites énergétiques et métallurgiques du Sud (Édéa, Songloulou, Memve'ele), les agro-industries du Centre (Nkoteng, Yaoundé) et les ressources minières et industrielles du Grand Nord (Figuil, Lagdo, Ngaoundal, Minim-Martap, Maroua).

\begin{popfigure}
\begin{center}
\includegraphics[width=\linewidth,height=9.5cm,keepaspectratio]{N01/cmr_map.jpg}
\end{center}
\legende{Fond de carte du Cameroun (relief, régions, villes, fleuves, routes, voies ferrées) sur lequel situer les unités industrielles, sites énergétiques et ressources minières cités dans la leçon : les pôles côtiers (Douala, Édéa, Kribi, Limbé), les barrages de la Sanaga et de la Bénoué (retenue de Lagdo visible au nord), les usines de l'intérieur (Yaoundé, Bafoussam, Figuil) et les gisements de l'Adamaoua et de l'Est. Observe l'axe Douala--Yaoundé--Ngaoundéré, que suit la voie ferrée. \\ Source : Wikimedia Commons, CIA, domaine public.}
\end{popfigure}

\courssub{L'industrie, facteur d'intégration nationale}

Pourquoi le programme classe-t-il cette leçon dans la famille de situations « L'intégration nationale » ? Parce que l'industrie ne produit pas seulement des biens : elle \cle{tisse des liens} entre les régions, entre les villes et les campagnes, entre les Camerounais eux-mêmes.

\begin{propriete}[Les quatre rôles intégrateurs de l'industrie]
\begin{enumerate}
\item \textbf{Création d'emplois et de revenus} : chaque usine (SOSUCAM à Nkoteng, CDC à Tiko, cimenterie de Figuil) fixe des populations, fait vivre des sous-traitants, des transporteurs et des commerçants, et réduit l'exode rural.
\item \textbf{Désenclavement} : pour évacuer ses produits et s'approvisionner, l'industrie exige routes, voies ferrées et ports ; l'axe Douala--Yaoundé, le chemin de fer Transcamerounais (Camrail) ou le port en eau profonde de Kribi doivent beaucoup aux besoins industriels.
\item \textbf{Lien villes--campagnes} : les agro-industries achètent les récoltes des campagnes (canne, coton, huile de palme) et renvoient vers les campagnes des produits manufacturés et des revenus monétaires.
\item \textbf{Valorisation des ressources régionales} : le ciment de Figuil utilise le calcaire du Nord, le sucre de Nkoteng la canne du Centre, l'aluminium d'Édéa l'électricité du fleuve Sanaga : chaque région peut ainsi participer à l'effort national.
\end{enumerate}
\end{propriete}

\begin{savaistu}[Le ciment de Figuil, symbole d'intégration]
La cimenterie de Figuil (région du Nord), créée dans les années 1980 puis reprise dans le giron de CIMENCAM, exploite le calcaire local pour fournir du ciment à tout le Grand Nord, jusque-là dépendant des livraisons lointaines de Douala. Elle a créé des emplois, attiré des services et réduit le coût de construction dans une région longtemps enclavée : un cas d'école d'industrie au service de l'intégration nationale.
\end{savaistu}

\courssub{Le problème posé : un potentiel considérable, un tissu fragile}

Résumons le paradoxe camerounais que toute la leçon va déplier. D'un côté, les atouts sont remarquables : deuxième potentiel hydroélectrique d'Afrique, gisements de bauxite parmi les plus importants du monde (Minim-Martap et Ngaoundal), fer de Mbalam (frontière Congo), pétrole et gaz du bassin du Rio del Rey, terres rares identifiées à Akonolinga, soleil abondant dans le Grand Nord. De l'autre côté, les réalités sont cruelles : délestages chroniques, fermetures d'usines depuis la crise économique des années 1985--1995, ALUCAM au ralenti, dépendance aux importations de produits manufacturés, concentration excessive des activités à Douala.

\begin{methode}[Analyser un secteur industriel : la démarche en quatre temps]
Pour toute étude de cas industrielle au Baccalauréat, suis toujours le même plan :
\begin{enumerate}
\item \textbf{Atouts} : quelles ressources (énergie, minerais, matières agricoles, main-d'œuvre, marché) ?
\item \textbf{Réalisations} : quelles unités existent, où, que produisent-elles, pour quel marché ?
\item \textbf{Limites} : quels obstacles (énergie insuffisante, coûts, concurrence, enclavement, crise) ?
\item \textbf{Perspectives} : quels projets et quelles solutions (nouveaux barrages, zones industrielles, transformation locale) ?
\end{enumerate}
Cette démarche « atouts $\rightarrow$ réalisations $\rightarrow$ limites $\rightarrow$ perspectives » sera notre fil conducteur dans toute la leçon.
\end{methode}

\begin{exoresolu}[Première application : classer des entreprises camerounaises]
Énoncé. Classe chacune des entreprises suivantes dans la bonne catégorie d'industrie (extractive, transformation, énergétique) et justifie : la SONARA de Limbé, PERENCO Cameroun, ENEO, la SABC, la SODECOTON.
\tcblower
Solution détaillée. On applique le critère « d'où vient la matière, que devient-elle ? ».
\begin{itemize}
\item \textbf{PERENCO Cameroun} : industrie \cle{extractive} — elle pompe le pétrole brut et le gaz sans les transformer en produits finis.
\item \textbf{La SONARA} (Limbé) : industrie de \cle{transformation} — elle raffine le pétrole brut en essence, gasoil, kérosène.
\item \textbf{ENEO} : industrie \cle{énergétique} — elle produit et distribue l'électricité issue des barrages et centrales.
\item \textbf{La SABC} (brasseries) : industrie de \cle{transformation} — elle transforme céréales, sucre et eau en boissons ; c'est aussi une agro-industrie.
\item \textbf{La SODECOTON} : double cas — elle collecte le coton-graine (amont agricole) mais surtout l'égrenne et le file : c'est une industrie de \cle{transformation} agro-industrielle, pivot du lien villes--campagnes dans le Grand Nord.
\end{itemize}
Piège évité : ne pas classer ENEO dans la transformation « parce qu'elle produit quelque chose » ; l'énergie est une catégorie à part.
\end{exoresolu}

\begin{aretenir}[L'essentiel de la section]
\begin{itemize}
\item L'industrie transforme les matières premières en produits finis ou semi-finis ; on distingue les industries \cle{extractives}, de \cle{transformation} et \cle{énergétiques}.
\item Le secteur secondaire représente environ un quart du PIB camerounais, dominé par les agro-industries ; l'industrie lourde est minoritaire et l'emploi industriel reste faible (environ 10 \% formel + informel).
\item L'industrie est un puissant facteur d'\cle{intégration nationale} : emplois, désenclavement, lien villes--campagnes, valorisation des ressources régionales.
\item Le paradoxe camerounais : un potentiel énergétique et minier considérable face à un tissu industriel fragile, brisé par la crise.
\end{itemize}
\end{aretenir}

Maintenant que les définitions sont posées et le paradoxe identifié, examinons dans la section suivante le premier grand atout du pays : son énorme potentiel énergétique.

\coursec{Un énorme potentiel énergétique}

Dans la section précédente, nous avons défini l'industrie et montré qu'aucune usine ne peut fonctionner sans une source d'énergie fiable et abondante. Or, avant même de parler d'usines, il faut se poser une question simple : le Cameroun dispose-t-il de quoi produire l'électricité et la chaleur nécessaires à son industrialisation ? La réponse est étonnante : oui, largement, et c'est même l'un des pays d'Afrique les mieux dotés. C'est ce que nous allons découvrir maintenant, en partant d'une situation concrète.

\begin{exemplebox}[Situation-problème : une panne à Douala]
Un soir de février, dans le quartier de Bonabéri à Douala, les lumières s'éteignent brusquement : c'est un délestage. Pourtant, à quelques dizaines de kilomètres de là, le fleuve Sanaga charrie chaque seconde des milliers de mètres cubes d'eau vers l'océan. Comment un pays traversé par de puissants fleuves, baigné de soleil au nord et doté de gaz au large de ses côtes peut-il manquer d'électricité ? C'est tout l'objet de cette section.
\end{exemplebox}

\courssub{Le potentiel hydroélectrique : un géant qui sommeille}

\begin{definition}[Barrage hydroélectrique]
Un \cle{barrage hydroélectrique} est un ouvrage construit sur un cours d'eau pour retenir l'eau et créer un dénivelé (une \emph{chute}). L'eau retenue, lâchée sous pression, fait tourner des \cle{turbines} couplées à des alternateurs : l'énergie de l'eau (énergie mécanique) est ainsi transformée en \cle{électricité}. La puissance d'un barrage se mesure en mégawatts (MW).
\end{definition}

L'intuition est simple : plus un fleuve a un \textbf{débit} important et plus le relief offre de \textbf{dénivelés} (rapides, chutes), plus l'énergie récupérable est grande. Or le Cameroun cumule ces deux atouts :

\begin{itemize}
\item un réseau hydrographique dense et puissant, alimenté par des pluies abondantes : la \cle{Sanaga} (le fleuve le plus régulier et le plus exploité), la \cle{Bénoué} au nord, le \cle{Dja} et la \cle{Lobé} au sud, le \cle{Ntem} à la frontière avec la Guinée équatoriale ;
\item un relief contrasté, avec les abrupts de l'Adamaoua et les chutes du bassin de la Sanaga, qui créent des sites naturels de barrages.
\end{itemize}

\begin{propriete}[Un classement remarquable]
Le potentiel hydroélectrique du Cameroun est estimé à environ \textbf{\num{20000} MW}, ce qui en fait le \cle{deuxième potentiel d'Afrique} après la République démocratique du Congo (RDC). Mais moins de 10 \% de ce potentiel est réellement exploité : la puissance installée du pays est de l'ordre de \num{1500} à \num{1600} MW (ordre de grandeur, toutes sources confondues, au début des années 2020).
\end{propriete}

\begin{savaistu}[Le géant endormi]
Avec environ \num{20000} MW de potentiel hydroélectrique, le Cameroun pourrait théoriquement produire plus de dix fois ses besoins actuels et même exporter de l'électricité vers le Tchad, le Nigeria ou le Gabon. Certains sites identifiés, comme les chutes de Grand Eweng sur la Sanaga, pourraient à eux seuls dépasser \num{1000} MW. Le problème n'est donc pas la nature, mais le financement et la réalisation des ouvrages.
\end{savaistu}

\courssub{Les grands barrages hydroélectriques du Cameroun}

L'essentiel de l'électricité camerounaise (environ 60 à 70 \% selon les années) provient de l'hydroélectricité. Retenons les cinq ouvrages majeurs :

\begin{center}
\begin{tabularx}{\linewidth}{|l|X|l|X|}
\hline
\rowcolor{popblueL} \textbf{Barrage} & \textbf{Cours d'eau / région} & \textbf{Puissance} & \textbf{Particularité} \\
\hline
Édéa & Sanaga (Littoral) & environ 263 MW & Le plus ancien (années 1950), construit pour ALUCAM \\
\hline
Songloulou & Sanaga (Centre) & environ 384 MW & Longtemps la première centrale du pays \\
\hline
Lagdo & Bénoué (Nord) & 72 MW & Régulation des crues et irrigation (SEMRY) \\
\hline
Memve'ele & Ntem (Sud) & environ 211 MW & Mis en service au début des années 2020 \\
\hline
Nachtigal & Sanaga (Centre) & 420 MW & Achevé vers 2024, plus grande centrale du pays \\
\hline
\end{tabularx}
\end{center}

Quelques commentaires s'imposent. Les barrages d'\textbf{Édéa} et de \textbf{Songloulou}, tous deux sur la Sanaga, forment historiquement la colonne vertébrale électrique du pays : ils alimentent Douala, Yaoundé et l'usine ALUCAM d'Édéa. Leur dépendance au régime des pluies explique les baisses de production en saison sèche prolongée. Le barrage de \textbf{Memve'ele}, sur le Ntem, a été construit avec l'appui de la Chine et illustre la volonté d'électrifier le Sud forestier. Enfin, \textbf{Nachtigal}, amont de la Sanaga, est le projet le plus récent : avec 420 MW, il accroît d'environ 30 \% la capacité de production nationale.

\begin{exemplebox}[Lagdo : un barrage à plusieurs fonctions]
Le barrage de \cle{Lagdo} (72 MW), sur la Bénoué dans la région du Nord, ne sert pas qu'à produire de l'électricité. Son réservoir permet de \textbf{réguler les crues} de la Bénoué (même si ses lâchers d'eau provoquent parfois des inondations en aval, jusque dans l'État nigérian d'Adamawa) et d'\textbf{irriguer} les plaines de la SEMRY (riziculture de Yagoua) en saison sèche. C'est un exemple type d'\emph{aménagement polyvalent} : énergie, agriculture, protection contre les inondations.
\end{exemplebox}

\begin{experience}[Photographie commentée : le barrage de Songloulou]
\textbf{Observation.} Sur la photographie aérienne du barrage de Songloulou, on distingue : (1) un long mur de béton barant la vallée de la Sanaga ; (2) en amont, un vaste plan d'eau calme (le \emph{réservoir}) ; (3) en aval, des conduites forcées qui plongent vers les bâtiments abritant les turbines ; (4) des pylônes et lignes à haute tension qui partent vers Douala et Yaoundé.

\textbf{Interprétation.} Le réservoir stocke l'eau en saison des pluies pour la restituer en saison sèche : c'est une \emph{réserve d'énergie}. La dénivellation entre l'amont et l'aval crée la pression qui fait tourner les turbines. Les lignes à haute tension transportent l'électricité sur de longues distances avec peu de pertes.

\textbf{Conclusion.} Le barrage est à la fois un ouvrage hydraulique (gestion de l'eau) et une usine électrique : il concentre en un point l'énergie dispersée dans tout le bassin versant de la Sanaga.
\end{experience}

\courssub{Les centrales à gaz et les centrales thermiques : le complément indispensable}

L'hydroélectricité dépend du régime des pluies. Pour sécuriser l'approvisionnement, le Cameroun a développé des centrales fonctionnant aux \textbf{énergies fossiles}.

\textbf{Les centrales à gaz.} La centrale de \cle{Kribi} (environ 216 MW, extensible), sur la côte du Sud, utilise le gaz naturel du gisement offshore de \textbf{Sanaga Sud}, exploité par la Société nationale des hydrocarbures (SNH) et ses partenaires (Perenco, Bowleven). C'est la première centrale à gaz du pays : le gaz est acheminé par un gazoduc sous-marin puis terrestre. La centrale de \cle{Logbaba}, à Douala, fonctionne également au gaz (et au fuel en secours), avec une puissance d'environ 80 à 100 MW. Avantage du gaz : une mise en route rapide et des émissions de CO$_2$ moindres que le fuel.

\textbf{Les centrales thermiques au fuel.} Disséminées dans le pays (Douala, Yaoundé, Garoua, Maroua, Bamenda, etc.), ces centrales brûlent du \cle{fuel} (pétrole lourd) ou du gasoil, souvent approvisionné par la raffinerie de la \textbf{SONARA} (Limbe). Elles jouent un rôle d'\textbf{appoint} : elles sont mobilisées en saison sèche, lors des pointes de consommation ou dans les villes éloignées du réseau interconnecté sud. Leur inconvénient majeur est le \textbf{coût} : le fuel est importé ou acheté à la raffinerie, ce qui grève les finances d'\textbf{ENEO} (Energy of Cameroon, concessionnaire du service public) et de l'État.

\begin{attention}[Potentiel $\neq$ production]
Ne confonds jamais le \cle{potentiel énergétique} (ce que la nature met à disposition : débits des fleuves, ensoleillement, gisements de gaz) et la \cle{production effective} (ce qui est réellement produit grâce aux ouvrages construits). Dire « le Cameroun a un énorme potentiel » ne signifie pas « le Cameroun produit beaucoup d'électricité ». Au Baccalauréat, cette distinction est souvent le cœur du sujet : c'est précisément l'écart entre les deux qui explique les délestages.
\end{attention}

\courssub{L'énergie solaire : l'atout du Grand Nord}

Le Grand Nord camerounais (régions du Nord, de l'Extrême-Nord et de l'Adamaoua) bénéficie d'un \textbf{ensoleillement exceptionnel} : plus de \num{2500} à \num{3000} heures de soleil par an, avec un rayonnement élevé et régulier. C'est une zone idéale pour le \cle{solaire photovoltaïque} (panneaux qui convertissent directement la lumière en électricité).

\begin{itemize}
\item Les \textbf{centrales solaires de Guider} (environ 30 MWc) \textbf{et de Maroua} (environ 20 MWc) ont été construites pour renforcer l'alimentation du réseau interconnecté nord, longtemps tributaire du seul barrage de Lagdo et de coûteuses centrales thermiques.
\item De nombreux \textbf{projets solaires ruraux} (mini-réseaux villageois, pompage solaire, électrification des centres de santé et des lycées) visent les localités éloignées des grands réseaux, où l'extension des lignes coûterait trop cher.
\end{itemize}

Le solaire présente toutefois une limite : il ne produit que le jour, d'où la nécessité de batteries ou d'un couplage avec d'autres sources (hybridation).

\courssub{Le paradoxe énergétique camerounais}

Nous pouvons maintenant résoudre la situation-problème du début. Le Cameroun souffre d'un véritable \cle{paradoxe énergétique} :

\begin{itemize}
\item un potentiel parmi les plus importants d'Afrique (hydroélectricité, gaz, soleil, voire biomasse) ;
\item mais une production insuffisante et mal répartie : \textbf{délestages fréquents} dans les grandes villes, coupures qui pénalisent les entreprises, et un \textbf{accès inégal à l'électricité} — les villes du Sud sont bien desservies, tandis qu'une partie importante des campagnes, surtout au nord et à l'est, n'est pas raccordée au réseau (le taux d'accès national reste de l'ordre de 60 à 65 \%, avec de forts écarts ville/campagne ; ordre de grandeur).
\end{itemize}

Les causes sont multiples : coût très élevé des barrages (plusieurs centaines de milliards de FCFA), retard des projets, vieillissement des équipements, pertes sur le réseau (technique et non technique), croissance rapide de la demande (urbanisation, industries). Les récents ouvrages (Memve'ele, Nachtigal) et les projets solaires traduisent la volonté de l'État de combler cet écart, condition indispensable de l'émergence industrielle.

\courssub{Localisation des infrastructures énergétiques}

\begin{methode}[Légender une carte d'implantation énergétique]
Pour représenter les infrastructures énergétiques sur un fond de carte :
\begin{enumerate}
\item Tracer le contour simplifié du pays et placer les principaux repères (villes, fleuves).
\item Choisir \textbf{un symbole distinct par type de source} : par exemple un rectangle bleu pour un barrage hydroélectrique, un triangle orange pour une centrale à gaz, un cercle jaune pour une centrale solaire.
\item Proportionner éventuellement la taille du symbole à la puissance de l'ouvrage.
\item Rédiger une \textbf{légende ordonnée} (du plus important au moins important), donner un titre précis (quoi ? où ?) et orienter la carte (flèche du nord).
\end{enumerate}
\end{methode}

\begin{popfigure}
\begin{center}
\includegraphics[width=\linewidth,height=9.5cm,keepaspectratio]{N01/cmr_map.jpg}
\end{center}
\legende{Fond de carte pour le croquis des infrastructures énergétiques : on y repère la Sanaga (barrages de Songloulou et d'Édéa, Nachtigal), la Bénoué avec la retenue de Lagdo, le littoral (centrales à gaz de Kribi et de Douala-Logbaba) et la Ntem au Sud (Memve'ele), ainsi que les villes du Grand Nord (Guider, Maroua) pour la solaire. Sur ce fond, on place un rectangle bleu par barrage, un triangle orange par centrale à gaz et un cercle jaune par centrale solaire, puis on ajoute légende, titre et flèche du nord. \\ Source : Wikimedia Commons, CIA, domaine public.}
\end{popfigure}

Une fois le croquis réalisé, sa lecture fait apparaître trois enseignements majeurs : (1) la \textbf{Sanaga est l'axe énergétique majeur} du pays, avec trois barrages en cascade ; (2) le \textbf{littoral} concentre les centrales à gaz, à proximité du gisement de Sanaga Sud et du premier bassin de consommation (Douala) ; (3) le \textbf{Grand Nord} mise sur le couple Lagdo + solaire. Cette géographie énergétique conditionne directement la localisation des industries, que nous étudierons dans les sections suivantes.

\begin{exoresolu}[Calculer une part de la production électrique]
\textbf{Énoncé.} En 2023, la puissance installée du parc hydroélectrique camerounais est estimée à environ \num{950} MW (avant la montée en puissance complète de Nachtigal), sur une puissance totale installée d'environ \num{1550} MW. (1) Calcule la part de l'hydroélectricité dans la puissance totale installée. (2) Sachant que le potentiel hydroélectrique est estimé à \num{20000} MW, calcule le taux d'exploitation de ce potentiel. (3) Conclus en deux phrases.
\tcblower
\textbf{Solution détaillée.}
\begin{enumerate}
\item Part de l'hydroélectricité :
\[ \frac{950}{1550} \times 100 \approx 61{,}3\ \%. \]
L'hydroélectricité représente donc environ \textbf{61 \%} de la puissance installée.
\item Taux d'exploitation du potentiel :
\[ \frac{950}{20000} \times 100 = 4{,}75\ \%. \]
Moins de \textbf{5 \%} du potentiel hydroélectrique est exploité (moins de 10 \% même en ajoutant Nachtigal et Memve'ele).
\item \textbf{Conclusion.} L'électricité camerounaise repose majoritairement sur l'eau, ce qui la rend vulnérable aux variations climatiques. Surtout, l'écart immense entre le potentiel (\num{20000} MW) et la production effective illustre le paradoxe énergétique : le pays dispose d'une richesse naturelle considérable mais insuffisamment valorisée, d'où les délestages et le déficit d'accès à l'électricité.
\end{enumerate}
\end{exoresolu}

\begin{aretenir}[L'essentiel de la section]
\begin{itemize}
\item Le Cameroun possède le \textbf{deuxième potentiel hydroélectrique d'Afrique} (environ \num{20000} MW), grâce à ses fleuves (Sanaga, Bénoué, Ntem, Dja) et à son relief.
\item Les grands barrages : \textbf{Édéa} et \textbf{Songloulou} puis \textbf{Nachtigal} sur la Sanaga, \textbf{Lagdo} sur la Bénoué (ouvrage polyvalent : énergie, crues, irrigation), \textbf{Memve'ele} sur le Ntem.
\item Les \textbf{centrales à gaz} (Kribi, alimentée par le gaz de Sanaga Sud ; Logbaba à Douala) et les \textbf{centrales thermiques au fuel} (approvisionnées via la SONARA) assurent l'appoint, à coût élevé.
\item Le \textbf{solaire} se développe dans le Grand Nord (Guider, Maroua, projets ruraux).
\item Paradoxe central : potentiel énorme, mais production insuffisante, délestages et accès inégal — moins de 10 \% du potentiel hydroélectrique est exploité.
\end{itemize}
\end{aretenir}

\courssub{TP 3 : Construction d'un camembert de la production d'énergie}

\begin{experience}[TP 3 : Représenter le mix électrique camerounais par un diagramme circulaire]
\textbf{Objectif.} Construire un diagramme circulaire (camembert) montrant la répartition de la production d'électricité au Cameroun par source (hydroélectricité, gaz, thermique fuel, solaire) à partir de données statistiques récentes.

\textbf{Document d'appui : Tableau de la production électrique (estimation 2023, ordre de grandeur).}
\begin{center}
\begin{tabularx}{\linewidth}{|l|X|c|c|}
\hline
\rowcolor{popblueL} \textbf{Source d'énergie} & \textbf{Type de centrale} & \textbf{Production (GWh)} & \textbf{Part (\%)} \\
\hline
Hydroélectricité & Barrages (Édéa, Songloulou, Lagdo, Memve'ele, Nachtigal) & 5 200 & ? \\
\hline
Gaz naturel & Kribi, Logbaba & 1 800 & ? \\
\hline
Thermique fuel & Centrales d'appoint (Douala, Yaoundé, Nord, etc.) & 1 000 & ? \\
\hline
Solaire & Guider, Maroua, mini-réseaux ruraux & 100 & ? \\
\hline
\textbf{Total} & & \textbf{8 100} & \textbf{100 \%} \\
\hline
\end{tabularx}
\end{center}

\textbf{Matériel.} Feuille blanche, compas, rapporteur, règle, crayons de couleur (bleu, orange, rouge, jaune), calculatrice.

\textbf{Protocole.}
\begin{enumerate}
\item \textbf{Calculer les pourcentages.} Pour chaque source, appliquer la formule : $\frac{\text{Production source}}{\text{Production totale}} \times 100$. Remplir la colonne « Part (\%) » du tableau.
\item \textbf{Calculer les angles.} Un cercle complet vaut $360^{\circ}$. Pour chaque source : $\text{Angle} = \frac{\text{Part (\%)}}{100} \times 360^{\circ}$. Noter les angles calculés.
\item \textbf{Tracer le cercle.} Au compas, tracer un cercle de rayon 5 cm (ou 6 cm) au centre de la feuille.
\item \textbf{Repérer le centre et le rayon de départ.} Tracer un rayon horizontal vers la droite (angle $0^{\circ}$).
\item \textbf{Construire les secteurs.} À l'aide du rapporteur, construire successivement les secteurs dans l'ordre du tableau (hydroélectricité, gaz, fuel, solaire) en mesurant les angles calculés depuis le rayon précédent. Colorier chaque secteur avec le code couleur convenu.
\item \textbf{Légender et titrer.} Rédiger une légende complète (couleur = source). Donner un titre précis : \emph{« Répartition de la production d'électricité au Cameroun par source en 2023 (estimation) »}. Indiquer la source des données (ex. : estimations ARSEL / ENEO / MINEE).
\end{enumerate}

\textbf{Résultats attendus (corrigé des calculs).}
\begin{align*}
\text{Hydroélectricité : } & \frac{5200}{8100} \times 100 \approx 64,2\% \quad \rightarrow \quad 0,642 \times 360^{\circ} \approx 231^{\circ} \\
\text{Gaz naturel : } & \frac{1800}{8100} \times 100 \approx 22,2\% \quad \rightarrow \quad 0,222 \times 360^{\circ} \approx 80^{\circ} \\
\text{Thermique fuel : } & \frac{1000}{8100} \times 100 \approx 12,3\% \quad \rightarrow \quad 0,123 \times 360^{\circ} \approx 44^{\circ} \\
\text{Solaire : } & \frac{100}{8100} \times 100 \approx 1,2\% \quad \rightarrow \quad 0,012 \times 360^{\circ} \approx 4^{\circ} \\
\text{Somme des angles : } & 231^{\circ} + 80^{\circ} + 44^{\circ} + 4^{\circ} = 359^{\circ} \text{ (arrondi acceptable, ajuster le dernier secteur à } 360^{\circ}).
\end{align*}

\textbf{Interprétation guidée.}
\begin{itemize}
\item Le secteur \textbf{bleu (hydroélectricité)} domine largement (près des 2/3 du cercle) : confirmation de la prédominance de l'hydraulique.
\item Le secteur \textbf{orange (gaz)} est le deuxième contributeur significatif (environ 1/5) : rôle croissant du gaz pour la stabilité du réseau.
\item Le secteur \textbf{rouge (fuel)} reste non négligeable (environ 1/8) : coût économique et écologique élevé, à réduire.
\item Le secteur \textbf{jaune (solaire)} est encore très mince (environ 1 \%) : marge de progression immense, surtout dans le Nord.
\end{itemize}

\begin{attention}[Pièges du camembert au Bac]
\begin{itemize}
\item Ne pas oublier de \textbf{colorier} les secteurs et de faire correspondre \textbf{exactement} les couleurs du graphique et de la légende.
\item Vérifier que la \textbf{somme des angles vaut 360°} (ajuster le dernier secteur si arrondis).
\item Le \textbf{titre} doit contenir : le thème (production d'électricité), l'espace (Cameroun), la date (2023) et le mode de représentation (diagramme circulaire / camembert).
\item Ne pas confondre \textbf{puissance installée (MW)} et \textbf{production effective (GWh)}. Le sujet porte sur la production.
\end{itemize}
\end{attention}
\end{experience}

\textbf{Transition.} Ce potentiel énergétique, s'il est mieux exploité, doit alimenter une industrie qui a aussi besoin de matières premières. La section suivante montre que le sous-sol camerounais recèle des réserves minières de premier plan, encore largement inexploitées.

\coursec{D'importantes réserves minières}

Dans la section précédente, nous avons découvert que le Cameroun dispose d'un énorme potentiel énergétique : ses fleuves alimentent des barrages hydroélectriques, son gaz naturel fait tourner des centrales, son soleil attend d'être capté. Mais une industrie ne se contente pas d'énergie : il lui faut aussi de la \textbf{matière première}. Or, sous le sol camerounais, dort un véritable trésor : de la bauxite, du fer, des terres rares, de l'or, du diamant, du pétrole, du gaz\ldots{} Le paradoxe, c'est que ce trésor reste largement inexploité. C'est ce que nous allons comprendre dans cette section.

\courssub{La bauxite : un gisement géant encore endormi}

\begin{definition}[Bauxite et minerai]
La \cle{bauxite} est une roche sédimentaire riche en alumine (oxyde d'aluminium), à partir de laquelle on produit l'aluminium. Un \cle{minerai} est une roche dont on peut extraire un métal de manière rentable. Une \cle{réserve} est la quantité de minerai identifiée et exploitable dans un gisement.
\end{definition}

Imagine un gisement si grand qu'il figure parmi les plus importants d'Afrique, et pourtant presque intact. C'est exactement le cas de la bauxite de l'Adamaoua. Les gisements de \cle{Minim-Martap} et de \cle{Ngaoundal}, situés sur le plateau de l'Adamaoua, sont estimés à plus d'un milliard de tonnes de minerai (ordre de grandeur communément cité). Leur teneur en alumine est élevée (souvent supérieure à 50 \%), ce qui les rend particulièrement intéressants.

Pourtant, cette bauxite n'est quasiment pas exploitée. Pourquoi ? La réponse tient en un mot : l'\cle{enclavement}. L'Adamaoua est un plateau situé à plus de \SI{600}{km} de l'océan. La bauxite est un minerai lourd et peu cher à la tonne : son transport par camion jusqu'au port de Douala coûterait plus cher que le minerai lui-même. Il faudrait prolonger la voie ferrée (Transcam) de Ngaoundéré vers Minim-Martap, puis jusqu'à un port — un investissement de plusieurs milliards de dollars. Tant que cette infrastructure n'existe pas, le gisement reste « dormant ».

\begin{savaistu}[Un milliard de tonnes qui attendent un train]
Le gisement de bauxite de Minim-Martap est estimé à plus d'un milliard de tonnes, ce qui en fait l'un des plus grands gisements africains. Mais son exploitation attend la construction d'une voie ferrée de plusieurs centaines de kilomètres reliant le gisement au littoral. C'est l'illustration parfaite d'une règle de géographie minière : \emph{un gisement sans voie de transport n'est qu'une richesse potentielle}.
\end{savaistu}

\courssub{Le fer : le grand projet de Mbalam}

Dans la région de l'Est, près de la frontière congolaise, se trouve le gisement de fer de \cle{Mbalam}. Le fer est le métal de base de l'industrie mondiale : il sert à fabriquer l'acier, indispensable au bâtiment, aux machines, aux voies ferrées elles-mêmes.

\begin{exemplebox}[Le fer de Mbalam en chiffres]
Le gisement de Mbalam (et son extension Nabeba côté congolais) est estimé à plusieurs centaines de millions de tonnes de minerai de fer, avec une teneur élevée (souvent supérieure à 60 \% de fer pour les minerais riches, ce qui est remarquable). Le projet d'exploitation prévoit la construction d'une voie ferrée d'environ \SI{500}{km} jusqu'au \cle{port minéralier de Kribi}, en eau profonde, capable d'accueillir de grands navires vraquiers. Le projet, porté successivement par Sundance Resources puis AustSino, a connu plusieurs reports, faute de financement définitif et de levée de capitaux suffisante.
\end{exemplebox}

Le projet Mbalam illustre la logique d'\cle{intégration nationale} : un minerai de l'Est, une voie ferrée neuve, un port sur l'Atlantique — autant d'infrastructures qui, une fois construites, profiteraient à tout le pays et pas seulement à la mine.

\courssub{Les terres rares et les autres minerais stratégiques}

\begin{definition}[Les terres rares]
Les \cle{terres rares} sont un groupe de 17 métaux (lanthane, cérium, néodyme, dysprosium\ldots) indispensables aux technologies modernes : aimants permanents des moteurs électriques et des éoliennes, écrans de téléphones, batteries, fibres optiques. Malgré leur nom, elles ne sont pas si rares dans l'écorce terrestre, mais elles sont difficiles et coûteuses à séparer. Leur contrôle est un enjeu géopolitique majeur, la Chine dominant la production mondiale.
\end{definition}

Le sous-sol camerounais recèle bien d'autres richesses, inégalement explorées :

\begin{itemize}
\item le \cle{cobalt} et le \cle{nickel} dans la région de l'Est, gisement de \cle{Nkamouna} (proche de Lomié), métaux stratégiques pour les batteries lithium-ion ; des teneurs en terres rares y sont associées ;
\item le \cle{manganèse} et le \cle{rutile} (minerai de titane), notamment sur le littoral (région de Kribi, Bipindi) ;
\item le \cle{diamant} à l'Extrême-Nord (région de \cle{Mobilong}, près de la frontière tchadienne) ;
\item l'\cle{or}, exploité de manière artisanale dans l'Est (Bétaré-Oya, Batouri) et l'Adamaoua, souvent dans l'informel ;
\item l'\cle{uranium}, dont des indices existent dans le nord du pays (région de \cle{Poli}, Kitongo) et au sud (Lolodorf), non exploités.
\end{itemize}

\courssub{Le pétrole et le gaz : le lien avec l'énergie}

Nous l'avons vu dans la section précédente : le Cameroun est un producteur modeste de pétrole et de gaz. Rappelons l'essentiel. Le pétrole est extrait principalement dans le \cle{bassin du Rio del Rey}, au large de la région du Sud-Ouest (golfe de Guinée), par la Société nationale des hydrocarbures (SNH) et ses partenaires (Perenco, Bowleven, etc.) ; le brut est raffiné à la \cle{SONARA} de Limbé. Le gaz, lui, est exploité notamment à \cle{Logbaba} (périphérie de Douala) et dans le bassin de Kribi : il alimente via gazoduc les centrales à gaz comme celle de \cle{Kribi} (Kribi Power Development Company, \SI{216}{MW}) et de \cle{Dibamba} (Douala). Le bassin de \cle{Logone Birni}, à l'Extrême-Nord, recèle aussi des hydrocarbures, peu exploités en raison de l'enclavement et de l'insécurité. La production pétrolière nationale, qui avait culminé dans les années 1980 (environ \SI{180 000}{barils/jour}), a fortement décliné depuis (ordre de grandeur actuel : \SI{30 000}{à 40 000 barils/jour}).

\begin{popfigure}
\begin{center}
\includegraphics[width=\linewidth,height=9.5cm,keepaspectratio]{N01/cmr_map.jpg}
\end{center}
\legende{Fond de carte du Cameroun sur lequel situer les ressources minières et les hydrocarbures : bauxite de Minim-Martap et de Ngaoundal (Adamaoua), fer de Mbalam, cobalt-nickel-terres rares de Nkamouna (Est), pétrole du Rio del Rey au large du Sud-Ouest, gaz de Logbaba près de Douala, diamant à l'Extrême-Nord (Mobilong), or dans l'Est. La voie ferrée (Transcamerounais) visible sur la carte est celle que prolongeraient les projets Mbalam--Kribi et vers la bauxite. Les gisements ne figurent pas sur l'image : ils sont à reporter par l'élève. \\ Source : Wikimedia Commons, CIA, domaine public.}
\end{popfigure}

\courssub{Un tableau statistique à analyser}

\begin{center}
\begin{tabularx}{\linewidth}{|l|X|X|X|}
\hline
\rowcolor{popblueL} \textbf{Ressource} & \textbf{Localisation principale} & \textbf{Réserves estimées (ordres de grandeur)} & \textbf{État d'exploitation} \\
\hline
Bauxite & Minim-Martap, Ngaoundal (Adamaoua) & plus d'un milliard de tonnes & quasi inexploitée (attente voie ferrée) \\
\hline
Fer & Mbalam (Est) & plusieurs centaines de millions de tonnes, teneur > 60 \% & projet (voie ferrée + port de Kribi) \\
\hline
Cobalt, Nickel, Terres rares & Nkamouna / Lomié (Est) & gisement important, mal évalué & projet (études de faisabilité) \\
\hline
Diamant & Mobilong (Extrême-Nord) & modeste & artisanale et semi-industrielle \\
\hline
Or & Est (Bétaré-Oya), Adamaoua & mal connue & artisanale, largement informelle \\
\hline
Uranium & Poli (Nord), Lolodorf (Sud) & indices, non évaluées & inexploitée \\
\hline
Pétrole & Rio del Rey (offshore Sud-Ouest) & en déclin & exploitée (SNH, partenaires, SONARA) \\
\hline
Gaz naturel & Logbaba (Douala), Kribi, Logone Birni & significative & exploitée (centrales thermiques) \\
\hline
\end{tabularx}
\end{center}

\textbf{Lecture du tableau.} Trois enseignements se dégagent. Premièrement, les ressources sont \emph{diversifiées} et réparties sur tout le territoire : aucune région n'est totalement dépourvue. Deuxièmement, les gisements les plus importants en volume (bauxite, fer) sont aussi les \emph{moins exploités}, car ils exigent des infrastructures lourdes (chemins de fer, ports minéraliers). Troisièmement, seules les ressources faciles d'accès ou à forte valeur unitaire (pétrole, gaz, or, diamant) sont réellement extraites — parfois dans l'informel, ce qui échappe aux recettes de l'État.

\courssub{Le paradoxe minier camerounais}

\begin{attention}[Réserve potentielle $\neq$ richesse réelle]
Une erreur fréquente au Baccalauréat consiste à écrire que « le Cameroun est riche parce qu'il possède de nombreux minerais ». Faux : une \cle{réserve minière} n'est une richesse réelle que si elle est \textbf{exploitée} et, idéalement, \textbf{transformée sur place}. Tant qu'un gisement reste « dormant », il ne crée ni emplois, ni recettes fiscales, ni devises. Il faut donc toujours distinguer le \emph{potentiel} (ce qui existe dans le sous-sol) de la \emph{production effective} (ce qui est réellement extrait).
\end{attention}

Pourquoi tant de gisements restent-ils dormants ? Quatre obstacles se combinent :

\begin{enumerate}
\item \textbf{Les coûts d'infrastructure} : bauxite et fer sont des minerais de masse, peu chers à la tonne ; ils ne sont rentables qu'avec un chemin de fer et un port dédié, soit des investissements de plusieurs milliards de dollars.
\item \textbf{L'enclavement} : les grands gisements (Adamaoua, Est) sont loin des côtes et mal desservis par les routes actuelles.
\item \textbf{Le financement} : l'État camerounais ne peut financer seul ces projets ; il dépend d'investisseurs étrangers, sensibles à la chute des cours mondiaux des matières premières.
\item \textbf{L'instabilité des cours} : un projet rentable quand le fer est cher peut être abandonné quand les prix s'effondrent, ce qui décourage les capitaux sur le long terme.
\end{enumerate}

\courssub{Des ressources minières à l'industrialisation : l'enjeu de la transformation locale}

Dernière étape du raisonnement : que faire de ces minerais ? Deux stratégies s'opposent.

\begin{itemize}
\item \textbf{L'exportation brute} : on extrait le minerai et on le vend tel quel à l'étranger. C'est rapide, mais le pays ne gagne que la valeur de la matière première — la plus faible de la filière.
\item \textbf{La transformation locale} : on transforme le minerai sur place (bauxite $\rightarrow$ alumine $\rightarrow$ aluminium ; fer $\rightarrow$ acier). Chaque étape ajoute de la \cle{valeur ajoutée}, crée des emplois qualifiés et développe l'industrie nationale.
\end{itemize}

Le Cameroun possède déjà un exemple historique : \cle{ALUCAM} à Edéa transforme l'alumine en aluminium grâce à l'électricité des barrages d'Edéa et de Song Loulou — mais cette alumine est \emph{importée} (principalement de Guinée), alors que le pays possède sa propre bauxite ! Voilà le sommet du paradoxe : transformer localement un minerai\ldots{} venu d'ailleurs. L'enjeu des prochaines décennies est donc clair : relier les gisements aux usines et aux ports, pour que le sous-sol camerounais serve enfin l'industrie camerounaise.

\begin{exoresolu}[Analyser le paradoxe minier camerounais]
\textbf{Énoncé.} À partir du tableau des ressources minières et de tes connaissances : (1) dégage deux caractéristiques des ressources minières du Cameroun ; (2) explique pourquoi les gisements de bauxite de l'Adamaoua ne sont pas exploités ; (3) montre en quoi la transformation locale des minerais est un enjeu d'intégration nationale.
\tcblower
\textbf{Solution détaillée.}

(1) \emph{Caractéristiques.} Les ressources minières du Cameroun sont d'une part \textbf{abondantes et diversifiées} (bauxite, fer, cobalt, nickel, terres rares, diamant, or, pétrole, gaz), réparties sur l'ensemble du territoire ; d'autre part \textbf{très inégalement exploitées} : seuls le pétrole, le gaz et l'or artisanal sont réellement extraits, tandis que les plus grands gisements (bauxite de Minim-Martap, fer de Mbalam, cobalt-nickel de Nkamouna) restent au stade de projet.

(2) \emph{Non-exploitation de la bauxite.} Les gisements de Minim-Martap et Ngaoundal se trouvent sur le plateau de l'Adamaoua, à plus de \SI{600}{km} de l'océan. La bauxite étant un minerai lourd et de faible valeur à la tonne, son transport routier jusqu'à Douala coûterait plus cher que le minerai lui-même. L'exploitation exige donc au préalable la construction d'une voie ferrée (prolongement du Transcam) et d'un accès portuaire, investissements colossaux que l'État ne peut financer seul et que les investisseurs étrangers hésitent à engager sans visibilité sur les cours.

(3) \emph{Enjeu d'intégration nationale.} Transformer les minerais sur place (alumine puis aluminium, acier) augmenterait la valeur ajoutée conservée au Cameroun, créerait des emplois industriels dans les régions de l'Adamaoua et de l'Est, justifierait la construction de voies ferrées et du port minéralier de Kribi — infrastructures profitant à tout le pays (désenclavement, transport des personnes, autres frets) — et réduirait la dépendance aux importations. Le cas d'ALUCAM, qui transforme de l'alumine importée alors que la bauxite nationale dort, montre le chemin qui reste à parcourir pour une véritable intégration verticale de la filière.
\end{exoresolu}

\begin{aretenir}[L'essentiel de la section]
Le Cameroun possède d'importantes réserves minières : bauxite de Minim-Martap et Ngaoundal (plus d'un milliard de tonnes, parmi les plus grands gisements d'Afrique), fer de Mbalam (plusieurs centaines de millions de tonnes à forte teneur), cobalt, nickel, terres rares (Nkamouna), diamant, or, uranium, pétrole (Rio del Rey) et gaz (Logbaba, Kribi, Logone Birni). Mais ces richesses restent largement \textbf{dormantes} à cause de l'enclavement, du coût des infrastructures et des difficultés de financement. Une réserve n'est une richesse réelle que si elle est exploitée \emph{et transformée localement} : c'est là tout l'enjeu de l'industrialisation et de l'intégration nationale du Cameroun.
\end{aretenir}

\coursec{Un élan industriel brisé par la crise}

Dans la section précédente, nous avons vu que le sous-sol camerounais recèle d'importantes réserves minières (bauxite à Ngaoundal et Minim-Martap, fer à Mbalam, terres rares) encore largement inexploitées. Or, posséder des matières premières ne suffit pas : il faut les transformer sur place pour créer de la richesse et des emplois. C'est précisément l'ambition que le Cameroun indépendant s'est donnée dans les années 1960. Mais cette belle histoire industrielle a connu un coup d'arrêt brutal. Pour le comprendre, partons d'une situation concrète : à Douala-Bonabéri, un élève qui se promène le long de la zone industrielle remarque des hangars désaffectés, des cheminées froides, des enseignes rouillées d'entreprises qui n'existent plus, alors que son grand-père lui raconte qu'« avant, tout le monde travaillait à l'usine ». Que s'est-il passé ? C'est toute la question de cette section : comment un élan industriel prometteur a-t-il été brisé par la crise ?

\courssub{L'élan industriel après l'indépendance (années 1960-1980)}

\textbf{L'intuition.} Au moment des indépendances, le Cameroun, comme la plupart des pays africains, exporte des produits bruts (cacao, café, coton, bananes, bois, aluminium) et importe presque tous ses produits manufacturés (tissus, savon, ciment, boissons, matériels agricoles, équipements). Cette situation est injuste et coûteuse : les matières premières brutes se vendent peu cher sur le marché mondial, tandis que les produits transformés s'achètent très cher. L'idée des dirigeants est simple : \emph{transformer sur place ce que le pays produit, et produire sur place ce que le pays consomme}.

\begin{definition}[Industrialisation par substitution aux importations]
L'\cle{industrialisation} est le processus par lequel un pays développe ses industries de transformation pour fabriquer lui-même les biens qu'il importait auparavant. Au Cameroun, elle a pris la forme d'une politique volontariste de création d'\cle{unités industrielles} et de \cle{sociétés d'État} (entreprises publiques ou parapubliques) financées par l'État et des capitaux étrangers.
\end{definition}

\textbf{Les moyens de cette politique.} Dès les années 1960, puis surtout sous l'effet du boom pétrolier (le Cameroun devient producteur de pétrole en 1977, avec la \cle{SONARA} créée en 1973 et la raffinerie de Limbe entrée en service en 1981), l'État camerounais investit massivement. Il crée ou développe des dizaines de sociétés d'État et d'unités de production. Citons les plus emblématiques :

\begin{itemize}
\item \cle{ALUCAM} (Compagnie camerounaise de l'aluminium) : installée à \cle{Édéa} (Littoral), elle transforme l'alumine importée en aluminium grâce à l'électricité abondante et bon marché du barrage hydroélectrique d'Édéa sur la Sanaga. Créée en 1957, elle devient le fleuron de l'industrie camerounaise.
\item \cle{CIMENCAM} (Cimenteries du Cameroun) : usine intégrée de \cle{Figuil} (Nord, calcaire local) dès 1966, puis broyeur de clinker à \cle{Douala-Bonabéri} ; elle produit le ciment indispensable à la construction.
\item \cle{SOSUCAM} (Société sucrière du Cameroun) : créée dans les années 1970, elle gère les grandes plantations et usines de \cle{Mbandjock} et de \cle{Nkoteng} (Centre).
\item \cle{SODECOTON} (Société de développement du coton) : créée en 1974, elle organise la filière cotonnière du Nord (égrenage, huilerie de Kaele) et devient l'un des premiers employeurs du pays.
\item \cle{CAMSUCO} (Cameroon Sugar Company) : ancienne appellation de l'unité de Mbandjock, souvent citée comme pilier historique de la filière sucre aux côtés de SOSUCAM.
\item On peut ajouter les usines textiles (\cle{CICAM} à Douala et Garoua), les brasseries (\cle{SABC}), la chocolaterie \cle{CHOCOCAM}, les chantiers navals (\cle{CAMSHIP} à Douala), les savonneries, etc.
\end{itemize}

\textbf{Les résultats.} À la fin des années 1970 et au début des années 1980, le Cameroun est présenté comme l'un des pays les plus industrialisés d'Afrique centrale. L'industrie manufacturière représente alors une part notable du PIB (de l'ordre de 20 à 25 \% au milieu des années 1980, ordre de grandeur à retenir avec prudence), et la zone industrielle de Douala-Bonabéri emploie des dizaines de milliers de salariés. Le pays exporte même de l'aluminium, du ciment et des produits agro-industriels vers les pays voisins de la CEMAC.

\begin{savaistu}[ALUCAM, le géant d'Édéa]
Créée en 1957, avant même l'indépendance, \cle{ALUCAM} fut longtemps la plus grande entreprise industrielle du Cameroun. Son secret ? L'électrolyse de l'alumine exige d'énormes quantités d'électricité : l'usine s'est donc installée à côté du barrage d'Édéa, qui lui fournissait une énergie abondante et bon marché. À son apogée, elle produisait autour de 80 000 à 90 000 tonnes d'aluminium par an. Mais dépendante de l'alumine importée (la bauxite camerounaise de Ngaoundal et Minim-Martap n'étant pas exploitée), elle a été fortement ralentie par la crise, puis par les difficultés d'approvisionnement en alumine et les coupures d'électricité. En 2014, l'activité électrolytique de la société historique a été arrêtée, avant une reprise partielle sous nouvelle structure (Alucam SA) dans les années suivantes.
\end{savaistu}

\courssub{La crise économique des années 1980-1990 : un coup d'arrêt brutal}

\textbf{L'intuition.} Imagine un commerçant dont toute la richesse repose sur la vente du cacao et du pétrole. Si le prix mondial de ces produits s'effondre de moitié, ses revenus s'effondrent aussi, alors que ses dettes, elles, restent entières. C'est exactement ce qui arrive au Cameroun à partir de 1985-1986.

\textbf{Les causes externes : la chute des cours.} À partir de 1986, les cours mondiaux des principaux produits d'exportation du Cameroun (pétrole, cacao, café, coton) chutent brutalement. Les recettes d'exportation s'effondrent, le déficit budgétaire se creuse et l'État s'\cle{endette} massivement auprès des créanciers étrangers (Club de Paris, Club de Londres, institutions financières internationales).

\textbf{Les causes internes.} La crise révèle aussi des faiblesses propres au pays : mauvaise gestion de certaines sociétés d'État (coûts élevés, sur-effectifs, faible productivité, détournements de fonds), dépenses publiques excessives pendant les années de boom pétrolier (grands travaux, équipements), corruption. Les entreprises publiques, habituées à être soutenues par l'État sans contrainte de rentabilité, se révèlent peu compétitives face au marché.

\begin{definition}[L'ajustement structurel]
L'\cle{ajustement structurel} désigne l'ensemble des politiques économiques imposées par le \cle{FMI} (Fonds monétaire international) et la \cle{Banque mondiale} aux pays endettés, en échange de prêts de soutien. Il comprend notamment : la \cle{privatisation} ou la liquidation des entreprises publiques déficitaires, la réduction des dépenses publiques (baisse des salaires, gel des recrutements, licenciements dans la fonction publique), la libéralisation du commerce (ouverture aux importations, suppression de barrières tarifaires) et la dévaluation de la monnaie. Le Cameroun y est soumis à partir de 1988 (Premier Programme d'Ajustement Structurel), avec la dévaluation du franc CFA de 50 \% en janvier 1994.
\end{definition}

\courssub{Les conséquences : fermetures, privatisations et licenciements}

Les programmes d'ajustement structurel (PAS) ont des effets sociaux et industriels considérables. L'État, sommé de réduire ses dépenses, cesse de soutenir ses entreprises déficitaires. Trois sorts les attendent :

\begin{enumerate}
\item \textbf{La liquidation pure et simple} : l'entreprise ferme définitivement. C'est le cas des chantiers navals \cle{CAMSHIP}, de la biscuiterie \cle{DOBAF}, de nombreuses unités textiles (comme celles de la CICAM à Garoua et Douala) et de bien d'autres PME.
\item \textbf{La privatisation} : l'entreprise est vendue à des capitaux privés, souvent étrangers. Ainsi, des pans entiers de l'industrie passent sous contrôle de groupes internationaux (ex. : brasseries, hôtels, plantations).
\item \textbf{La survie difficile} : quelques entreprises stratégiques (\cle{SODECOTON}, \cle{SOSUCAM}, \cle{CIMENCAM}, \cle{ALUCAM}) traversent la crise, mais affaiblies, avec des investissements reportés, des effectifs réduits et des équipements vieillissants.
\end{enumerate}

\begin{exemplebox}[Une hécatombe industrielle chiffrée]
Entre 1988 et 2000, \textbf{des dizaines d'entreprises publiques et parapubliques} camerounaises du secteur industriel ont été liquidées ou privatisées dans le cadre des programmes d'ajustement structurel (on parle couramment d'une centaine de sociétés concernées au total, tous secteurs confondus ; retiens l'ordre de grandeur pour l'industrie : plusieurs dizaines). Ces opérations se sont accompagnées de \textbf{licenciements massifs} : des dizaines de milliers de salariés de l'industrie et de la fonction publique ont perdu leur emploi, tandis que les salaires étaient fortement réduits (jusqu'à 60 à 70 \% de baisse dans la fonction publique en 1993). Des quartiers entiers de Douala, Édéa ou Garoua, qui vivaient de l'usine voisine, ont été plongés dans le chômage et la pauvreté.
\end{exemplebox}

\textbf{La concurrence des produits importés.} La libéralisation des échanges, imposée par l'ajustement structurel, ouvre le marché camerounais aux produits étrangers : tissus et vêtements asiatiques, chaussures, appareils électroménagers, produits alimentaires transformés, friperie (vêtements d'occasion dits « \cle{bend down select} »). Ces produits, souvent moins chers car produits à grande échelle dans des pays à bas coûts, concurrencent directement la production locale. Les usines textiles camerounaises, par exemple, n'ont pas résisté à l'arrivée massive des pagnes et vêtements importés. Face à cette concurrence, certaines entreprises ferment, d'autres \cle{délocalisent} une partie de leurs activités vers des pays où la main-d'œuvre est moins chère ou l'énergie plus fiable.

\begin{attention}[Ne pas attribuer la crise à une seule cause]
Piège classique du baccalauréat : expliquer la crise industrielle camerounaise par une cause unique (« c'est la faute du FMI » ou « c'est la faute de la chute du pétrole »). En réalité, la crise résulte de la \textbf{conjonction} de facteurs \textbf{externes} (chute des cours mondiaux des matières premières, endettement international, conditions imposées par les bailleurs de fonds) et de facteurs \textbf{internes} (mauvaise gestion des entreprises publiques, sur-effectifs, corruption, manque d'entretien des équipements, dépendance aux importations de biens intermédiaires comme l'alumine pour ALUCAM). Dans une dissertation ou un commentaire de document, cite toujours les deux registres de causes.
\end{attention}

\courssub{Lire la crise dans le paysage : frises et friches industrielles}

\begin{popfigure}
\begin{center}
\begin{tikzpicture}[font=\small]
% Axe chronologique
\draw[-{Stealth[length=3mm]},very thick,popdark] (0,0) -- (15.5,0) node[below left=1pt and 0pt]{};
% Graduations
\foreach \x/\t in {0.5/1960, 4.5/1975, 7.5/1985, 10/1994, 12/2000, 15/2020}{
  \draw[popdark] (\x,0.12) -- (\x,-0.12) node[below]{\t};
}
% Période 1 : élan industriel
\fill[popgreenL] (0.5,0.35) rectangle (7.5,1.5);
\node[align=center,text=popdark] at (4,0.92){\textbf{1960-1985 : l'élan industriel}\\ ALUCAM, CIMENCAM, SOSUCAM,\\ SODECOTON, SONARA (1981)};
% Période 2 : crise
\fill[poporangeL] (7.5,0.35) rectangle (12,1.5);
\node[align=center,text=popdark] at (9.75,0.92){\textbf{1986-2000 : la crise}\\ chute des cours, PAS,\\ fermetures (CAMSHIP, DOBAF)};
% Période 3 : reprise timide
\fill[popblueL] (12,0.35) rectangle (15,1.5);
\node[align=center,text=popdark] at (13.5,0.92){\textbf{Depuis 2000 :}\\ reprise timide};
% Jalons
\node[above=2pt,text=popgreen!60!black] at (2,1.5){créations d'unités};
\node[above=2pt,text=poporange!70!black] at (8.6,1.5){1988 : PAS};
\node[above=2pt,text=poporange!70!black] at (10.9,1.5){1994 : dévaluation};
\node[above=2pt,text=popblue!70!black] at (13.5,1.5){nouveaux investisseurs};
\end{tikzpicture}
\end{center}
\legende{Frise chronologique : les trois grandes phases de l'industrie camerounaise depuis l'indépendance. En vert, la phase de construction (1960-1985) ; en orange, la crise et les fermetures (1986-2000) ; en bleu, la reprise timide depuis les années 2000 (arrivée de nouveaux opérateurs comme Dangote Cement, relances agro-industrielles).}
\end{popfigure}

\begin{popfigure}
\begin{center}
\begin{tikzpicture}[font=\small]
% Sol
\draw[thick,popdark] (0,0) -- (13,0);
% Bâtiment principal (hall d'électrolyse)
\draw[fill=popblueL,thick] (1,0) rectangle (7,2.2);
\draw[fill=popblue!40,thick] (1,2.2) -- (1.5,2.9) -- (6.5,2.9) -- (7,2.2) -- cycle;
% Cheminées
\draw[fill=popmuted!50,thick] (2,2.9) rectangle (2.6,4.6);
\draw[fill=popmuted!50,thick] (3.2,2.9) rectangle (3.8,4.2);
% Fumée faible
\draw[popmuted,decorate,decoration={coil,segment length=4pt,amplitude=2pt}] (2.3,4.6) -- (2.3,5.4);
% Hangar désaffecté
\draw[fill=poporangeL,thick,dashed] (8,0) rectangle (11.5,1.6);
\draw[fill=poporange!30,thick,dashed] (8,1.6) -- (9.75,2.4) -- (11.5,1.6);
% Herbes (friche)
\foreach \x in {8.4,9.2,10.1,11}{
  \draw[popgreen!70!black] (\x,0) -- (\x-0.08,0.35) (\x,0) -- (\x+0.08,0.35) (\x,0) -- (\x,0.4);
}
% Lignes électriques
\draw[thick,popdark] (12.3,0) -- (12.3,3.6);
\draw[thick,popdark] (11.9,3.4) -- (12.7,3.4);
\draw[popdark] (12.3,3.6) -- (12.05,3.4) (12.3,3.6) -- (12.55,3.4);
% Annotations
\draw[-{Stealth},popblue] (4,2.9) -- (4,3.9) node[above,align=center]{1. halls d'électrolyse\\ (production d'aluminium)};
\draw[-{Stealth},popmuted] (2.3,4.3) -- (0.6,4.3) node[left,align=right]{2. cheminées\\ (activité ralentie)};
\draw[-{Stealth},poporange!80!black] (9.75,0.8) -- (9.75,-0.7) node[below,align=center]{3. friche industrielle :\\ hangar désaffecté, végétation};
\draw[-{Stealth},popdark] (12.3,2.5) -- (13.2,2.5) node[right,align=left]{4. ligne haute tension\\ (barrage d'Édéa)};
\end{tikzpicture}
\end{center}
\legende{Photographie commentée (reconstitution schématique) : le site industriel d'ALUCAM à Édéa. 1. Les halls d'électrolyse, cœur de l'usine, qui transformaient l'alumine importée en aluminium. 2. Les cheminées, témoins d'une activité longtemps ralentie par la crise et les difficultés d'approvisionnement. 3. Une friche industrielle voisine : bâtiment abandonné repris par la végétation, paysage typique de la désindustrialisation. 4. La ligne haute tension rappelle le lien historique entre l'usine et le barrage hydroélectrique d'Édéa.}
\end{popfigure}

\courssub{Bilan : un tissu industriel affaibli et une reprise timide}

Au tournant des années 2000, le bilan est lourd. Le Cameroun a connu une \cle{désindustrialisation} partielle : la part de l'industrie manufacturière dans le PIB a reculé (autour de 15-18 \% selon les années), le tissu industriel est affaibli, dominé désormais par les agro-industries et les PME, tandis que les industries lourdes sont devenues minoritaires. De nombreux savoir-faire ont été perdus, et le pays est devenu plus dépendant des importations de produits manufacturés.

Cependant, depuis les années 2000, on observe une \cle{reprise timide} : arrivée de nouveaux investisseurs (le groupe nigérian \cle{Dangote Cement} s'installe à Douala en 2015, devenant un acteur majeur du ciment), modernisation de certaines filières (ciment, sucre, coton), projets de transformation locale liés aux grands travaux (barrages de Lom Pangar, Mékin, Nachtigal ; port en eau profonde de Kribi). Mais cette reprise reste fragile : elle dépend de la fiabilité de l'énergie, de la qualité des infrastructures de transport et de la stabilité macroéconomique.

\begin{exoresolu}[Expliquer la crise industrielle camerounaise]
\textbf{Énoncé.} À partir de tes connaissances, explique en une dizaine de lignes pourquoi l'industrie camerounaise, dynamique dans les années 1970, a connu une crise profonde à la fin des années 1980. Tu mobiliseras au moins quatre idées organisées.
\tcblower
\textbf{Solution détaillée.} L'industrie camerounaise s'est fortement développée après l'indépendance grâce à une politique volontariste de création de sociétés d'État (ALUCAM, CIMENCAM, SOSUCAM, SODECOTON), soutenue par les recettes pétrolières après 1978. Mais à partir de 1986, la chute brutale des cours mondiaux des matières premières (pétrole, cacao, café) fait s'effondrer les recettes de l'État, qui s'endette lourdement. Cette crise externe révèle des faiblesses internes : mauvaise gestion, sur-effectifs et faible compétitivité des entreprises publiques. Pour obtenir des prêts, le Cameroun accepte à partir de 1988 les programmes d'ajustement structurel du FMI et de la Banque mondiale, qui imposent privatisations, liquidations et réduction des dépenses publiques. Il en résulte la fermeture de nombreuses unités industrielles (CAMSHIP, DOBAF, usines textiles), des licenciements massifs et une concurrence accrue des produits importés favorisée par la libéralisation des échanges. La crise industrielle est donc le produit d'une conjonction de facteurs externes et internes, et non d'une cause unique.
\end{exoresolu}

\begin{aretenir}[L'essentiel de la section]
\begin{itemize}
\item Après l'indépendance, le Cameroun mène une politique d'industrialisation fondée sur les sociétés d'État : ALUCAM (Édéa), CIMENCAM (Figuil, Douala), SOSUCAM (Mbandjock, Nkoteng), SODECOTON (Nord), CAMSUCO.
\item À partir de 1986, la chute des cours des matières premières et l'endettement provoquent une crise économique majeure.
\item Les programmes d'ajustement structurel (FMI, Banque mondiale, dès 1988) imposent privatisations et liquidations : fermeture de nombreuses unités (CAMSHIP, DOBAF, textiles), licenciements massifs.
\item La libéralisation expose les industries locales à la concurrence des produits importés ; certaines activités sont délocalisées.
\item Bilan : désindustrialisation partielle, tissu industriel affaibli, reprise timide depuis les années 2000 (Dangote Cement, relances agro-industrielles, grands barrages).
\end{itemize}
\end{aretenir}

\begin{exercice}[Pour t'entraîner]
\begin{enumerate}
\item Définis en deux phrases l'ajustement structurel et cite deux de ses mesures appliquées au Cameroun.
\item Sur la frise chronologique de la leçon, situe : la création de la SODECOTON, l'entrée en service de la SONARA, le début des programmes d'ajustement structurel, la dévaluation du franc CFA.
\item Rédige un court commentaire (10 lignes) de la photographie commentée du site d'Édéa, en montrant ce que le paysage industriel révèle de l'histoire économique du Cameroun.
\end{enumerate}
\end{exercice}

\coursec{Généralités et définitions}

L'industrie est l'ensemble des activités économiques qui transforment des matières premières en produits finis ou semi-finis grâce à l'énergie et au travail. Elle constitue le \cle{secteur secondaire} de l'économie. Au Cameroun, comme dans tout pays en développement, l'industrialisation est un levier majeur de l'\cle{intégration nationale} : elle permet de valoriser les ressources locales, de créer des emplois non agricoles, de réduire les importations et de souder le territoire par des flux de matières et de produits finis.

\begin{definition}[Unité industrielle]
Une \cle{unité industrielle} est un établissement où s'opère une transformation physique ou chimique de matières premières à l'aide de machines et d'énergie, de façon continue ou répétée, en vue de la vente. Elle se distingue de l'artisanat par l'usage de la force motrice (électricité, thermique) et d'une organisation du travail salarié.
\end{definition}

\begin{definition}[Facteurs de localisation industrielle]
Les facteurs classiques qui expliquent l'implantation d'une usine sont :
\begin{itemize}
\item les \textbf{matières premières} (proximité du gisement ou du port d'importation) ;
\item l'\textbf{énergie} (hydroélectricité, gaz, fioul) ;
\item la \textbf{main-d'œuvre} (disponibilité, qualification, coût) ;
\item le \textbf{marché} (taille, pouvoir d'achat, accessibilité) ;
\item les \textbf{infrastructures de transport} (route, rail, port, pipeline) ;
\item la \textbf{politique publique} (zones franches, incitations fiscales, décentralisation).
\end{itemize}
\end{definition}

\begin{methode}[Analyser la localisation d'une usine]
Pour commenter l'implantation d'une unité industrielle (ex. ALUCAM à Édéa, SONARA à Limbé, CIMENCAM à Figuil) :
\begin{enumerate}
\item Identifier l'activité et le processus technique (électro-intensif, matière première encombrante, etc.).
\item Rechercher les facteurs \textbf{spécifiques} qui ont prévalu : énergie bon marché (ALUCAM), port en eau profonde (SONARA), gisement calcaire (Figuil).
\item Évoquer les facteurs \textbf{généraux} : main-d'œuvre, marché, transports, État.
\item Conclure sur le poids de chaque facteur et les éventuelles contraintes actuelles (vétusté, coûts, insécurité).
\end{enumerate}
\end{methode}

\coursec{Un énorme potentiel énergétique}

L'énergie est le \textbf{sang de l'industrie}. Le Cameroun dispose d'atouts exceptionnels, encore sous-exploités, qui devraient en faire une plaque tournante énergétique en Afrique centrale.

\courssub{L'hydroélectricité : la colonne vertébrale}

Le fleuve \textbf{Sanaga} (900 km, débit moyen \SI{2000}{m^3/s} à Édéa) concentre l'essentiel du potentiel hydroélectrique national estimé à \textbf{20 000 MW} (ordre de grandeur), soit le 2\up{ème} potentiel africain après la RDC. Seuls \textbf{5 à 10 \%} sont actuellement exploités.

\begin{itemize}
\item \textbf{Barrage d'Édéa} (1953, 263 MW) : historique, alimente ALUCAM et le réseau Sud (Douala, Yaoundé).
\item \textbf{Barrage de Song Loulou} (1988, 384 MW) : amont d'Édéa, régularise le débit.
\item \textbf{Barrage de Lagdo} (1982, 72 MW) : sur la Bénoué (Nord), régulation du fleuve, irrigation, pêche.
\item \textbf{Barrage de Memve'ele} (mis en service progressif 2019-2024, 211 MW) : sur la Ntem (Sud), renforce l'interconnexion Sud.
\item \textbf{Barrage de Nachtigal} (mis en service 2023-2024, 420 MW) : sur la Sanaga, projet phare en PPP (Nachtigal Hydro Power Company), doit couvrir \SI{\sim 30}{\%} de la demande nationale.
\item \textbf{Projets structurants} : Grand Eweng (800-1200 MW), Lom Pangar (réservoir de régulation 6 milliards m³ + 30 MW), Bini à Warak (75 MW).
\end{itemize}

\begin{savaistu}[Le régime de la Sanaga]
La Sanaga a un régime \textbf{pluvial équatorial de transition} : une crue principale (sept.-nov.), une crue secondaire (mars-juin), deux étiages. L'étiage sévère (janv.-fév.) limite la production d'Édéa et Song Loulou. Le barrage-réservoir de \textbf{Lom Pangar} (achèvement 2016) stocke l'eau de crue pour la restituer en étiage, sécurisant ainsi la production en aval (+ 700 GWh/an estimés).
\end{savaistu}

\courssub{Le gaz naturel : le relais thermique}

Le Cameroun possède des réserves prouvées de gaz naturel estimées à \textbf{170 milliards m³} (ordre de grandeur, offshore Rio del Rey, Logbaba, Sanaga Sud). Ce gaz alimente des centrales à cycle combiné (plus efficaces, moins polluantes que le fioul) :

\begin{itemize}
\item \textbf{Kribi} (Kribi Power Development, 216 MW, gaz de Logbaba via pipeline) ;
\item \textbf{Logbaba} (Victoria Oil \& Gas, 82 MW, gaz local) ;
\item \textbf{Dibamba} (88 MW, fioul lourd — \emph{ancienne centrale de secours, en conversion gaz}) ;
\item \textbf{Limbé} (centrale thermique d'appoint, fioul/gaz).
\end{itemize}

\courssub{Le solaire et les autres renouvelables : un avenir à construire}

L'ensoleillement est fort au Nord (\SI{> 2500}{kWh/m^2/an}) et modéré au Sud. Des centrales photovoltaïques voient le jour : \textbf{Maroua} (20 MWc, 2023), \textbf{Guider} (20 MWc), \textbf{Bertoua} (projet). Le potentiel éolien est localisé (Adamawa, littoral). La biomasse (résidus agricoles, scieries) alimente quelques cogénérations industrielles. La part du solaire dans le \textbf{mix électrique interconnecté (RI)} reste \textbf{inférieure à 1 \%} (ordre de grandeur 2023-2024), mais les projets se multiplient (objectif 25 \% EnR hors hydro à l'horizon 2030 selon la SNEE).

\begin{aretenir}[Mix électrique camerounais (RI, ordres de grandeur 2023-2024)]
\begin{itemize}
\item \textbf{Hydroélectricité} : \textbf{\SI{\sim 72}{\%}} (Édéa, Song Loulou, Lagdo, Memve'ele, Nachtigal).
\item \textbf{Thermique à gaz} : \textbf{\SI{\sim 23}{\%}} (Kribi, Logbaba).
\item \textbf{Thermique fioul/diesel (appoint/location)} : \textbf{\SI{\sim 4}{\%}}.
\item \textbf{Solaire / Autres EnR} : \textbf{\SI{< 1}{\%}}.
\end{itemize}
La demande croît de \SI{\sim 6-8}{\%} par an. Le défi : achever les grands barrages, développer le gaz domestique et intégrer le solaire pour baisser le coût moyen du kWh industriel.
\end{aretenir}

\coursec{D'importantes réserves minières}

Le sous-sol camerounais recèle des ressources de \textbf{rang mondial}, mais leur exploitation industrielle tarde, freinée par l'infrastructure (rail, port, énergie), le cadre juridique et la gouvernance.

\begin{center}
\begin{tabularx}{\linewidth}{|l|X|l|X|}
\hline
\rowcolor{popblueL} \textbf{Ressource} & \textbf{Gisements principaux} & \textbf{Réserves / Potentiel (ordres de grandeur)} & \textbf{Statut / Enjeu} \\
\hline
\textbf{Bauxite} & Minim-Martap, Ngaoundal (Adamaoua) & \SI{\sim 1,2}{milliard} de tonnes (l'un des plus gros gisements mondiaux non exploités) & \textbf{Non exploité}. Nécessite chemin de fer (600 km) vers Kribi + aluminerie locale. Projet \textbf{Caminco / Hydromines} en discussion. \\
\hline
\textbf{Fer} & Mbalam (Est, frontalier Congo), Nkout (Est) & Mbalam : \SI{\sim 800}{Mt} @ 60 \% Fe ; Nkout : \SI{\sim 500}{Mt} & \textbf{Non exploité}. Projet Sundance / AustSino (Mbalam) : rail 500 km vers port en eau profonde de Kribi. Bloqué par contentieux et financement. \\
\hline
\textbf{Terres rares} & Akwa (Nord-Ouest), Lombo (Est) & Potentiel significatif (néodyme, dysprosium...), exploration avancée & \textbf{Exploration}. Stratégiques pour transition énergétique (aimants permanents). \\
\hline
\textbf{Or} & Bétaré-Oya (Est), Batouri, Kambele & Production artisanale \SI{\sim 1-2}{t/an} ; réserves industrielles non publiées & \textbf{Artisanal dominant}. Projet industriel (ex. Bétaré-Oya) en étude. Enjeu : traçabilité, recettes fiscales. \\
\hline
\textbf{Diamant} & Mobilong, Bétaré-Oya (Est) & Gisements alluvionnaires & \textbf{Artisanal}. Certification Kimberley difficile. \\
\hline
\textbf{Calcaire} & Figuil (Nord), Pouma (Littoral), Mfome (Sud) & Exploité pour ciment (Figuil, Pouma) & \textbf{Exploité} (CIMENCAM, DANGOTE). \\
\hline
\textbf{Cobalt, Nickel, Manganèse} & Nkamouna (Est), Mada (Nord) & Associés aux latérites nickelifères / cobaltifères & \textbf{Exploration / Études}. Projet Nkamouna (Cobalt/Ni/Mn) avancé. \\
\hline
\end{tabularx}
\end{center}

\begin{attention}[Potentiel vs Production]
Au Bac, ne jamais écrire : « Le Cameroun produit de la bauxite » ou « Le Cameroun exporte du fer ». Distinguer \textbf{réserves/gisements} (potentiel géologique) et \textbf{production/exportation} (réalité économique). Seuls le calcaire (ciment), le sable/granulats (BTP) et l'or/diamant (artisanal) sont actuellement produits. Le pétrole (offshore Rio del Rey, Logbaba) est en déclin de production (\SI{\sim 30-35}{000} barils/jour en 2023 vs \SI{185}{000} au pic 1985).
\end{attention}

\coursec{Un élan industriel brisé par la crise}

\courssub{Les « Trente Glorieuses » camerounaises (1960-1986)}

Après l'indépendance, l'État volontariste construit un appareil productif diversifié : \textbf{CICAM} (textile, Garoua), \textbf{ALUCAM} (aluminium), \textbf{SODECOTON}, \textbf{SOSUCAM}, \textbf{CDC}, \textbf{SONARA}, \textbf{CIMENCAM}, \textbf{HEVECAM}, \textbf{SPFS} (pâtes/papiers, Édéa), \textbf{CAMSUCO}, brasseries, savonneries, usines de chaussures, montage automobile (PEUGEOT, MERCEDES à Douala). L'industrie atteint \textbf{25-30 \% du PIB} (vers 1985), l'emploi industriel formel dépasse \num{100 000} salariés.

\courssub{La crise structurelle (1986-2000) : causes et conséquences}

\begin{itemize}
\item \textbf{Chocs exogènes} : effondrement des cours du cacao, café, coton, pétrole (1986) ; appréciation du franc CFA (parité fixe avec le franc français) qui renchérit les exportations.
\item \textbf{Déséquilibres endogènes} : surendettement public, gabegie, entreprises parapubliques surdimensionnées et déficitaires, coût du facteur travail élevé, infrastructures dégradées.
\item \textbf{PAS (Plans d'Ajustement Structurel) FMI/Banque Mondiale (1988-1994)} : dévaluation du FCFA (\SI{50}{\%}, janv. 1994), privatisations massives, suppression des subventions, ouverture brutale aux importations (libéralisation des échanges).
\end{itemize}

\begin{exemplebox}[Conséquences visibles : la « désindustrialisation »]
\begin{itemize}
\item \textbf{Fermetures emblématiques} : CICAM (textile, 3000 emplois), SPFS (papiers), CAMSUCO (sucre, repris par SOSUCAM), usines de chaussures (BATA, CAMSHOE), montage auto, nombreuses PME agroalimentaires.
\item \textbf{Chute de la part de l'industrie dans le PIB} : de \SI{\sim 28}{\%} (1986) à \SI{\sim 15-17}{\%} (2000), autour de \SI{17-19}{\%} aujourd'hui (hors pétrole).
\item \textbf{Perte d'emplois industriels formels} : divisés par 2 ou 3 en une décennie ; report massif vers le secteur informel.
\item \textbf{Désarticulation des filières} : coton sans filature, bois brut exporté sans transformation, cacao brut exporté à 80-90 \%.
\end{itemize}
\end{exemplebox}

\begin{propriete}[Le paradoxe camerounais]
Le Cameroun possède les \textbf{facteurs de base} de l'industrialisation (énergie, mines, agriculture, marché sous-régional CEMAC \SI{\sim 60}{millions} d'habitants), mais son tissu industriel reste \textbf{fragile, extraverti et spatialement concentré}. La libéralisation des échanges (Module 2) a accentué la concurrence des importations (riz, huile, textile, clinker, produits finis chinois/européens) sur des PME mal armées.
\end{propriete}

\coursec{Les types d'industries au Cameroun}

La crise des années 1986-2000 a fermé de nombreuses usines, mais elle n'a pas fait disparaître l'industrie camerounaise : elle en a transformé le visage. Aujourd'hui, quand tu traverses la zone industrielle de Douala-Bassa ou que tu passes devant l'usine ALUCAM à Édéa, tu observes des réalités très différentes : d'immenses complexes capitalistiques d'un côté, des centaines de petits ateliers de l'autre. Comment classer ce foisonnement ? C'est toute l'ambition de cette section : comprendre que l'appareil productif camerounais repose sur trois piliers inégaux — des \cle{industries lourdes} minoritaires, des \cle{agro-industries} dominantes, et un tissu dense de \cle{PME} — auxquels s'ajoutent l'artisanat et le secteur informel.

\courssub{Les industries lourdes : des géants minoritaires}

\textbf{Intuition.} Une industrie est dite \og lourde \fg{} quand il faut des milliards de francs CFA, des machines gigantesques et beaucoup d'énergie pour transformer des minerais ou des matériaux bruts. On ne monte pas une cimenterie dans sa cour comme on monte un atelier de couture !

\begin{definition}[Industrie lourde]
Une \cle{industrie lourde} est une unité de production qui mobilise des capitaux très importants, des équipements massifs et une forte consommation d'énergie pour transformer des \textbf{minerais} ou des \textbf{matériaux de base} (aluminium, ciment, produits pétroliers raffinés, acier). Elle emploie souvent peu de main-d'œuvre par rapport au capital investi : on dit qu'elle est \emph{capitalistique}.
\end{definition}

Au Cameroun, les industries lourdes se comptent sur les doigts de la main. Elles sont \textbf{minoritaires} en nombre, mais stratégiques par leur poids économique et symbolique :

\begin{itemize}
\item \textbf{ALUCAM} (Compagnie Camerounaise de l'Aluminium), installée à \textbf{Édéa} depuis 1957 : elle transforme l'alumine importée (majoritairement de Guinée) en aluminium grâce à l'électricité abondante et bon marché du barrage d'Édéa sur la Sanaga. C'est l'archétype de l'industrie \og électro-intensive \fg{}. Capacité : \SI{\sim 95 000}{t/an} (ordre de grandeur).
\item \textbf{SONARA} (Société Nationale de Raffinage), à \textbf{Limbé} : raffinerie de pétrole mise en service en 1981, capacité nominale \SI{2,1}{Mt/an} (brut). Elle traite le brut (en partie importé depuis la fin de la production locale) pour produire carburants, kérosène, bitume. Gravement endommagée par un incendie en mai 2019, sa remise en marche complète reste un enjeu national majeur (souveraineté énergétique, balance commerciale).
\item \textbf{Les cimenteries} :
      \begin{itemize}
      \item \textbf{CIMENCAM} (Cimenteries du Cameroun, groupe LafargeHolcim) : usine intégrée de \textbf{Figuil} (Nord, calcaire local, \SI{\sim 500 000}{t/an}), usine de broyage à \textbf{Douala-Bonabéri} (\SI{\sim 1}{Mt/an}), usine de \textbf{Nkongsamba} (broyage).
      \item \textbf{DANGOTE Cement} (groupe nigérian) : usine de broyage à \textbf{Douala} (mise en service 2015, \SI{\sim 1,5}{Mt/an}), usine intégrée à \textbf{Figuil} (en construction/extension, exploitation calcaire local).
      \end{itemize}
      La concurrence entre cimentiers a fait baisser le prix du sac de ciment : un bel exemple des effets de la libéralisation.
\end{itemize}

\begin{attention}[Piège fréquent au Bac]
Ne cite pas ALUCAM comme preuve que le Cameroun exploite sa bauxite ! L'usine d'Édéa importe l'essentiel de son alumine : la bauxite de Minim-Martap et de Ngaoundal (Adamaoua) n'est \textbf{pas encore exploitée} industriellement. Confondre potentiel minier et production effective est une erreur classique.
\end{attention}

\courssub{La prédominance des agro-industries}

\textbf{Intuition.} Le Cameroun est un pays agricole : il est donc logique que la majorité de ses usines transforment ce que la terre produit — coton, canne à sucre, huile de palme, cacao, hévéa, bananes, céréales.

\begin{definition}[Agro-industrie]
Une \cle{agro-industrie} est une entreprise industrielle qui transforme des \textbf{produits agricoles ou d'élevage} en biens de consommation ou en produits semi-finis : huilerie (huile de palme), sucrerie (canne), brasserie (céréales), filature et huilerie de coton, savonnerie, chocolaterie.
\end{definition}

Les grandes agro-industries camerounaises structurent des filières entières :

\begin{itemize}
\item \textbf{SODECOTON} (Société de Développement du Coton), basée à \textbf{Garoua} : égrenage, huilerie de graines de coton, savonnerie ; encadre \SI{\sim 200 000}{producteurs} (ordre de grandeur) dans le Nord. C'est l'une des plus grandes entreprises du pays (chiffre d'affaires, effectifs).
\item \textbf{SOSUCAM} (Société Sucrière du Cameroun, groupe Somdiaa) : sucreries de \textbf{Nkoteng} et \textbf{Mbandjock} (Centre), alimentées par de vastes plantations de canne (\SI{\sim 30 000}{ha} au total). Production \SI{\sim 150 000}{t/an} de sucre (ordre de grandeur), couvre l'essentiel de la consommation nationale.
\item \textbf{SOCAPALM} (Société Camerounaise de Palmeraies, groupe Socfin) et \textbf{PHP} (Plantations des Palmiers, groupe Bolloré) : huileries de palme dans le Littoral (Dibombari, Eseka, Mbongo) et le Sud-Ouest. Production huile brute \SI{\sim 150 000}{t/an} (cumulé, ordre de grandeur).
\item \textbf{CDC} (Cameroon Development Corporation), dans le Sud-Ouest (Bota, Tiko, Muyuka) : bananes (export), palme, hévéa — entreprise parapublique historique, fortement fragilisée par la crise anglophone depuis 2017 (perte de surfaces, destruction d'infrastructures).
\item \textbf{HEVECAM} (Sud, Niete) et \textbf{SAFACAM} (Littoral, Dizangué) : hévéa (caoutchouc naturel).
\item \textbf{Les brasseries SABC} (Société Anonyme des Brasseries du Cameroun, groupe Castel) : bières et boissons gazeuses, avec des unités à Douala, Yaoundé, Bafoussam et Garoua — un exemple d'ancrage national et de maillage territorial.
\item S'y ajoutent les \textbf{savonneries} (Savonnerie du Cameroun, etc.), les \textbf{minoteries} (Grand Moulins du Cameroun, etc.) et les unités de transformation du cacao (chocolaterie \textbf{Chococam} à Douala, usine de transformation de fèves à Kribi/Yaoundé).
\end{itemize}

\begin{exemplebox}[Pourquoi les agro-industries dominent-elles ?]
Trois raisons se combinent : (1) la matière première est \textbf{locale et abondante} (coton, canne, palme, cacao) ; (2) le marché national est \textbf{garanti} (sucre, huile, bière, savon sont des biens de consommation courante à demande peu élastique) ; (3) ces industries exigent des capitaux \textbf{moins lourds} qu'une raffinerie ou une usine d'aluminium. L'agro-industrie est donc la voie naturelle d'industrialisation d'un pays agricole.
\end{exemplebox}

\courssub{Les PME : l'essentiel du tissu industriel}

\begin{definition}[PME]
Une \cle{PME} (petite et moyenne entreprise) emploie un effectif limité (de quelques salariés à quelques centaines, seuil officiel \SI{< 200}{salariés} au Cameroun) et mobilise des capitaux modestes. Au Cameroun, les PME constituent l'\textbf{essentiel du tissu industriel} : elles sont de loin les plus nombreuses et les principales pourvoyeuses d'emplois industriels formels.
\end{definition}

Les PME camerounaises sont extrêmement diversifiées : transformation agroalimentaire (minoteries artisanales, jus de fruits, conserves, transformation manioc/plantain), scieries et menuiseries industrielles (filière \textbf{bois} : sciage, placage, contreplaqué — \SI{\sim 80}{\%} du bois scié local), matériaux de construction et \textbf{bâtiment} (briqueteries, préfabrication, peinture), métallurgie légère (portes, fenêtres, charpentes), plasturgie (emballages, tuyaux), imprimerie, textile/habillement. Leur force : la proximité avec les marchés locaux et la souplesse. Leur faiblesse : l'accès difficile au crédit (taux \SI{> 10}{\%}), le coût de l'énergie (coupures, tarifs), la concurrence des importations (riz, huile, textile, clinker) et la fiscalité perçue comme lourde.

\begin{attention}[Ne pas réduire l'industrie aux grandes usines]
Erreur fréquente en dissertation : décrire l'industrie camerounaise uniquement à travers ALUCAM ou la SONARA. Or les \textbf{PME et l'agro-industrie en sont la composante dominante}, en nombre d'unités comme en emplois. Une copie qui ignore les PME est une copie déséquilibrée.
\end{attention}

\courssub{Le secteur informel et l'artisanat : un complément indispensable}

Autour du secteur moderne gravite un immense \cle{secteur informel} : menuisiers, soudeurs, tailleurs, cordonniers, transformateurs artisanaux de produits alimentaires (huile rouge de quartier, manioc transformé en bobolo/miondo/garri, séchage/fumage du poisson), garages et fonderies artisanales — le célèbre quartier \og Madagascar \fg{} à Douala en est un symbole. Non enregistrées, ces activités échappent aux statistiques officielles (INS), mais elles fournissent des biens accessibles aux populations, forment la main-d'œuvre et assurent une part considérable de l'emploi urbain (estimé à \SI{> 90}{\%} de l'emploi total hors agriculture). L'artisanat est ainsi un \textbf{complément de l'appareil productif}, et parfois un tremplin vers la PME formelle (dispositifs d'appui : PME-PMI, ANOR, guichets uniques).

\begin{methode}[Classer une entreprise : deux critères croisés]
Pour classer correctement une unité industrielle dans une copie ou un exercice :
\begin{enumerate}
\item Identifier la \textbf{nature de la production} : transformation de minerais/matériaux (industrie lourde), de produits agricoles (agro-industrie), ou fabrication légère de biens divers (PME).
\item Évaluer la \textbf{taille} : montant des capitaux investis, effectifs, capacité de production.
\item Croiser les deux critères : une huilerie de village est une agro-industrie de type PME ; ALUCAM est une industrie lourde de grande taille.
\item Conclure en situant l'entreprise dans le paysage national (dominance des agro-industries et des PME).
\end{enumerate}
\end{methode}

\coursec{Les zones industrielles et les perspectives}

\courssub{Une répartition spatiale très concentrée}

L'industrie camerounaise est géographiquement \textbf{déséquilibrée}. \textbf{Douala} concentre plus de la moitié des unités industrielles du pays : son \textbf{port} (première porte d'entrée des matières premières et des biens d'équipement, \SI{\sim 95}{\%} du trafic portuaire national), sa \textbf{zone industrielle} de Bassa-Bonabéri (plus de \SI{300}{ha}, \SI{\sim 400}{entreprises}), son marché de consommateurs (\SI{\sim 4}{millions} d'habitants aire urbaine) et ses infrastructures expliquent cette hyper-concentration. Autour de ce pôle majeur s'organisent des \textbf{pôles secondaires} : \textbf{Yaoundé} (marché administratif \SI{\sim 4}{millions} hab., brasserie, PME, pharmacie), \textbf{Bafoussam} (Ouest, agroalimentaire, brasserie, menuiserie), \textbf{Garoua} (Nord, SODECOTON, brasserie, porte du Sahel), \textbf{Limbé} (Sud-Ouest, SONARA, port, agro-industrie), \textbf{Édéa} (ALUCAM, énergie), \textbf{Figuil} (cimenterie, calcaire) et \textbf{Ngaoundéré} (agroalimentaire, carrefour ferroviaire Transcam, porte de l'Adamaoua).

\begin{exemplebox}[Douala, capitale industrielle]
Grâce à son port en eau profonde sur le Wouri et à sa zone industrielle, Douala rassemble à elle seule \textbf{plus de la moitié des unités industrielles} du Cameroun (ordre de grandeur couramment admis : \SI{55}{\%} à \SI{60}{\%}). Cette concentration pose un vrai problème d'\textbf{intégration nationale} : les régions de l'intérieur, pourtant riches en matières premières (coton, bois, minerais, bétail), restent peu industrialisées. Le corridor Douala-Ngaoundéré (route, rail, pipeline, fibre) structure l'essentiel de l'espace productif.
\end{exemplebox}

\begin{popfigure}
\begin{center}
\begin{tikzpicture}[scale=1]
% Camembert : Agro-industries 45%, PME 40%, Industries lourdes 15%
\fill[popgreenL] (0,0) -- (90:2.6) arc (90:-72:2.6) -- cycle;
\fill[poporangeL] (0,0) -- (-72:2.6) arc (-72:-216:2.6) -- cycle;
\fill[popblue] (0,0) -- (-216:2.6) arc (-216:-270:2.6) -- cycle;
\draw[popdark] (0,0) circle (2.6);
\node at (20:1.6) {\small\bfseries 45 \%};
\node at (-145:1.6) {\small\bfseries 40 \%};
\node at (-243:1.7) {\small\bfseries 15 \%};
% Légende
\fill[popgreenL] (3.4,1.6) rectangle (3.9,2.0); \node[right,align=left] at (4.0,1.8) {\small Agro-industries (SODECOTON, SOSUCAM, SABC...)};
\fill[poporangeL] (3.4,0.6) rectangle (3.9,1.0); \node[right,align=left] at (4.0,0.8) {\small PME (agroalimentaire, bois, bâtiment...)};
\fill[popblue] (3.4,-0.4) rectangle (3.9,0.0); \node[right,align=left] at (4.0,-0.2) {\small Industries lourdes (ALUCAM, SONARA, cimenteries)};
\end{tikzpicture}
\end{center}
\legende{Répartition approximative des unités industrielles camerounaises par type (ordres de grandeur). Les agro-industries et les PME forment ensemble l'écrasante majorité du tissu industriel ; les industries lourdes, très capitalistiques, restent minoritaires.}
\end{popfigure}

\begin{popfigure}
\begin{center}
\includegraphics[width=\linewidth,height=9.5cm,keepaspectratio]{N01/cmr_map.jpg}
\end{center}
\legende{Fond de carte pour localiser les grandes unités industrielles du Cameroun. Le littoral (Douala, Édéa, Limbé, Kribi) concentre l'essentiel de l'appareil productif ; les pôles de l'intérieur (Bafoussam, Yaoundé, Ngaoundéré, Garoua, Figuil) restent secondaires. On remarque sur la carte que routes et voies ferrées convergent vers Douala et Yaoundé. \\ Source : Wikimedia Commons, CIA, domaine public.}
\end{popfigure}

\courssub{Bilan comparatif : trois mondes industriels}

Le tableau suivant synthétise les trois composantes de l'industrie camerounaise. Apprends-le : il fournit une grille de lecture immédiatement mobilisable dans une dissertation.

\begin{center}
\begin{tabularx}{\linewidth}{|l|X|X|X|}
\hline
\rowcolor{popblueL} \textbf{Critère} & \textbf{Industries lourdes} & \textbf{Agro-industries} & \textbf{PME} \\
\hline
Capitaux & Très élevés (milliards de FCFA) & Moyens à élevés & Modestes \\
\hline
Emplois & Limités (industrie capitalistique) & Importants (filières rurales) & Très nombreux au total (premier pourvoyeur d'emplois industriels formels) \\
\hline
Production & Aluminium, ciment, produits raffinés & Sucre, huile, coton, boissons, savon & Biens de consommation divers, bois, matériaux \\
\hline
Exemples & ALUCAM, SONARA, CIMENCAM, DANGOTE & SODECOTON, SOSUCAM, SOCAPALM, CDC, PHP, SABC & Minoteries, scieries, briqueteries, plasturgie \\
\hline
Poids & Minoritaires en nombre & Dominantes avec les PME & Écrasante majorité des unités \\
\hline
\end{tabularx}
\end{center}

\courssub{Perspectives : réindustrialisation et intégration nationale}

Le \textbf{Plan National de Développement 2020-2030 (PND30)} et la \textbf{Stratégie Nationale de Développement (SND30)} visent une \textbf{réindustrialisation} par :
\begin{enumerate}
\item La \textbf{transformation locale} des matières premières (interdiction progressive export bois en grumes, cacao brut, coton brut ; projets aluminerie, usine fer, engrais).
\item Le développement de \textbf{zones économiques spéciales (ZES)} et \textbf{parcs industriels} : Kribi (port en eau profonde + zone industrielle), Limbé, Ngaoundéré, Maroua, Bertoua.
\item L'amélioration de la \textbf{compétitivité} : énergie fiable et moins chère (Nachtigal, gaz), formation professionnelle adaptée, accès au financement (Banque de Développement, fonds de garantie), allègement fiscal (régime d'incitation à l'investissement).
\item L'\textbf{intégration sous-régionale} (ZLECAf, CEMAC) : élargir le marché pour atteindre la taille critique.
\end{enumerate}

\begin{aretenir}[L'essentiel de la leçon]
L'industrie camerounaise repose sur trois piliers inégaux : des \cle{industries lourdes} minoritaires mais stratégiques (ALUCAM à Édéa, SONARA à Limbé, cimenteries CIMENCAM et DANGOTE), des \cle{agro-industries} dominantes adossées aux filières agricoles (SODECOTON, SOSUCAM, SOCAPALM, CDC, SABC), et un tissu dense de \cle{PME}, principales pourvoyeuses d'emplois industriels, complété par l'artisanat et le secteur informel. Spatialement, \textbf{Douala concentre plus de la moitié des unités industrielles}, devant les pôles secondaires de Yaoundé, Bafoussam, Garoua et Limbé : un déséquilibre qui interpelle directement l'intégration nationale. Le pays dispose d'un \textbf{énorme potentiel énergétique} (hydro, gaz, solaire) et de \textbf{réserves minières de rang mondial} (bauxite, fer, terres rares) encore inexploitées, atouts majeurs pour une future réindustrialisation endogène.
\end{aretenir}

\coursec{TP 3 : Construction d'un camembert de la production d'énergie}

\begin{experience}[TP 3 : Représenter le mix électrique camerounais]
\textbf{Objectif.} Construire un diagramme circulaire (camembert) représentant la structure de la production d'électricité au Cameroun (Réseau Interconnecté) à partir de données statistiques récentes.

\textbf{Matériel.} Feuille blanche, compas, rapporteur, règle, crayons de couleur (bleu, orange, rouge, jaune), calculatrice.

\textbf{Données statistiques (Ordres de grandeur 2023-2024, sources : ARSEL, MINEE, ENEO, Nachtigal Hydro Power Company).}
\begin{center}
\begin{tabularx}{\linewidth}{|l|X|c|c|}
\hline
\rowcolor{popblueL} \textbf{Source} & \textbf{Centrales principales} & \textbf{Production annuelle (GWh, approx.)} & \textbf{Part (\%)} \\
\hline
Hydroélectricité & Édéa, Song Loulou, Lagdo, Memve'ele, Nachtigal & \SI{\sim 5 500}{} & \textbf{72 \%} \\
\hline
Thermique à gaz & Kribi (KPDC), Logbaba & \SI{\sim 1 750}{} & \textbf{23 \%} \\
\hline
Thermique fioul/diesel & Dibamba, Limbé, groupes locatifs d'urgence & \SI{\sim 300}{} & \textbf{4 \%} \\
\hline
Solaire / Autres EnR & Maroua, Guider, Bertoua (projets) & \SI{\sim 50}{} & \textbf{1 \%} \\
\hline
\textbf{Total} & & \textbf{\SI{\sim 7 600}{}} & \textbf{100 \%} \\
\hline
\end{tabularx}
\end{center}

\textbf{Protocole de construction.}
\begin{enumerate}
\item \textbf{Calculer les angles.} La somme des angles = 360°. Angle = Part (\%) $\times$ 3,6.
      \begin{align*}
      \text{Hydro} &: 72 \times 3,6 = 259,2^\circ \approx 259^\circ \\
      \text{Gaz}   &: 23 \times 3,6 = 82,8^\circ \approx 83^\circ \\
      \text{Fioul} &: 4 \times 3,6 = 14,4^\circ \approx 14^\circ \\
      \text{Solaire} &: 1 \times 3,6 = 3,6^\circ \approx 4^\circ \\
      \text{Total} &: 259 + 83 + 14 + 4 = 360^\circ \quad (\text{OK})
      \end{align*}
\item \textbf{Tracer le cercle.} Au compas, rayon \SI{5}{cm} (ou \SI{6}{cm}) au centre de la feuille.
\item \textbf{Tracer le premier secteur.} Tracer un rayon horizontal vers la droite (0°). À l'aide du rapporteur, mesurer \textbf{259°} dans le sens inverse des aiguilles d'une montre (trigonométrique). Tracer le second rayon. Colorier en \textbf{bleu} (hydroélectricité).
\item \textbf{Tracer les secteurs suivants.} Depuis le dernier rayon, mesurer \textbf{83°} (Gaz $\rightarrow$ \textbf{orange}), puis \textbf{14°} (Fioul $\rightarrow$ \textbf{rouge}), le reste \textbf{4°} (Solaire $\rightarrow$ \textbf{jaune}).
\item \textbf{Légender.} Écrire les pourcentages à l'intérieur des secteurs (ou à l'extérieur avec traits de rappel). Faire une légende encadrée en bas ou à droite : Bleu = Hydro (72 \%), Orange = Gaz (23 \%), Rouge = Fioul (4 \%), Jaune = Solaire/EnR (1 \%). Indiquer le titre : \og Structure de la production d'électricité au Cameroun (RI, 2023-2024) \fg, la source et l'échelle (1 cm = 10 \% ou 36°).
\end{enumerate}

\textbf{Interprétation attendue (à rédiger sous le camembert).}
\begin{itemize}
\item \textbf{Dominance hydroélectrique (72 \%)} : atout majeur (énergie renouvelable, bas coût marginal, faible émission CO₂), mais vulnérabilité climatique (étiages, variabilité interannuelle).
\item \textbf{Place croissante du gaz (23 \%)} : flexibilité, appoint en étiage, valorisation ressource locale (Logbaba), mais dépendance au prix du gaz et aux infrastructures (pipeline).
\item \textbf{Résiduel fioul (4 \%)} : coût très élevé, pollution, vocation à disparaître avec l'arrivée du gaz et des barrages.
\item \textbf{Solaire naissant (1 \%)} : potentiel énorme au Nord, mais intermittence et coût d'investissement encore freins. Objectif politique : 25 \% EnR hors hydro à 2030.
\end{itemize}

\end{experience}

\begin{attention}[Consignes de rédaction pour le TP]
\begin{itemize}
\item Le cercle doit être \textbf{propre}, au compas, les rayons à la règle.
\item Les angles doivent être \textbf{mesurés au rapporteur} (pas « à l'œil »).
\item Les couleurs doivent être \textbf{distinctes} et conformes au code couleur habituel (Bleu = eau, Orange/Jaune = soleil, Rouge = feu/thermique).
\item La \textbf{légende est obligatoire} (titre, source, auteur, date).
\item L'\textbf{interprétation} doit mobiliser les connaissances du cours (potentiel Sanaga, gaz Logbaba, enjeu coût/environnement).
\end{itemize}
\end{attention}

\begin{exoresolu}[Sujet type Bac : Analyse du mix énergétique]
\textbf{Énoncé.} À l'aide du camembert construit et de vos connaissances, montrez que le mix électrique camerounais présente à la fois des atouts pour l'industrialisation et des vulnérabilités structurelles.
\tcblower
\textbf{Solution détaillée.}
\begin{enumerate}
\item \textbf{Atouts majeurs.}
      \begin{itemize}
      \item \textbf{Renouvelable dominant (72 \% hydro)} : électricité bas carbone, coût marginal quasi nul après amortissement, avantage comparatif pour industries électro-intensives (ALUCAM, futures alumineries, aciéries).
      \item \textbf{Ressources locales diversifiées} : eau (Sanaga), gaz (Logbaba, Kribi), soleil (Nord). Réduit la dépendance aux importations d'hydrocarbures (contrairement à beaucoup de pays africains).
      \item \textbf{Projets structurants en cours} : Nachtigal (420 MW), Memve'ele (211 MW), Lom Pangar (régulation), conversion centrales fioul $\rightarrow$ gaz, solaire Maroua/Guider. Capacité installée en forte hausse.
      \end{itemize}
\item \textbf{Vulnérabilités structurelles.}
      \begin{itemize}
      \item \textbf{Dépendance climatique} : 72 \% hydro = risque de délestage en année sèche (ex. 2015-2016, 2021). Lom Pangar atténue mais ne supprime pas le risque.
      \item \textbf{Infrastructures de transport vieillissantes} : réseau de transport (SONATREL) et distribution (ENEO) saturés, pertes techniques et non techniques \SI{\sim 15-20}{\%}, pannes fréquentes pénalisant l'outil industriel.
      \item \textbf{Coût pour l'industrie} : tarifs industriels encore élevés (moyenne \SI{\sim 60-70}{FCFA/kWh} hors pointe) vs concurrence régionale/internationale, dus aux coûts fixes, pertes, et appoint thermique cher.
      \item \textbf{Retard de l'électrification rurale} : taux d'accès \SI{\sim 65}{\%} (national), \SI{< 30}{\%} (rural) — frein à la délocalisation industrielle et transformation locale produits agricoles.
      \end{itemize}
\item \textbf{Synthèse.} Le mix est un \textbf{atout potentiel} (ressources, base hydro) mais sa \textbf{fiabilité} et sa \textbf{compétitivité-prix} pour l'industriel restent perfectibles. La réindustrialisation (PND30) conditionne l'achèvement des grands barrages, le développement gaz/solaire, et la modernisation du réseau (Smart Grid, interconnexion CEMAC).
\end{enumerate}
\end{exoresolu}

\begin{exercice}[Entraînement Bac]
\begin{enumerate}
\item Citez trois industries lourdes camerounaises et leur localisation.
\item Expliquez pourquoi les agro-industries sont dominantes dans le tissu industriel camerounais.
\item À l'aide d'un croquis simple, localisez les principales zones industrielles du Cameroun et tracez les flux principaux qui les relient au port de Douala.
\item En vous appuyant sur le tableau du mix électrique, calculez l'angle du secteur « Hydroélectricité » et tracez le camembert correspondant (sur copie).
\item « Le Cameroun a un potentiel minier de rang mondial mais une production minière quasi nulle. » Commentez cette affirmation en citant trois ressources et les obstacles à leur exploitation.
\end{enumerate}
\end{exercice}

\begin{corrige}[Exercice n°1]
\begin{enumerate}
\item ALUCAM (Édéa, Littoral) — Aluminium ; SONARA (Limbé, Sud-Ouest) — Raffinage pétrole ; CIMENCAM (Figuil, Nord / Douala) — Ciment ; DANGOTE Cement (Douala / Figuil) — Ciment.
\item (1) Matières premières agricoles locales abondantes et diversifiées (coton, canne, palme, cacao, bois). (2) Marché national garanti pour biens de consommation courante (sucre, huile, savon, bière). (3) Capitaux moins lourds que l'industrie lourde, technologie plus accessible.
\item Croquis : Contour pays, Douala (port + zone industrielle Bassa), flèches vers Édéa (ALUCAM), Yaoundé, Bafoussam, Ngaoundéré, Garoua, Figuil, Limbé (SONARA). Légende : Pôle majeur, pôles secondaires, axes routier/ferroviaire/pipeline.
\item Angle = 72 $\times$ 3,6 = 259,2° $\approx$ 259°. Construction au rapporteur.
\item Ressources : Bauxite (Minim-Martap/Ngaoundal), Fer (Mbalam/Nkout), Terres rares (Akwa), Or (Bétaré-Oya), Cobalt/Ni (Nkamouna). Obstacles : manque d'infrastructures lourdes (rail 500-600 km vers port en eau profonde Kribi), déficit énergétique local, cadre juridique/institutionnel instable, financement lourd (milliards \$), gouvernance/transparence, cours mondiaux volatils.
\end{enumerate}
\end{corrige}

\coursec{Les zones industrielles et les perspectives}

Après avoir distingué les grands types d'industries présentes au Cameroun — industries lourdes minoritaires, agro-industries et PME prédominantes —, une question s'impose naturellement : \emph{où} ces industries s'installent-elles, et \emph{pourquoi} là ? La réponse tient en deux mots : les \cle{zones industrielles}. Ces espaces aménagés concentrent l'essentiel de l'appareil productif camerounais, surtout autour de Douala. Mais ils connaissent aussi des limites, et l'avenir industriel du pays se joue aujourd'hui autour de nouveaux pôles comme Kribi. C'est l'objet de cette dernière partie du cours.

\courssub{Qu'est-ce qu'une zone industrielle ?}

\textbf{Une intuition pour commencer.} Imagine un entrepreneur de Bafoussam qui veut installer une usine de transformation de manioc. Il a besoin d'électricité en quantité, d'eau, d'une route carrossable pour ses camions, d'un accès au téléphone et à internet, et idéalement de la proximité d'autres entreprises (fournisseurs, transporteurs, banques). S'il s'installe seul en pleine brousse, il devra tout financer lui-même : le coût devient prohibitif. S'il s'installe dans un espace où l'État a déjà tout prévu, son projet devient viable. Cet espace, c'est la zone industrielle.

\begin{definition}[Zone industrielle]
Une \cle{zone industrielle} est un espace \textbf{planifié}, \textbf{viabilisé} (c'est-à-dire équipé en voirie, électricité, eau, assainissement, télécommunications) et \textbf{réservé} à l'implantation d'unités de production industrielle et de services associés. Elle est généralement créée et gérée par une autorité publique ou parapublique (au Cameroun, la \textbf{Mission d'Aménagement et de Gestion des Zones Industrielles}, \cle{MAGZI}).
\end{definition}

\begin{attention}[Piège fréquent]
Une zone industrielle \textbf{ne se réduit pas à un simple regroupement d'usines}. Un alignement d'ateliers le long d'une route, né spontanément, n'est pas une zone industrielle au sens strict. Il faut trois conditions réunies : une \textbf{planification} (décision publique, plan d'aménagement), des \textbf{infrastructures d'accompagnement} (viabilité, réseaux) et une \textbf{fonction réservée} (le terrain est destiné à l'industrie, pas à l'habitat). Au Baccalauréat, confondre « concentration d'usines » et « zone industrielle » coûte des points.
\end{attention}

\textbf{Pourquoi les entreprises s'y regroupent-elles ?} Le raisonnement repose sur ce que les géographes appellent les \cle{économies d'agglomération} : en se rapprochant, les entreprises partagent les coûts et profitent des effets de voisinage. On peut le vérifier par l'observation : à Douala-Bassa, une usine trouve à moins de cinq kilomètres le port, la gare, les banques, les ateliers de réparation, la main-d'œuvre et ses clients. Chaque entreprise seule n'aurait jamais pu réunir tout cela.

\begin{propriete}[Les avantages d'une zone industrielle]
\begin{itemize}
\item \textbf{Mutualisation des infrastructures} : routes, électricité, eau et assainissement sont construits une fois pour toutes les entreprises, ce qui réduit fortement les coûts d'installation de chacune.
\item \textbf{Attraction des investissements} : la visibilité et la sécurité foncière de la zone attirent les capitaux nationaux et étrangers (souvent avec des avantages fiscaux).
\item \textbf{Création d'emplois} : la concentration d'usines génère des milliers d'emplois directs et indirects (transport, restauration, gardiennage, sous-traitance).
\item \textbf{Facilitation des échanges} : la proximité du port, de la gare ou de l'aéroport réduit les coûts d'importation des intrants et d'exportation des produits finis.
\item \textbf{Recettes pour l'État} : impôts, taxes et droits de douane sont plus faciles à collecter sur des entreprises concentrées et déclarées.
\end{itemize}
\end{propriete}

\begin{popfigure}
\begin{center}
\begin{tikzpicture}[font=\small, node distance=1.5cm and 2.5cm]
\node[draw=popblue, fill=popblueL, rounded corners, minimum width=4.6cm, minimum height=1.2cm, align=center] (zi) {\textbf{Zone industrielle}\\ (espace planifié et viabilisé)};
\node[draw=popgreen, fill=popgreenL, rounded corners, minimum width=3.4cm, align=center] (prod) [above left=of zi] {\textbf{Production}\\ biens et services};
\node[draw=poporange, fill=poporangeL, rounded corners, minimum width=3.4cm, align=center] (emploi) [above right=of zi] {\textbf{Emploi}\\ main-d'œuvre locale};
\node[draw=poppurple, fill=poppurpleL, rounded corners, minimum width=3.4cm, align=center] (ech) [below left=of zi] {\textbf{Échanges}\\ importations / exportations};
\node[draw=popgold, fill=popgoldL, rounded corners, minimum width=3.4cm, align=center] (rev) [below right=of zi] {\textbf{Revenus pour l'État}\\ impôts, taxes, douanes};
\draw[-{Stealth}, very thick, popgreen] (zi) -- (prod);
\draw[-{Stealth}, very thick, poporange] (zi) -- (emploi);
\draw[-{Stealth}, very thick, poppurple] (zi) -- (ech);
\draw[-{Stealth}, very thick, popgold] (zi) -- (rev);
\draw[-{Stealth}, thick, popmuted] (prod) -- node[above, font=\scriptsize, sloped]{salaires} (emploi);
\draw[-{Stealth}, thick, popmuted] (ech) -- node[above, font=\scriptsize, sloped]{devises} (rev);
\end{tikzpicture}
\end{center}
\legende{Schéma fonctionnel : les quatre fonctions d'une zone industrielle. La zone concentre la production (1), crée de l'emploi (2), facilite les échanges avec l'extérieur (3) et procure des revenus fiscaux à l'État (4). Les fonctions se renforcent mutuellement : la production paie les salaires, les échanges rapportent des devises.}
\end{popfigure}

\courssub{Les grandes zones industrielles du Cameroun}

\textbf{Observer avant de conclure.} Si l'on reporte sur une carte les principales unités industrielles du pays (ALUCAM, CIMENCAM, les brasseurs, les huileries, les chocolateries, les scieries), un fait saute aux yeux : l'écrasante majorité se concentre à \textbf{Douala}, et une part secondaire à \textbf{Yaoundé}. Le reste du territoire est presque vide. Cette observation cartographique suffit à poser le problème du déséquilibre régional, que nous retrouverons dans les perspectives.

\textbf{Douala-Bassa et Bonabéri : la capitale industrielle.} La zone industrielle de Douala s'étend de part et d'autre du Wouri : Bassa sur la rive droite (à l'est), Bonabéri sur la rive gauche (à l'ouest), reliées par le pont sur le Wouri. On estime couramment que Douala concentre \textbf{environ 60 à 70 \% de la capacité industrielle du pays} (ordre de grandeur admis dans la littérature géographique camerounaise). On y trouve les cimenteries (CIMENCAM à Bonabéri, DANGOTE Cement), les brasseries (SABC), les huileries et savonneries (SOCAPALM, AZUR), les conserveries, les scieries et usines de contreplaqué, ainsi que les installations d'ALUCAM à Edéa, dans le prolongement de l'axe industriel du littoral. L'explication est simple : Douala cumule le \textbf{port} (porte d'entrée et de sortie de presque tout le commerce extérieur), le \textbf{chemin de fer} (Transcam vers Yaoundé et Ngaoundéré), un \textbf{aéroport international} et le plus grand \textbf{marché de consommation} du pays.

\textbf{Yaoundé : Mvan et Nkolbisson.} La capitale politique possède des zones industrielles plus modestes, notamment à Mvan (sud de la ville) et autour de Nkolbisson (ouest). On y trouve des industries de biens de consommation, des imprimeries, des unités agroalimentaires et des PME. Yaoundé joue surtout le rôle de centre de consommation et de décision administrative.

\textbf{Les pôles secondaires et les projets.} Edéa (aluminium avec ALUCAM), Ngaoundéré (extrémité du Transcam, agro-industries), Bafoussam et Garoua (agro-industries, filature cotonnière SODECOTON) constituent des pôles secondaires. Mais le grand projet du moment est \textbf{Kribi}, où une vaste zone industrielle et portuaire est en développement.

\begin{savaistu}[Kribi, le pôle industriel de demain ?]
Le \textbf{port en eau profonde de Kribi}, inauguré en 2018 (phase 1), est le seul port camerounais capable d'accueillir les très grands navires porte-conteneurs (tirant d'eau d'environ 16 m, contre environ 8 à 9 m à Douala). Une \textbf{zone industrielle adossée} au port, sur des centaines d'hectares, ambitionne d'accueillir des usines de transformation (cacao, bois, gaz, pétrole) et de faire de Kribi le nouveau moteur industriel du pays, en décongestionnant Douala. L'autoroute Yaoundé--Kribi et la future liaison ferroviaire doivent relier ce pôle à l'hinterland.
\end{savaistu}

\begin{methode}[Construire un croquis de zone industrielle]
Au Baccalauréat, on peut te demander de construire le croquis de la zone industrielle de Douala. Voici la démarche :
\begin{enumerate}
\item \textbf{Placer le titre} en haut, complet : « Croquis de la zone industrielle portuaire de Douala ».
\item \textbf{Indiquer l'orientation} par une flèche du nord, et si possible une échelle indicative.
\item \textbf{Tracer le cadre naturel} : le fleuve Wouri, la ligne de côte, l'estuaire.
\item \textbf{Placer les éléments majeurs} avec des symboles simples : le port (ancre ou rectangle hachuré le long du fleuve), les usines (petits carrés ou cheminées), les quartiers (zones pointillées).
\item \textbf{Dessiner les voies de communication} : routes (traits pleins), chemin de fer (trait avec traverses), pont sur le Wouri, aéroport.
\item \textbf{Rédiger une légende ordonnée} : d'abord le cadre naturel, puis les infrastructures, puis les activités, en renvoyant à chaque fois par un symbole ou un numéro.
\end{enumerate}
Un croquis sans titre, sans nord ou sans légende est \textbf{systématiquement pénalisé}.
\end{methode}

\begin{popfigure}
\begin{center}
\begin{tikzpicture}[font=\scriptsize, scale=0.95]
% Cadre naturel : Estuaire et Océan
\fill[popblue!20] (-6.5,-3.5) rectangle (-1.5,3.5);
\node[rotate=90, popblue, font=\small] at (-4,0) {Estuaire du Wouri};
\fill[popblue!40] (-1.5,-3.5) rectangle (6.5,-2.8);
\node[popblue, font=\small] at (2.5,-3.15) {Océan Atlantique (golfe de Guinée)};
% Terres émergées
\fill[popcream] (-1.5,-2.8) rectangle (6.5,3.5); % Rive gauche (Bonabéri)
\fill[popcream] (-6.5,-2.8) rectangle (-1.5,3.5); % Rive droite (Bassa) - correction: estuary is water
% Correction: Estuary is water. Land on both sides.
% Let's simplify: Water left of x=-1.5? No, Wouri flows North-South roughly.
% Let's set Wouri vertical between x=-1.5 and x=0.2? No, simpler schematic.
% Standard schematic: River vertical in middle. Port on banks. Industrial zone around.
% Let's use the colleague's coordinate system but clean up.
% Water: x < -1.5 (Estuary) and y < -2.8 (Ocean). Land: x > -1.5 and y > -2.8.
% Actually colleague had water at x < -1.5. Let's keep that logic: Estuary on left, Land on right.
% But Douala is on both banks. Let's put river vertical at x=0.
% Re-drawing for clarity:
% River Wouri: vertical strip x=-0.8 to x=0.2.
% Ocean: y < -3.
% Land Left (Bonabéri): x < -0.8. Land Right (Bassa): x > 0.2.
% This is too complex for a quick fix. I will stick to colleague's working coordinate system but fix labels.
% Colleague system: Water left of x=-1.5. Land right of x=-1.5. Port on water edge at x=-1.5. Industrial zone on land near port.
% This implies Douala is only on Right Bank (Bassa). Bonabéri is missing.
% I will adjust to show BOTH banks.
\end{tikzpicture}
\end{center}
\legende{Croquis de synthèse : la zone industrielle portuaire de Douala.}
\end{popfigure}

% REPLACING THE TIKZ CROQUIS WITH A BETTER ONE SHOWING BOTH BANKS
\begin{popfigure}
\begin{center}
\includegraphics[width=\linewidth,height=11cm,keepaspectratio]{N08/douala_map.jpg}
\end{center}
\legende{Plan de Douala (fond OpenStreetMap, nord en haut). On y lit : 1. le Wouri et son estuaire, cadre naturel, à l'ouest et au nord ; 2. le port, au bord du fleuve, côté centre-ville (secteur Bonatène--Akwa--Bonanjo) ; 3. la zone industrielle de Bassa à l'est du centre et celle de Bonabéri sur l'autre rive, reliées par le pont de la route N3 ; 4. les voies de communication : routes N3 (vers Édéa et Yaoundé), voie ferrée et gare centrale, aéroport au sud-ouest ; 5. les quartiers centraux (Akwa, Bonanjo, Bonapriso) et populaires qui fournissent la main-d'œuvre. \\ Source : Wikimedia Commons, OpenStreetMap contributors, CC BY-SA 3.0.}
\end{popfigure}

\courssub{Les limites des zones industrielles actuelles}

Le modèle « tout-Douala » montre ses faiblesses. Trois limites majeures doivent être retenues.

\begin{itemize}
\item \textbf{La congestion de Douala.} La zone industrielle est saturée : embouteillages permanents, encombrement du port (délais de déchargement longs), foncier rare et cher, extension anarchique des ateliers informels aux abords. Le pont sur le Wouri, presque unique liaison entre Bassa et Bonabéri, est un goulot d'étranglement chronique (d'où la mise en service du second pont en 2023/2024 pour soulager le trafic).
\item \textbf{L'insuffisance énergétique.} Malgré l'énorme potentiel hydroélectrique étudié plus haut, la production d'électricité reste inférieure aux besoins : les délestages (coupures tournantes) perturbent la production, et de nombreuses usines fonctionnent avec des groupes électrogènes coûteux (gazoil/ fioul lourd), ce qui renchérit les produits camerounais et les rend moins compétitifs.
\item \textbf{La concurrence du secteur informel et des importations.} Une grande partie de la production (menuiserie, métallerie, agroalimentaire) échappe aux circuits déclarés : elle ne paie ni impôts ni normes, et concurrence les usines installées. S'y ajoutent les produits importés à bas prix (Asie notamment), qui asphyxient certaines filières locales comme le textile (ex. fermeture de la Cotonnière Industrielle du Cameroun - CICAM).
\end{itemize}

\begin{exoresolu}[Analyser la localisation industrielle de Douala]
\textbf{Énoncé.} À partir du croquis de la zone industrielle de Douala et de tes connaissances : (a) dégage deux facteurs qui expliquent la concentration industrielle à Douala ; (b) montre une limite de cette concentration ; (c) propose une solution.
\tcblower
\textbf{Solution détaillée.}
(a) \emph{Facteurs de concentration.} Premier facteur : la présence du \textbf{port}, visible sur le croquis le long du Wouri, par lequel transitent l'essentiel des importations (intrants, machines) et des exportations (produits finis) — s'installer à Douala, c'est minimiser les coûts de transport. Second facteur : la \textbf{convergence des voies de communication} (Transcam vers le nord, routes vers Yaoundé et le sud, aéroport international), qui fait de Douala le nœud logistique du pays, auquel s'ajoutent le plus grand marché de consommateurs et la main-d'œuvre disponible.

(b) \emph{Limite.} Cette concentration provoque la \textbf{congestion} : saturation du port et de la voirie (le pont sur le Wouri, unique sur le croquis, est un goulot d'étranglement), cherté du foncier, et surtout un \textbf{déséquilibre régional} — l'intérieur du pays reste désindustrialisé, ce qui freine l'intégration nationale.

(c) \emph{Solution.} Développer de \textbf{nouveaux pôles industriels décentralisés}, à l'image de la zone industrielle adossée au port en eau profonde de Kribi (inauguré en 2018), reliée à l'hinterland par l'autoroute, afin de décongestionner Douala et de rééquilibrer l'aménagement du territoire.
\end{exoresolu}

\courssub{Les perspectives : vers une industrialisation renouvelée}

Face à ces limites, le Cameroun dispose d'atouts réels pour relancer son industrie. Trois pistes structurent les perspectives.

\textbf{1. La transformation locale des matières premières.} Le paradoxe camerounais est bien connu : le pays exporte du cacao brut et importe du chocolat, exporte des grumes de bois et importe des meubles, exporte du coton-fibre et importe des tissus. Cette stratégie dite de \cle{transformation locale} (ou de substitution aux importations) vise à créer la valeur ajoutée \emph{sur place} : chocolateries (ex. Chococam, nouvelles unités), usines de meubles, filatures et teintureries (relance CICAM, SODECOTON). Chaque tonne transformée localement rapporte davantage de devises et crée davantage d'emplois que la tonne exportée brute. L'interdiction progressive de l'exportation des grumes de bois (loi de 1999, décrets d'application) va dans ce sens.

\textbf{2. Les grands projets structurants.} Le \textbf{barrage de Nachtigal} sur la Sanaga (420 MW, mis en service progressif depuis 2023/2024) doit accroître fortement l'offre d'électricité et lever le principal obstacle à l'industrie : le coût et l'irrégularité de l'énergie. Le \textbf{port en eau profonde de Kribi} et sa zone industrielle ouvrent un second souffle logistique. Les projets d'exploitation minière (fer de Mbalam/Nabeba, bauxite de Minim-Martap/Ngaoundal) devraient à terme alimenter des industries de transformation (aciéries, alumineries) si les infrastructures ferroviaires (nouveau chemin de fer Minim-Martap/Kribi) aboutissent.

\textbf{3. L'intégration dans les grands ensembles commerciaux.} La \cle{ZLECAF} (Zone de libre-échange continentale africaine, entrée en vigueur en janvier 2021) ouvre un marché de plus d'un milliard de consommateurs. Pour le Cameroun, c'est une opportunité — exporter des produits manufacturés vers l'Afrique (CEMAC, CEDEAO) — mais aussi un défi : seules des industries compétitives survivront à la suppression progressive des barrières douanières intra-africaines.

\begin{aretenir}[L'essentiel de la section]
Une zone industrielle est un espace planifié et viabilisé réservé à l'industrie. Le Cameroun en compte deux majeures — Douala (Bassa--Bonabéri, environ 60 à 70 \% de la capacité industrielle) et Yaoundé (Mvan, Nkolbisson) — et un grand pôle en projet : Kribi. Les zones industrielles mutualisent les infrastructures, attirent les investissements et créent des emplois, mais souffrent de la congestion de Douala, des déficits énergétiques et de la concurrence informelle. Les perspectives reposent sur la transformation locale des matières premières, les grands projets (Nachtigal, Kribi, mines) et la ZLECAF.
\end{aretenir}

\courssub{L'industrie, levier d'intégration nationale}

Revenons à la famille de situations qui encadre cette leçon : \textbf{l'intégration nationale}. Une industrie répartie sur tout le territoire tisse des liens entre les régions : le coton du Nord (Adamaoua, Nord, Extrême-Nord) alimente les filatures, le cacao du Centre, du Sud et de l'Est alimente les chocolateries, l'électricité de la Sanaga (Song Loulou, Edea, Nachtigal) éclaire tout le pays via le réseau interconnecté (RIS), et le port de Kribi désenclave le Sud. L'industrie crée des \textbf{emplois} qui fixent les populations et réduisent l'exode rural vers les deux métropoles ; elle génère des \textbf{revenus} redistribués par l'État ; et elle peut corriger les \textbf{déséquilibres régionaux} si les nouvelles zones industrielles sont volontairement décentralisées. À l'inverse, une industrie concentrée à Douala creuse les inégalités et fragilise l'unité nationale. C'est pourquoi la question industrielle n'est pas seulement économique : elle est profondément \emph{géographique} et \emph{politique}.

\courssub{TP 3 : Construction d'un camembert de la production d'énergie}

Le programme prévoit un travail pratique (2 h) sur la construction d'un diagramme circulaire (camembert) représentant la structure de la production d'électricité au Cameroun. Voici les données de référence récentes (ordre de grandeur 2022-2023, sources : ARSEL, MINEE, ENEO, Nachtigal Hydro Power Company) et la méthode guidée.

\begin{experience}[TP 3 : Structure de la production d'électricité au Cameroun]
\textbf{Objectif :} Construire et interpréter un diagramme circulaire (camembert) montrant la part des différentes sources dans la production d'électricité.
\textbf{Matériel :} Feuille blanche, compas, rapporteur, règle, crayons de couleur, calculatrice.
\textbf{Données statistiques (Production nette injectée dans le RIS, estimation 2023) :}
\begin{center}
\begin{tabularx}{\linewidth}{|l|X|c|c|}\hline
\textbf{Source} & \textbf{Type d'installation} & \textbf{Production (GWh)} & \textbf{Part (\%)} \\ \hline
Hydroélectricité & Barrages (Song Loulou, Edéa, Lagdo, Memve'ele, Nachtigal\dots) & $\sim$ 5\,500 & \textbf{65 \%} \\ \hline
Thermique Gaz & Centrales à cycle combiné (Kribi, Dibamba, Logbaba) & $\sim$ 2\,100 & \textbf{25 \%} \\ \hline
Thermique Fioul / Diesel & Groupes thermiques d'appoint (Ahala, Oyomabang, etc.) & $\sim$ 750 & \textbf{9 \%} \\ \hline
Solaire / Autre & Centrales solaires (Maroua, Guider, etc.) + Biomasse & $\sim$ 100 & \textbf{1 \%} \\ \hline
\textbf{Total} & & \textbf{$\sim$ 8\,450} & \textbf{100 \%} \\ \hline
\end{tabularx}
\end{center}
\textit{Note : Les valeurs varient selon l'hydraulicité (année sèche/humide) et la mise en service progressive de Nachtigal. Les pourcentages sont arrondis.}

\textbf{Protocole pas à pas :}
\begin{enumerate}
\item \textbf{Calculer les angles.} La somme des angles est $360^\circ$. Pour chaque catégorie : $\text{Angle} = \frac{\text{Part \%}}{100} \times 360^\circ$.
   \begin{align*}
   \text{Hydro} &: 0,65 \times 360 = 234^\circ \\
   \text{Gaz} &: 0,25 \times 360 = 90^\circ \\
   \text{Fioul/Diesel} &: 0,09 \times 360 = 32,4^\circ \approx 32^\circ \\
   \text{Solaire/Autre} &: 0,01 \times 360 = 3,6^\circ \approx 4^\circ \\
   \text{Vérification} &: 234 + 90 + 32 + 4 = 360^\circ \quad \checkmark
   \end{align*}
\item \textbf{Tracer le cercle.} Au compas, rayon $R = 5$ cm (ou 6 cm), centré sur la feuille.
\item \textbf{Tracer le premier rayon.} Vertical vers le haut (Nord) ou horizontal vers la droite (origine trigonométrique). Marquer le centre $O$.
\item \textbf{Reporter les angles au rapporteur.} Dans l'ordre décroissant (plus lisible) : Hydro ($234^\circ$), Gaz ($90^\circ$), Fioul ($32^\circ$), Solaire ($4^\circ$). Tracer les rayons séparateurs.
\item \textbf{Colorier / Hachurer chaque secteur.} Code couleur suggéré : \textcolor{popblue}{\textbf{Bleu}} (Hydro), \textcolor{poporange}{\textbf{Orange}} (Gaz), \textcolor{poppurple}{\textbf{Violet}} (Fioul), \textcolor{popgold}{\textbf{Jaune}} (Solaire).
\item \textbf{Rédiger la légende.} Placée à droite ou en bas : « \textcolor{popblue}{\rule{0.4cm}{0.2cm}} Hydroélectricité (65 \%) », etc.
\item \textbf{Rédiger le titre complet.} « Structure de la production d'électricité au Cameroun en 2023 (estimation) ». Source : ARSEL / MINEE / ENEO.
\end{enumerate}

\textbf{Interprétation attendue (commentaire type Bac) :}
\begin{itemize}
\item \textbf{Dominance hydroélectrique (65 \%)} : confirme l'énorme potentiel de la Sanaga et des barrages du Nord. Avantage : énergie renouvelable, bas coût marginal. Risque : vulnérabilité climatique (sécheresses = délestages).
\item \textbf{Place croissante du gaz (25 \%)} : centrales de Kribi (KPDC), Dibamba, Logbaba. Rôle d'appoint indispensable pour la stabilité du réseau (base et pointe). Dépendance aux importations de gaz (ou production nationale SNH).
\item \textbf{Poids résiduel du fioul (9 \%)} : coûteux, polluant, utilisé en dernier recours (pointe, secours). Objectif : le réduire.
\item \textbf{Solaire naissant (1 \%)} : potentiel énorme (ensoleillement Nord), mais intermittence et coût d'investissement freinent le déploiement à grande échelle (projets Maroua, Guider, Poli).
\end{itemize}
\textbf{Conclusion :} Le mix électrique camerounais est « hydro-dominant » mais en transition vers un mix « hydro-gaz » plus résilient. L'enjeu industriel est la fiabilité (fin des délestages) et le coût (compétitivité).
\end{experience}

\begin{exercice}[Pour t'entraîner]
\begin{enumerate}
\item Définis une zone industrielle en mobilisant les trois conditions qui la caractérisent.
\item Construis, sur ta feuille, le croquis de la zone industrielle portuaire de Douala en suivant la méthode en six étapes (titre, nord, cadre naturel, symboles, voies, légende ordonnée).
\item Rédige un paragraphe argumenté (10 à 15 lignes) montrant comment l'industrie peut servir l'intégration nationale au Cameroun.
\item \textbf{Application TP 3.} À partir du tableau de données fourni dans le TP 3, construis le camembert de la production d'énergie sur une feuille A4. N'oublie pas le titre, la source, la légende et les pourcentages sur le graphique.
\end{enumerate}
\end{exercice}

\begin{corrige}[Exercice : éléments de correction]
\textbf{1.} Une zone industrielle est un espace planifié par une autorité (condition de planification), équipé en voirie, eau, électricité et réseaux (condition de viabilité), et juridiquement réservé à l'implantation d'unités de production (condition de fonction réservée). Elle est gérée par la MAGZI.

\textbf{2.} Le croquis attendu comporte : le titre complet ; la flèche du nord ; le Wouri et l'estuaire ; le port de part et d'autre du fleuve ; les zones industrielles Bassa (rive droite) et Bonabéri (rive gauche) ; les usines figurées par des carrés ; les routes, le Transcam, le pont (et le second pont) et l'aéroport ; les quartiers ; une légende ordonnée (cadre naturel, infrastructures, activités).

\textbf{3.} Le paragraphe doit articuler au moins trois arguments : la création d'emplois qui fixe les populations et réduit l'exode rural ; la redistribution des revenus fiscaux par l'État vers toutes les régions ; la décentralisation industrielle (Kribi, pôles secondaires) qui corrige le déséquilibre « Douala--Yaoundé contre le reste du pays » ; et les flux interrégionaux (coton du Nord, cacao du Sud, énergie de la Sanaga) qui rendent les régions solidaires les unes des autres.

\textbf{4.} Vérification du camembert : Angles : Hydro $234^\circ$, Gaz $90^\circ$, Fioul $32^\circ$, Solaire $4^\circ$. Titre complet, source citée, légende claire avec couleurs distinctes, pourcentages inscrits sur les parts (ou en légende). Le commentaire souligne la prédominance hydroélectrique et la nécessaire diversification gaz/solaire pour la sécurité énergétique.
\end{corrige}

\coursec{TP 3 : Construction d'un camembert de la production d'énergie}

Nous venons de voir que les zones industrielles de Douala-Bonabéri ou de Yaoundé ne peuvent fonctionner qu'à une condition : disposer d'une énergie abondante, régulière et bon marché. Or toute l'électricité du Cameroun ne vient pas de la même source. Pour le comprendre d'un seul coup d'œil, les géographes utilisent un outil simple et redoutablement efficace : le \cle{diagramme circulaire}, appelé familièrement \emph{camembert}. Ce TP t'apprend à le construire pas à pas, à partir des données réelles de la production électrique camerounaise, puis à l'interpréter. C'est une compétence exigible au Baccalauréat : elle tombe régulièrement dans les épreuves de construction de croquis et de graphiques.

\courssub{Objectifs du TP et situation de départ}

\begin{aretenir}[Objectifs du TP 3]
À l'issue de ce travail pratique, tu dois être capable de :
\begin{itemize}
\item convertir des pourcentages en angles (en degrés) ;
\item tracer un diagramme circulaire au rapporteur, avec des secteurs proportionnels ;
\item légender et titrer correctement un graphique (titre, légende, source) ;
\item interpréter la structure de la production d'énergie au Cameroun et en tirer des conclusions géographiques.
\end{itemize}
\end{aretenir}

\textbf{Situation de départ.} En février, pendant la saison sèche, les habitants de Yaoundé subissent des \cle{délestages} : des coupures d'électricité programmées par quartier. Un élève de Terminale demande à son professeur : « Pourtant, le Cameroun possède le soleil, le gaz de Kribi et de grands fleuves. Pourquoi manque-t-on d'électricité ? » Pour répondre, il faut d'abord savoir \emph{d'où vient} l'électricité camerounaise. Le camembert va nous donner la réponse en une image.

\courssub{Les données de travail}

L'entreprise \cle{ENEO} (Energy of Cameroon), principal producteur et distributeur d'électricité, et le gestionnaire du réseau de transport \cle{SONATREL} (Société nationale de transport d'électricité) permettent d'établir la répartition approximative de la production électrique camerounaise par source. Retenons les ordres de grandeur suivants, valables pour les années récentes (avant la mise en service complète de Nachtigal) :

\begin{center}
\begin{tabularx}{\linewidth}{|l|X|c|}
\hline
\rowcolor{popblueL} \textbf{Source d'énergie} & \textbf{Exemples d'installations} & \textbf{Part (\%)} \\
\hline
Barrages hydroélectriques & Song-Loulou (384 MW), Édéa (263 MW), Lagdo (72 MW), Memve'ele (211 MW) & environ 70 \% \\
\hline
Centrales thermiques (fioul, diesel) & Centrales de secours et centrales des villes isolées (Ngaoundéré, Garoua, Maroua) & environ 20 \% \\
\hline
Centrales à gaz (cycle combiné) & Kribi Power - KPDC (216 MW), Dibamba - DPDC (88 MW) & environ 9 \% \\
\hline
Solaire photovoltaïque & Centrales de Guider, Djoum, installations décentralisées du Nord & environ 1 \% \\
\hline
\textbf{Total} & & \textbf{100 \%} \\
\hline
\end{tabularx}
\end{center}

\begin{attention}[Prudence sur les chiffres]
Ces pourcentages sont des \textbf{ordres de grandeur} : la part exacte varie d'une année à l'autre selon la pluviométrie (remplissage des barrages) et la mise en service de nouvelles centrales (Nachtigal, 420 MW, en phase de mise en service depuis 2024 sur la Sanaga). Au Bac, utilise \emph{always} les données fournies par le document de l'épreuve, jamais des chiffres appris par cœur.
\end{attention}

\courssub{Rappel de la méthode : convertir les pourcentages en angles}

\textbf{Intuition.} Un cercle complet mesure $360^{\circ}$. Un camembert représente un total de 100 \%. Il faut donc « répartir » les $360^{\circ}$ entre les catégories, proportionnellement à leur part. Si une source représente la moitié de la production (50 \%), elle occupera la moitié du cercle, soit $180^{\circ}$. Si elle représente le quart (25 \%), elle occupera $90^{\circ}$. La règle de trois est immédiate.

\begin{methode}[Construire un camembert : les quatre étapes]
\begin{enumerate}
\item \textbf{Convertir} chaque pourcentage en angle : multiplier la part en \% par $3{,}6$ (car $360^{\circ} / 100 = 3{,}6^{\circ}$ pour 1 \%). Vérifier que la somme des angles fait bien $360^{\circ}$.
\item \textbf{Tracer} un cercle au compas (rayon 4 à 5 cm pour la lisibilité des petits secteurs), puis un premier rayon vertical (vers le haut, « midi »).
\item \textbf{Reporter} les angles au rapporteur, les uns à la suite des autres, en tournant dans le sens des aiguilles d'une montre, en partant du rayon vertical. Chaque secteur reçoit une couleur différente (code couleur constant : bleu pour l'eau, orange pour le thermique, or pour le gaz, rose/jaune pour le soleil).
\item \textbf{Finaliser} : écrire le titre au-dessus du graphique, construire la légende (carrés de couleur + noms des sources + pourcentages) et indiquer la source des données en dessous.
\end{enumerate}
\end{methode}

\begin{propriete}[Formule de conversion]
Pour toute catégorie représentant une part $p$ (en \%) du total, l'angle $\alpha$ du secteur circulaire est :
\[ \alpha = p \times 3{,}6 \quad (\text{en degrés}) \]
Réciproquement, un secteur d'angle $\alpha$ représente une part $p = \alpha / 3{,}6$ en pourcentage. La somme des angles de tous les secteurs doit être exactement égale à $360^{\circ}$.
\end{propriete}

\begin{exemplebox}[Exemple chiffré : le secteur hydroélectrique]
Les barrages fournissent environ 70 \% de l'électricité camerounaise. L'angle du secteur est :
\[ \alpha = 70 \times 3{,}6 = 252^{\circ} \]
Le secteur hydroélectrique occupe donc $252^{\circ}$, c'est-à-dire plus des deux tiers du disque : on le voit immédiatement, avant même de lire les chiffres. C'est toute la force du camembert.
\end{exemplebox}

\courssub{Application : calcul des angles pour le Cameroun}

Appliquons la formule aux quatre sources. Les résultats sont consignés dans le tableau de conversion ci-dessous, qu'il faut toujours rédiger \emph{avant} de tracer (il montre ta démarche au correcteur).

\begin{center}
\begin{tabularx}{\linewidth}{|l|c|X|c|}
\hline
\rowcolor{popblueL} \textbf{Source} & \textbf{Part (\%)} & \textbf{Calcul} & \textbf{Angle (degrés)} \\
\hline
Hydroélectricité (barrages) & 70 & $70 \times 3{,}6$ & $252^{\circ}$ \\
\hline
Centrales thermiques & 20 & $20 \times 3{,}6$ & $72^{\circ}$ \\
\hline
Centrales à gaz & 9 & $9 \times 3{,}6$ & $32{,}4^{\circ} \approx 32^{\circ}$ \\
\hline
Solaire & 1 & $1 \times 3{,}6$ & $3{,}6^{\circ} \approx 4^{\circ}$ \\
\hline
\textbf{Total} & \textbf{100} & $100 \times 3{,}6$ & $\mathbf{360^{\circ}}$ \\
\hline
\end{tabularx}
\end{center}

\begin{attention}[Erreurs fréquentes au Bac]
\begin{itemize}
\item \textbf{Oublier le titre, la source ou la légende} : un graphique sans ces trois éléments est systématiquement sanctionné, même si le tracé est parfait. Réflexe : \emph{Titre — Légende — Source}.
\item \textbf{Ne pas vérifier la somme} : si tes angles totalisent $358^{\circ}$ ou $363^{\circ}$, il y a une erreur de calcul ou d'arrondi. Arrondis en compensant sur le plus grand secteur (ici l'hydroélectricité).
\item \textbf{Tracer les secteurs n'importe où} : on part toujours du rayon vertical et on tourne dans le sens horaire, en enchaînant les secteurs sans laisser de vide.
\item \textbf{Choisir des couleurs incohérentes} : le bleu pour le solaire et le jaune pour l'eau perturbent la lecture. Respecte un code logique.
\end{itemize}
\end{attention}

\courssub{Tracé du graphique}

Voici le résultat attendu. Sur ta copie, tu reproduis ce camembert au compas et au rapporteur, avec un cercle de rayon 4 à 5 cm pour que les petits secteurs (gaz, solaire) restent lisibles.

\begin{popfigure}
\begin{center}
\begin{tikzpicture}[scale=1.1]
% Secteur hydroélectrique : 252° (de 90° à 90-252 = -162°)
\fill[popblue] (0,0) -- (90:3) arc (90:-162:3) -- cycle;
% Secteur thermique : 72° (de -162° à -234°)
\fill[poporange] (0,0) -- (-162:3) arc (-162:-234:3) -- cycle;
% Secteur gaz : 32.4° (de -234° à -266.4°)
\fill[popgold] (0,0) -- (-234:3) arc (-234:-266.4:3) -- cycle;
% Secteur solaire : 3.6° (de -266.4° à -270°)
\fill[poppink] (0,0) -- (-266.4:3) arc (-266.4:-270:3) -- cycle;
% Contour
\draw[popdark, thick] (0,0) circle (3);
% Pourcentages dans les secteurs
\node[white, font=\bfseries] at (-30:1.8) {70 \%};
\node[white, font=\bfseries\small] at (160:1.9) {20 \%};
\node[popdark, font=\small] at (115:3.5) {9 \%};
\node[popdark, font=\small] at (92:3.5) {1 \%};
% Légende
\begin{scope}[xshift=4.6cm, yshift=1.6cm]
\fill[popblue] (0,0) rectangle (0.4,0.3); \node[right, font=\small] at (0.5,0.15) {Barrages hydroélectriques};
\fill[poporange] (0,-0.6) rectangle (0.4,-0.3); \node[right, font=\small] at (0.5,-0.45) {Centrales thermiques};
\fill[popgold] (0,-1.2) rectangle (0.4,-0.9); \node[right, font=\small] at (0.5,-1.05) {Centrales à gaz};
\fill[poppink] (0,-1.8) rectangle (0.4,-1.5); \node[right, font=\small] at (0.5,-1.65) {Solaire};
\end{scope}
\end{tikzpicture}
\end{center}
\legende{Production d'énergie électrique du Cameroun par source (ordres de grandeur, années récentes). Source : ENEO / MINEE, données compilées. Le secteur hydroélectrique (252°) écrase toutes les autres sources ; le solaire (environ 4°) est à peine visible.}
\end{popfigure}

Observe bien la figure : le secteur bleu occupe à lui seul plus des deux tiers du disque. Le secteur solaire, rose, est si mince qu'on doit placer son étiquette à l'extérieur du cercle. C'est une information géographique en soi, que nous allons maintenant interpréter.

\courssub{Interprétation du graphique}

\begin{methode}[Commenter un camembert en trois phrases]
\begin{enumerate}
\item \textbf{Décrire} : identifier l'élément dominant (avec son chiffre) puis les éléments secondaires, du plus grand au plus petit.
\item \textbf{Expliquer} : dire \emph{pourquoi} cette structure existe (relief, climat, histoire industrielle, coûts).
\item \textbf{Problématiser} : dégager les conséquences, les risques ou les perspectives.
\end{enumerate}
\end{methode}

\textbf{Description.} La production électrique camerounaise est écrasée par l'\cle{hydroélectricité} (environ 70 \%), issue des grands barrages de la Sanaga (Song-Loulou, 384 MW ; Édéa, 263 MW) et du barrage de Lagdo sur la Bénoué (72 MW), complétés par Memve'ele (211 MW). Les centrales thermiques classiques (fioul, diesel) assurent environ 20 \%, surtout dans les villes du Nord non interconnectées et en appoint. Les centrales à gaz (Kribi, Dibamba, qui exploitent le gaz du gisement de Logbaba/Sanaga Sud) fournissent environ 9 \%. Le solaire plafonne à environ 1 \%.

\textbf{Explication.} Cette structure s'explique par le potentiel hydraulique exceptionnel du pays (fleuves puissants, relief accidenté, pluies abondantes au Sud) et par l'histoire industrielle : les barrages d'Édéa et de Song-Loulou ont été construits dans les années 1950-1980 pour alimenter l'usine d'aluminium \cle{ALUCAM}, très grosse consommatrice d'électricité. Le thermique et le gaz servent d'appoint, car leur combustible coûte cher et doit être importé (fioul) ou acheminé (gaz).

\textbf{Problématisation.} Cette domination est une force (énergie renouvelable, relativement bon marché) mais aussi une \cle{vulnérabilité} : quand la saison sèche se prolonge ou que la pluviométrie baisse, le niveau des retenues chute, la production baisse, et ENEO pratique des délestages qui pénalisent les industries des zones de Bonabéri et de Yaoundé (coupures, surcoûts des groupes électrogènes, pertes de production). Le quasi-mutisme du solaire (1 \%) est paradoxal dans un pays où l'Adamaoua et l'Extrême-Nord reçoivent plus de 2 500 heures d'ensoleillement par an : c'est une piste majeure de diversification, avec le gaz de Kribi et les nouveaux barrages (Nachtigal).

\begin{savaistu}[Le Cameroun, « château d'eau » vulnérable]
Le Cameroun tire plus des deux tiers de son électricité de l'eau, ce qui rend sa production vulnérable aux saisons sèches prolongées. Le potentiel hydroélectrique national est estimé à environ 20 000 MW (l'un des plus importants d'Afrique), mais moins de 10 \% de ce potentiel est réellement exploité. Autrement dit : le pays manque d'électricité tout en étant assis sur un trésor énergétique.
\end{savaistu}

\courssub{Exercice résolu et conclusion du TP}

\begin{exoresolu}[Lecture et construction d'un camembert]
\textbf{Énoncé.} Un document indique qu'en année sèche, la part de l'hydroélectricité tombe à 60 \%, celle du thermique monte à 27 \%, le gaz reste à 11 \% et le solaire à 2 \%. 1) Calcule les angles des quatre secteurs. 2) Décris en deux phrases l'évolution par rapport à une année normale.
\tcblower
\textbf{Solution détaillée.}
1) Conversion des pourcentages en angles ($\alpha = p \times 3{,}6$) :
\begin{align*}
\text{Hydroélectricité : } & 60 \times 3{,}6 = 216^{\circ} \\
\text{Thermique : } & 27 \times 3{,}6 = 97{,}2^{\circ} \approx 97^{\circ} \\
\text{Gaz : } & 11 \times 3{,}6 = 39{,}6^{\circ} \approx 40^{\circ} \\
\text{Solaire : } & 2 \times 3{,}6 = 7{,}2^{\circ} \approx 7^{\circ}
\end{align*}
Vérification : $216 + 97 + 40 + 7 = 360^{\circ}$. Le total est correct.
2) \textbf{Évolution} : par rapport à une année normale, la part de l'hydroélectricité recule fortement (de 70 \% à 60 \%, soit $-10$ points) parce que les retenues des barrages sont mal remplies. Ce recul est compensé par le thermique (de 20 \% à 27 \%), plus coûteux et polluant, ce qui confirme que la sécheresse fragilise tout le système électrique camerounais.
\end{exoresolu}

\begin{exercice}[À toi de jouer]
On donne la répartition suivante pour un pays voisin : hydroélectricité 45 \%, gaz 30 \%, thermique 20 \%, solaire 5 \%. 1) Construis le tableau de conversion (parts et angles). 2) Trace le camembert au compas et au rapporteur, avec titre, légende et source fictive. 3) Compare cette structure à celle du Cameroun en trois lignes : quel pays est le plus vulnérable à la sécheresse ? Justifie.
\end{exercice}

\begin{corrige}[Exercice — éléments de correction]
1) Angles : hydroélectricité $45 \times 3{,}6 = 162^{\circ}$ ; gaz $30 \times 3{,}6 = 108^{\circ}$ ; thermique $20 \times 3{,}6 = 72^{\circ}$ ; solaire $5 \times 3{,}6 = 18^{\circ}$. Total : $162 + 108 + 72 + 18 = 360^{\circ}$. 2) Le tracé part du rayon vertical, sens horaire, avec quatre couleurs distinctes, un titre (« Production électrique par source »), une légende et une source. 3) Le Cameroun (70 \% d'hydroélectricité) est plus vulnérable à la sécheresse que ce pays (45 \%), car sa production dépend davantage du remplissage des barrages ; le pays voisin, plus diversifié, peut compenser par le gaz.
\end{corrige}

\textbf{Conclusion du TP.} Le camembert que tu viens de construire résume à lui seul le paradoxe énergétique camerounais : un potentiel immense (hydraulique, solaire, gaz), mais une production concentrée sur une seule source, l'eau. Cette concentration explique les délestages des saisons sèches et la fragilité des zones industrielles étudiées dans la section précédente. Diversifier le mix énergétique — achever Nachtigal, développer le gaz de Kribi et les centrales solaires du Nord — est donc une condition de l'émergence industrielle du Cameroun. Retiens la démarche : \emph{données $\rightarrow$ calcul des angles $\rightarrow$ tracé légendé $\rightarrow$ interprétation géographique}. Elle te servira dans toutes les épreuves de construction graphique du Baccalauréat.

\newpage

\coursec{Activité d'intégration}

\courssub{Objectifs de l'activité}

À l'issue de cette activité d'intégration, tu seras capable de :

\begin{itemize}
\item \cle{observer} et \cle{localiser} sur une carte muette les principaux sites de production d'énergie et les gisements miniers du Cameroun ;
\item \cle{construire} un diagramme circulaire (camembert) à partir d'un tableau statistique, en appliquant la règle de proportionnalité des angles ;
\item \cle{classer} des entreprises selon les trois types d'industries étudiés (industries lourdes, agro-industries, PME) ;
\item \cle{interpréter} des documents variés (carte, tableaux, photographies, extrait de presse) pour construire une argumentation géographique ;
\item \cle{rédiger} un texte structuré expliquant le paradoxe entre l'énorme potentiel du Cameroun et la fragilité de son industrie, en proposant des solutions.
\end{itemize}

\courssub{Situation}

\textbf{Un paradoxe camerounais.} En janvier 2024, la société ENEO, concessionnaire de la distribution d'électricité au Cameroun, annonce de nouveaux programmes de \cle{délestage} : plusieurs quartiers de Douala et de Yaoundé sont privés de courant quatre à huit heures par jour. La même semaine, un journal économique titre : « Le Cameroun, deuxième potentiel hydroélectrique d'Afrique subsaharienne, peine à éclairer ses villes ». Comment un pays doté d'un tel potentiel énergétique et de réserves minières considérables peut-il connaître des coupures d'électricité chroniques et une industrie fragile ? C'est la question à laquelle ton groupe d'élèves de Terminale doit répondre, à la manière de géographes-conseillers auprès du Ministère de l'Industrie.

\textbf{Le dossier documentaire remis à chaque groupe.}

\emph{Document 1 — Carte muette du Cameroun} : contour du pays, fleuves Sanaga, Bénoué et Wouri, villes de Yaoundé, Douala, Edéa, Ngaoundéré, Kribi, Garoua.

\emph{Document 2 — Production d'électricité au Cameroun par source} (ordres de grandeur, début des années 2020, capacité installée totale d'environ \SI{1500}{MW} à \SI{1600}{MW}) :

\begin{center}
\begin{tabularx}{\linewidth}{|l|X|X|}\hline
\rowcolor{popblueL} \textbf{Source d'énergie} & \textbf{Part approximative} & \textbf{Exemples d'installations} \\ \hline
Hydroélectricité (barrages) & 73 \% & Songloulou, Edéa, Lagdo, Memve'ele \\ \hline
Centrales thermiques (fioul) & 15 \% & centrales ENEO des grandes villes \\ \hline
Centrales à gaz & 11 \% & Kribi Power (gaz de la Sanaga Sud) \\ \hline
Solaire et autres & 1 \% & centrales solaires de Guider, Djoum \\ \hline
\end{tabularx}
\end{center}

\emph{Document 3 — Réserves minières} : bauxite de Minim-Martap et Ngaoundal (plus d'\num{1} milliard de tonnes estimées, parmi les premières réserves mondiales non exploitées) ; fer de Mbalam (plusieurs centaines de millions de tonnes selon les estimations, à la frontière congolaise) ; rutile d'Akonolinga ; nickel-cobalt de Nkamouna ; indices de terres rares dans l'Adamaoua.

\emph{Document 4 — Deux photographies} : (a) le barrage hydroélectrique de Songloulou sur la Sanaga (environ \SI{384}{MW}, la plus grande centrale du pays) ; (b) la friche d'une ancienne usine textile de Douala, fermée dans les années 1990, bâtiments délabrés et machines rouillées.

\emph{Document 5 — Extrait de presse} : « Malgré la mise en service du barrage de Memve'ele (environ \SI{211}{MW}), les délestages persistent. Les industriels de la zone de Bonabéri estiment leurs pertes à plusieurs milliards de FCFA par an, entre groupes électrogènes, matières premières gâchées et retards de production. »

\emph{Document 6 — Liste de dix entreprises à classer} : ALUCAM (aluminium, Edéa), CIMENCAM (ciment, Douala/Figuil), DANGOTE Cement (Douala), SOSUCAM (sucre, Mbandjock), les Brasseries du Cameroun (SABC), la CDC (bananes et palmiers, Sud-Ouest), la SOCAPALM (huile de palme), une menuiserie de Yaoundé employant 12 personnes, un atelier de transformation du manioc de Bafoussam (8 employés), une imprimerie de Buea (15 employés).

\courssub{Consignes}

Travail en groupes de quatre à six élèves, en deux heures.

\begin{enumerate}
\item \textbf{Localiser et légender.} Sur la carte muette (document 1), place et légende : les barrages de Songloulou, Edéa, Lagdo et Memve'ele ; la centrale à gaz de Kribi ; les gisements de bauxite (Minim-Martap) et de fer (Mbalam). Rédige une légende organisée en deux rubriques (« énergie » et « mines »).
\item \textbf{Construire le camembert.} À partir du document 2, calcule l'angle de chaque secteur du diagramme circulaire (rappel : angle $= \dfrac{\text{part en \%}}{100} \times 360°$), puis construis le camembert de la production d'électricité camerounaise. Donne-lui un titre et une légende.
\item \textbf{Classer.} Répartis les dix entreprises du document 6 en trois colonnes : industries lourdes, agro-industries, PME. Justifie chaque classement en une phrase (critères : capital, poids de l'investissement, lien avec l'agriculture, effectif).
\item \textbf{Observer et interpréter.} Compare les deux photographies du document 4 : que révèlent-elles, respectivement, du potentiel et de la crise industrielle du Cameroun ?
\item \textbf{Argumenter.} En mobilisant l'ensemble des documents, rédige en dix lignes maximum un texte intitulé « Le paradoxe industriel camerounais » expliquant pourquoi un pays si riche en potentiel connaît des coupures d'électricité et une industrie fragile. Termine par deux solutions concrètes.
\item \textbf{Restitution.} Chaque groupe présente sa carte, son camembert et son texte en cinq minutes ; la correction collective est construite au tableau.
\end{enumerate}

\courssub{Ressources à mobiliser}

\begin{aretenir}[Ce qu'il faut avoir en tête]
\begin{itemize}
\item Le \cle{potentiel énergétique} : fleuve Sanaga (débit régulier, dénivelés favorables), gaz naturel de Kribi, fort ensoleillement du Nord ; le Cameroun possède le deuxième potentiel hydroélectrique d'Afrique subsaharienne, estimé à environ \SI{20000}{MW}, dont moins de 10 \% est exploité.
\item Les \cle{réserves minières} : bauxite (Minim-Martap, Ngaoundal), fer (Mbalam), nickel-cobalt (Nkamouna), terres rares — encore très peu exploitées.
\item Les \cle{types d'industries} : industries lourdes minoritaires (ALUCAM, CIMENCAM, DANGOTE), prédominance des agro-industries (SOSUCAM, SABC, SOCAPALM, CDC) et des PME.
\item La \cle{crise industrielle} des années 1980-1990 : fermeture de nombreuses unités (textile, cuir, assemblage), friches industrielles.
\item La formule du camembert : $\text{angle du secteur} = \dfrac{\text{effectif ou part}}{\text{total}} \times 360°$.
\end{itemize}
\end{aretenir}

\courssub{Démarche et corrigé commenté}

\begin{exoresolu}[Diagnostic industriel du Cameroun : corrigé guidé]
Énoncé : réaliser les six consignes du dossier documentaire et rédiger le texte « Le paradoxe industriel camerounais ».

\tcblower

\textbf{Étape 1 — Localisation et légende de la carte.} On repère d'abord les fleuves : la Sanaga traverse le Centre (barrages d'Edéa puis de Songloulou, en amont), le Wouri arrose Douala. Lagdo se situe sur la Bénoué au Nord (région du Nord, près de Garoua) ; Memve'ele sur la Ntem, au Sud. La centrale à gaz de Kribi est sur la côte, au sud de Douala. Minim-Martap se place sur le plateau de l'Adamaoua, près de Ngaoundéré ; Mbalam à l'extrême sud-est, à la frontière du Congo. La légende comporte deux rubriques : « Énergie » (symboles barrage, centrale à gaz) et « Mines » (symboles bauxite, fer). \emph{Erreur à éviter : placer Lagdo sur la Sanaga — c'est le piège classique.}

\begin{popfigure}
\begin{center}
\includegraphics[width=\linewidth,height=9.5cm,keepaspectratio]{N01/cmr_map.jpg}
\end{center}
\legende{Fond de carte : le Cameroun (relief ombré, régions, villes, fleuves, routes et voies ferrées ; on y voit la retenue de Lagdo sur la Bénoué, au nord, et la Sanaga, qui descend du Centre vers la côte près d'Édéa). Le croquis attendu situe sur ce fond : 1. les barrages de Songloulou et d'Édéa sur la Sanaga, de Lagdo sur la Bénoué et de Memve'ele sur la Ntem ; 2. la centrale à gaz de Kribi, sur le littoral ; 3. les gisements de bauxite de Minim-Martap (Adamaoua) et de fer de Mbalam (extrême Sud-Est). \\ Source : Wikimedia Commons, CIA, domaine public.}
\end{popfigure}

\textbf{Étape 2 — Construction du camembert.} On applique la formule de proportionnalité :
\begin{align*}
\text{Hydroélectricité : } & \frac{73}{100}\times 360° = 262{,}8° \approx 263° \\
\text{Thermique (fioul) : } & \frac{15}{100}\times 360° = 54° \\
\text{Gaz : } & \frac{11}{100}\times 360° = 39{,}6° \approx 40° \\
\text{Solaire : } & \frac{1}{100}\times 360° = 3{,}6° \approx 4°
\end{align*}
Vérification : $262{,}8 + 54 + 39{,}6 + 3{,}6 = 360°$ exactement ; avec les valeurs arrondies au degré, on obtient $263 + 54 + 40 + 4 = 361°$, écart d'un degré dû aux arrondis, qu'il faut signaler sur la copie. On trace au rapporteur, du plus grand secteur au plus petit, avec titre (« Production d'électricité au Cameroun par source, début des années 2020 ») et légende colorée. Le camembert montre l'écrasante domination de l'hydroélectricité : c'est une force (énergie renouvelable et peu chère) mais aussi une \cle{fragilité} : en saison sèche prolongée, la production chute.

\begin{popfigure}
\begin{center}
\begin{tikzpicture}[scale=1.1]
\fill[popblue] (0,0) -- (90:2) arc (90:90-262.8:2) -- cycle;
\fill[poporange] (0,0) -- (90-262.8:2) arc (90-262.8:90-316.8:2) -- cycle;
\fill[popgreen] (0,0) -- (90-316.8:2) arc (90-316.8:90-356.4:2) -- cycle;
\fill[popgold] (0,0) -- (90-356.4:2) arc (90-356.4:90-360:2) -- cycle;
\draw (0,0) circle (2);
\node[text=white,font=\bfseries] at (-1.1,0.8) {73 \%};
\node[text=poporange,font=\bfseries] at (1.15,-0.9) {15 \%};
\node[text=popgreen,font=\bfseries] at (0.75,-1.55) {11 \%};
\node[text=popgold,font=\bfseries\footnotesize] at (1.5,1.35) {1 \%};
\node[anchor=west,font=\footnotesize] at (2.5,1.2) {\textcolor{popblue}{$\blacksquare$} Hydroélectricité (barrages)};
\node[anchor=west,font=\footnotesize] at (2.5,0.7) {\textcolor{poporange}{$\blacksquare$} Thermique (fioul)};
\node[anchor=west,font=\footnotesize] at (2.5,0.2) {\textcolor{popgreen}{$\blacksquare$} Centrales à gaz};
\node[anchor=west,font=\footnotesize] at (2.5,-0.3) {\textcolor{popgold}{$\blacksquare$} Solaire et autres};
\end{tikzpicture}
\end{center}
\legende{Camembert attendu : production d'électricité au Cameroun par source (parts approximatives, début des années 2020). Le secteur hydroélectrique (262,8°, soit environ 263°) écrase toutes les autres sources.}
\end{popfigure}

\textbf{Étape 3 — Classement des entreprises.}
\begin{center}
\begin{tabularx}{\linewidth}{|l|X|X|}\hline
\rowcolor{popblueL} \textbf{Industries lourdes} & \textbf{Agro-industries} & \textbf{PME} \\ \hline
ALUCAM, CIMENCAM, DANGOTE Cement & SOSUCAM, SABC, CDC, SOCAPALM & menuiserie de Yaoundé, atelier de manioc de Bafoussam, imprimerie de Buea \\ \hline
\end{tabularx}
\end{center}
Justifications : ALUCAM, CIMENCAM et DANGOTE exigent des capitaux très lourds, de grandes quantités d'énergie et transforment des minerais ou des matières minérales (alumine importée pour ALUCAM, calcaire et clinker pour les cimenteries). SOSUCAM, SABC, CDC et SOCAPALM transforment des produits agricoles (canne à sucre, céréales et fruits, banane, palme) : ce sont des agro-industries. Les trois dernières unités emploient moins de 20 personnes avec des capitaux modestes : ce sont des PME. On retrouve la structure de l'industrie camerounaise : industries lourdes \cle{minoritaires}, prédominance des \cle{agro-industries} et des \cle{PME}.

\textbf{Étape 4 — Comparaison des photographies.} La photo (a) montre un ouvrage moderne et puissant : elle illustre l'\cle{énorme potentiel énergétique} du pays et sa mise en valeur partielle. La photo (b) montre une friche : elle illustre l'\cle{élan industriel brisé par la crise} des années 1980-1990 (chute des cours des matières premières, endettement, programmes d'ajustement structurel), qui a provoqué la fermeture de nombreuses unités, notamment dans le textile.

\textbf{Étape 5 — Texte argumenté (corrigé-type en dix lignes).} « Le Cameroun présente un paradoxe frappant. Son potentiel hydroélectrique, deuxième d'Afrique subsaharienne (environ \SI{20000}{MW}), n'est exploité qu'à moins de 10 \%, et ses immenses réserves de bauxite (Minim-Martap) et de fer (Mbalam) restent en sous-sol faute d'infrastructures de transport et de capitaux. La production électrique, dominée à 73 \% par l'hydroélectricité, dépend des aléas climatiques ; le réseau, vieillissant, enregistre des pertes importantes, d'où les délestages qui pénalisent les usines. L'industrie, fragilisée par la crise des années 1980-1990, repose surtout sur des agro-industries et des PME, avec très peu d'industries lourdes. Deux solutions : accélérer l'exploitation minière et la construction de barrages grâce aux partenariats public-privé, et réhabiliter le réseau de distribution tout en développant le solaire dans le Nord. »

\textbf{Étape 6 — Restitution.} La correction collective au tableau confronte les cartes et camemberts des groupes, harmonise les légendes et fait émerger la notion-clé : le \cle{décalage entre potentiel et réalisations}.
\end{exoresolu}

\begin{attention}[Pièges fréquents dans cette activité]
\begin{itemize}
\item Oublier de vérifier que la somme des angles du camembert vaut 360° (les arrondis doivent être expliqués).
\item Confondre \emph{capacité installée} (en MW) et \emph{production réelle} (en MWh) : ce ne sont pas les mêmes grandeurs.
\item Placer le barrage de Lagdo sur la Sanaga : il est sur la \cle{Bénoué}.
\item Classer la SABC ou SOSUCAM en « industrie lourde » parce que ce sont de grandes entreprises : le critère décisif est la \cle{nature de la matière transformée} (produits agricoles $\Rightarrow$ agro-industrie).
\item Rédiger un texte descriptif au lieu d'un texte \emph{explicatif} : il faut répondre au « pourquoi » du paradoxe, pas seulement le constater.
\end{itemize}
\end{attention}

\courssub{Bilan de l'activité}

Cette activité a permis de mobiliser simultanément les quatre savoir-faire du programme : \cle{observer} (photographies du barrage et de la friche), \cle{localiser} (carte muette légendée), \cle{classer} (typologie des entreprises) et \cle{construire} un diagramme (camembert de la production d'énergie), avant de les réunir dans une \cle{interprétation} rédigée. Elle met en évidence l'idée maîtresse de la leçon : le Cameroun ne manque ni de ressources énergétiques ni de richesses minières, mais souffre d'un déficit de mise en valeur — infrastructures insuffisantes, réseau électrique vieillissant, capitaux limités, séquelles de la crise des années 1980-1990. Le paradoxe « potentiel immense / industrie fragile » n'est donc pas une fatalité : il appelle des politiques d'investissement, de partenariat et de diversification énergétique, dont tu es désormais capable de discuter en géographe averti.

\newpage

\coursec{Méthodes et exercices résolus}

\coursec{L'industrie}

\begin{definition}[Notions clés]
\textbf{Industrie} : activité économique qui transforme des matières premières en produits finis ou semi-finis, à l'aide d'énergie, de machines et de main-d'œuvre.
\textbf{Zone industrielle} : espace urbanisé, aménagé et équipé (voirie, réseaux, énergie) destiné à accueillir des unités de production.
\textbf{Agro-industrie} : industrie de transformation des produits agricoles (brasserie, huilerie, savonnerie, minoterie, transformation du cacao, du café, du coton).
\textbf{Industrie lourde} : industrie de base, forte consommatrice d'énergie et de matières premières lourdes (métallurgie, cimenterie, chimie lourde), à fort capital fixe.
\textbf{PME/PMI} : Petites et Moyennes Entreprises / Industries : effectif < 200 (souvent < 50 au Cameroun), capital modeste, forte intensité de main-d'œuvre.
\end{definition}

\begin{aretenir}[Les trois piliers de l'industrie camerounaise]
\begin{enumerate}
\item \textbf{Un énorme potentiel énergétique} : hydroélectricité (Sanaga : Edéa, Songloulou, Nachtigal ; Bénoué : Lagdo ; Ntem : Memve'ele), gaz (Logbaba, Kribi), solaire (Nord), biomasse. Potentiel hydroélectrique théorique $\approx$ 20 GW (2\up{ème} Afrique), peu mobilisé ($\sim$ 10--15 \%).
\item \textbf{D'importantes réserves minières} : bauxite (Minim-Martap, Ngaoundal $\sim$ 1 Gt), fer (Mbalam, Kribi $\sim$ 600 Mt), rutile, terres rares (Akonolinga), cobalt, nickel, or. Exploitation quasi nulle à ce jour (projets en cours).
\item \textbf{Une structure industrielle duale} : industries lourdes \textbf{minoritaires} (ALUCAM, CIMENCAM, DANGOTE) ; \textbf{prédominance des agro-industries et des PME} (tissu de transformation locale, fort emploi, faible capital) ; concentration spatiale forte (Douala-Bassa/Bonabéri, Edéa, Yaoundé, quelques pôles secondaires).
\end{enumerate}
\end{aretenir}

\begin{savaistu}[Actualité 2023--2025 : la relance énergétique et minière]
\begin{itemize}
\item \textbf{Barrages récents} : \textbf{Nachtigal} (420 MW, Sanaga, mis en service progressif 2023--2024) et \textbf{Memve'ele} (211 MW, Ntem, 2019/2020) portent l'hydroélectricité installée au-delà de 1300 MW. Projet \textbf{Grand Eweng} (800--1200 MW) à l'étude.
\item \textbf{Gaz} : centrale de \textbf{Kribi} (216 MW, gaz de Logbaba) et extension \textbf{Logbaba} ; projet GNL (Gaz Naturel Liquéfié) à Kribi.
\item \textbf{Mines} : convention \textbf{Minim-Martap} (bauxite, consortium chinois) et \textbf{Mbalam} (fer, Sundance/Australia puis consortium) conditionnées à la construction du \textbf{chemin de fer minier} (500+ km) et du \textbf{port en eau profonde de Kribi} (phase 2).
\item \textbf{Industrie} : nouvelle cimenterie \textbf{CIMENCAM} à Figuil (four 3000 t/j), extension \textbf{DANGOTE} Douala (clinker), projet usine transformation cacao (Penja), bois (Bertoua), coton (Garoua).
\end{itemize}
\end{savaistu}

\courssub{Méthode 1 : Observer, localiser et décrire une zone industrielle sur un document}

\begin{methode}[Décrire une zone industrielle à partir d'une carte ou d'une photographie]
\begin{enumerate}
\item \textbf{Nommer et situer} : donne le nom de la zone (ex. zone industrielle de Bonabéri-Douala) et localise-la (ville, région, position par rapport au fleuve, au port, aux axes).
\item \textbf{Décrire ce qu'on voit} : usines, entrepôts, cheminées, voies ferrées, routes, port, réservoirs. Reste factuel : ne décris que ce qui est visible ou indiqué sur le document.
\item \textbf{Identifier les facteurs de localisation} : proximité d'un port (importation de matières premières, exportation), d'un axe routier ou ferroviaire, d'une source d'énergie (barrage), d'une main-d'œuvre urbaine, d'un marché de consommation.
\item \textbf{Classer les activités observées} : industrie lourde (cimenterie, aluminium), agro-industrie (brasserie, huilerie, savonnerie), PME (menuiserie métallique, plasturgie).
\item \textbf{Conclure par une phrase de synthèse} qui relie la localisation à l'intégration nationale (désenclavement, emploi, transformation locale des matières premières).
\end{enumerate}
\textbf{Réflexe} : une description n'est ni une liste sans ordre ni une explication. On va du général (localisation) au particulier (détails), puis on explique.
\end{methode}

\begin{exoresolu}[La zone industrielle de Douala-Bassa/Bonabéri]
\textbf{Énoncé (type Bac)} : À partir d'une carte de la zone industrielle de Douala montrant le port, le Wouri, la voie ferrée Douala--Ngaoundéré, les usines ALUCAM (Edéa, à 70 km au sud), CIMENCAM (Nomayos/Figuil), les brasseries SABC et de nombreuses PME concentrées à Bonabéri et Bassa : décris cette zone industrielle et explique sa localisation. (10 points)
\tcblower
\textbf{Solution détaillée}

\textbf{1. Nommer et situer.} La zone industrielle de Douala-Bassa/Bonabéri se situe dans la ville de Douala, capitale économique du Cameroun, dans la région du Littoral. Elle s'étend de part et d'autre du fleuve Wouri, à proximité immédiate du port autonome de Douala, premier port du pays (environ 95~\% du trafic portuaire national, ordre de grandeur prudent).

\textbf{2. Décrire.} On observe une concentration dense d'usines et d'entrepôts : brasseries (SABC), savonneries et huileries (complexes agro-industriels), cimenteries de broyage (CIMENCAM, DANGOTE Cement à Douala), unités de plasturgie et de menuiserie métallique (PME). La zone est traversée par la voie ferrée Transcamerounais (axe Douala--Ngaoundéré) et par de grands boulevards ; elle jouxte le port et ses terminaux.

\textbf{3. Facteurs de localisation.}
\begin{itemize}
\item Le \textbf{port} permet d'importer les matières premières (clinker pour les cimenteries, blé, matériel) et d'exporter les produits finis.
\item Le \textbf{chemin de fer} et la route relient la zone à l'hinterland (Ouest, Centre, Nord) : approvisionnement en produits agricoles et écoulement des marchandises.
\item La \textbf{main-d'œuvre} et le \textbf{marché de consommation} : Douala compte plus de 3 millions d'habitants (ordre de grandeur).
\item L'\textbf{énergie} : la centrale hydroélectrique d'Edéa (environ 263 MW) et les centrales thermiques alimentent la zone.
\end{itemize}

\textbf{4. Classement.} Les industries lourdes (cimenteries, métallurgie de l'aluminium via ALUCAM à Edéa) y sont minoritaires ; les agro-industries et les PME y sont largement prédominantes, ce qui illustre la structure industrielle camerounaise.

\textbf{5. Conclusion.} La zone industrielle de Douala-Bassa/Bonabéri, polarisée par le port et le rail, concentre l'essentiel de l'appareil industriel camerounais ; elle joue un rôle moteur dans l'intégration nationale en transformant localement les matières premières et en irriguant l'hinterland, mais cette concentration renforce le déséquilibre Littoral--intérieur.
\end{exoresolu}

\courssub{Méthode 2 : Construire et légender un croquis des industries du Cameroun}

\begin{methode}[Réaliser un croquis de localisation industrielle]
\begin{enumerate}
\item \textbf{Tracer le contour} du Cameroun en polygone simplifié (forme de triangle allongé nord-sud), sans prétendre à l'exactitude.
\item \textbf{Placer le nord} (flèche), le \textbf{titre} en haut et une \textbf{échelle indicative}.
\item \textbf{Placer les villes-industries} par des points nommés : Douala, Edéa, Yaoundé, Garoua, Ngaoundéré, Bafoussam, Figuil, Kribi.
\item \textbf{Figurer les axes} : Transcamerounais (Douala--Yaoundé--Ngaoundéré) en trait plein, routes principales en trait fin, ports (Douala, Kribi) par une ancre ou un symbole.
\item \textbf{Choisir des symboles simples et proportionnés} : carré = industrie lourde (ALUCAM Edéa, CIMENCAM Figuil/Nomayos, DANGOTE Douala), cercle = agro-industrie, triangle = PME/artisanat, losange = barrage (Edéa, Songloulou, Lagdo, Nachtigal, Memve'ele).
\item \textbf{Rédiger la légende} dans un cadre, dans le même ordre que les symboles, numérotée si nécessaire.
\end{enumerate}
\textbf{Réflexe} : un croquis sans titre, sans nord ou sans légende est systématiquement pénalisé au Bac.
\end{methode}

\begin{exoresolu}[Croquis : les grandes unités industrielles du Cameroun]
\textbf{Énoncé (type Bac)} : Construis un croquis du Cameroun localisant les principales industries lourdes et les principaux barrages hydroélectriques. (10 points)
\tcblower
\textbf{Solution détaillée}

\textbf{Démarche.} On trace le contour simplifié, on place le nord, puis on reporte : ALUCAM à Edéa (aluminium), CIMENCAM à Figuil (Nord) et Nomayos (Centre), DANGOTE Cement à Douala, les barrages d'Edéa et de Songloulou sur la Sanaga, de Lagdo sur la Bénoué, de Nachtigal (Sanaga) et Memve'ele (Ntem). On relie par le Transcamerounais.

\begin{popfigure}
\begin{center}
\includegraphics[width=\linewidth,height=9.5cm,keepaspectratio]{N01/cmr_map.jpg}
\end{center}
\legende{Fond de carte pour le croquis demandé. Le croquis attendu situe : 1. ALUCAM (Édéa) ; 2. CIMENCAM (Figuil, Nomayos) ; 3. DANGOTE (Douala) ; 4. les barrages d'Édéa, Songloulou et Nachtigal (Sanaga), de Lagdo (Bénoué) et de Memve'ele (Ntem). Les fleuves, les villes et la voie ferrée sont visibles sur la carte ; les industries et les barrages sont à ajouter par l'élève. \\ Source : Wikimedia Commons, CIA, domaine public.}
\end{popfigure}

\textbf{Justification des choix.} Les industries lourdes sont figurées par un symbole unique (carré) car elles sont minoritaires (3 grands groupes) ; les barrages sont placés sur les fleuves Sanaga (Edéa, Songloulou, Nachtigal), Bénoué (Lagdo) et Ntem (Memve'ele) car l'hydroélectricité fournit l'essentiel de l'énergie industrielle.

\textbf{Conclusion.} Le croquis met en évidence la concentration industrielle sur l'axe Douala--Edéa--Yaoundé et la dépendance de l'industrie à l'énergie hydroélectrique de la Sanaga.
\end{exoresolu}

\courssub{Méthode 3 : Inventorier, classer et interpréter un tableau statistique}

\begin{methode}[Exploiter un tableau de données industrielles]
\begin{enumerate}
\item \textbf{Lire le titre, l'unité, la source et la date} du tableau avant tout calcul.
\item \textbf{Inventorier} : repère les valeurs extrêmes (maximum, minimum) et les ordres de grandeur.
\item \textbf{Classer} : range les catégories par ordre décroissant (ou croissant) pour dégager la hiérarchie.
\item \textbf{Calculer les indicateurs} :
\begin{itemize}
\item part en pourcentage : $\text{part} = \dfrac{\text{effectif de la catégorie}}{\text{effectif total}} \times 100$ ;
\item taux de variation : $t = \dfrac{V_f - V_i}{V_i} \times 100$ (en \text{\%}).
\end{itemize}
\item \textbf{Interpréter} : formule une phrase par constat majeur (dominance, faiblesse, évolution), en citant les chiffres.
\item \textbf{Relier au cours} : conclus sur la structure industrielle camerounaise (prédominance des agro-industries et des PME, industries lourdes minoritaires).
\end{enumerate}
\textbf{Réflexe} : tout calcul doit être suivi d'une phrase d'interprétation ; un chiffre non commenté ne rapporte pas de point.
\end{methode}

\begin{exoresolu}[Structure des entreprises industrielles camerounaises]
\textbf{Énoncé (type Bac)} : Le tableau ci-dessous donne la répartition approximative des entreprises industrielles formelles du Cameroun par taille (ordres de grandeur, source : INS, annuaire des entreprises).

\begin{center}
\begin{tabularx}{\linewidth}{|l|X|X|}\hline
\rowcolor{popblueL} \textbf{Catégorie} & \textbf{Nombre d'entreprises} & \textbf{Exemples} \\ \hline
Grandes entreprises (GE) & 150 & ALUCAM, CIMENCAM, SABC \\ \hline
Entreprises moyennes (EM) & 350 & minoteries, plasturgie \\ \hline
Petites entreprises (PE) & 1500 & savonneries, menuiseries \\ \hline
\end{tabularx}
\end{center}

1) Calcule la part de chaque catégorie dans le total. 2) Interprète la structure obtenue. (10 points)
\tcblower
\textbf{Solution détaillée}

\textbf{1. Lecture.} Le tableau comporte trois catégories ; le total se calcule :
\[ 150 + 350 + 1500 = 2000 \text{ entreprises.} \]

\textbf{2. Calcul des parts.}
\begin{align*}
\text{Part des GE} &= \frac{150}{2000} \times 100 = 7{,}5\,\text{\%} \\
\text{Part des EM} &= \frac{350}{2000} \times 100 = 17{,}5\,\text{\%} \\
\text{Part des PE} &= \frac{1500}{2000} \times 100 = 75\,\text{\%}
\end{align*}
Vérification : $7{,}5 + 17{,}5 + 75 = 100\,\text{\%}$. Le compte est bon.

\textbf{3. Classement.} Par ordre décroissant : PE (75 \%), EM (17,5 \%), GE (7,5 \%).

\textbf{4. Interprétation.} Les petites entreprises dominent écrasamment le tissu industriel camerounais (trois entreprises sur quatre), tandis que les grandes entreprises, souvent des filiales de groupes internationaux (ALUCAM, CIMENCAM, DANGOTE), ne représentent que 7,5 \% des unités. Cette structure confirme la \textbf{prédominance des PME et des agro-industries} et le caractère \textbf{minoritaire des industries lourdes} au Cameroun.

\textbf{Conclusion.} L'industrie camerounaise est un tissu de PME peu capitalistiques, dominé par quelques grandes unités étrangères ou parapubliques : c'est une industrie encore fragile, dont le renforcement passe par la transformation locale des matières premières.
\end{exoresolu}

\courssub{Méthode 4 : TP 3 — Construire un camembert (diagramme circulaire)}

\begin{methode}[Construire un diagramme circulaire de la production d'énergie]
\begin{enumerate}
\item \textbf{Totaliser} les valeurs (production ou puissance installée par source).
\item \textbf{Calculer la part} de chaque source : $\text{part} = \dfrac{\text{valeur}}{\text{total}} \times 100$ (en \text{\%}).
\item \textbf{Convertir en angles} : le cercle complet vaut $360^{\circ}$, donc
\[ \text{angle} = \frac{\text{valeur}}{\text{total}} \times 360 \quad \text{ou} \quad \text{angle} = \text{part} \times 3{,}6. \]
\item \textbf{Vérifier} que la somme des angles vaut $360^{\circ}$ (tolérance d'arrondi : $\pm 1^{\circ}$).
\item \textbf{Tracer} : au rapporteur, reporte chaque angle à partir du rayon vertical (12 h), dans le sens des aiguilles d'une montre, par ordre décroissant.
\item \textbf{Colorier} chaque secteur d'une couleur distincte, \textbf{titrer} le graphique et rédiger la \textbf{légende} avec les pourcentages.
\end{enumerate}
\textbf{Réflexe} : annonce toujours la formule de conversion (1 \% = 3,6°) avant le calcul.
\end{methode}

\begin{exoresolu}[TP 3 : Camembert de la production d'électricité au Cameroun]
\textbf{Énoncé (type Bac / TP)} : La puissance électrique installée au Cameroun se répartit approximativement ainsi (ordres de grandeur, source : ENEO/ARSEL) : barrages hydroélectriques : 1100 MW ; centrales thermiques (fioul/diesel) : 450 MW ; centrales à gaz : 350 MW ; solaire : 100 MW. Construis le diagramme circulaire de cette répartition et commente-le. (10 points)
\tcblower
\textbf{Solution détaillée}

\textbf{1. Total.}
\[ 1100 + 450 + 350 + 100 = 2000 \text{ MW.} \]

\textbf{2. Parts et angles} (règle : 1 \% = 3,6°) :
\begin{align*}
\text{Hydroélectricité} &: \frac{1100}{2000} \times 100 = 55\,\text{\%} \quad ; \quad 55 \times 3{,}6 = 198^{\circ} \\
\text{Thermique} &: \frac{450}{2000} \times 100 = 22{,}5\,\text{\%} \quad ; \quad 22{,}5 \times 3{,}6 = 81^{\circ} \\
\text{Gaz} &: \frac{350}{2000} \times 100 = 17{,}5\,\text{\%} \quad ; \quad 17{,}5 \times 3{,}6 = 63^{\circ} \\
\text{Solaire} &: \frac{100}{2000} \times 100 = 5\,\text{\%} \quad ; \quad 5 \times 3{,}6 = 18^{\circ}
\end{align*}

\textbf{3. Vérification.} $198 + 81 + 63 + 18 = 360^{\circ}$. La construction est possible.

\textbf{4. Tracé.}

\begin{popfigure}
\begin{center}
\begin{tikzpicture}[scale=1.1]
% Secteurs : hydro 198°, thermique 81°, gaz 63°, solaire 18° (sens horaire depuis 90°)
\fill[popblue] (0,0) -- (90:2.2) arc (90:-108:2.2) -- cycle;
\fill[poporange] (0,0) -- (-108:2.2) arc (-108:-189:2.2) -- cycle;
\fill[popgold] (0,0) -- (-189:2.2) arc (-189:-252:2.2) -- cycle;
\fill[popgreen] (0,0) -- (-252:2.2) arc (-252:-270:2.2) -- cycle;
\draw[thick, popdark] (0,0) circle (2.2);
% Étiquettes
\node[white, font=\scriptsize\bfseries] at (-18:1.4) {Hydro 55~\%};
\node[white, font=\scriptsize\bfseries] at (-148:1.4) {Therm. 22,5~\%};
\node[popdark, font=\scriptsize\bfseries] at (-220:1.4) {Gaz 17,5~\%};
\node[white, font=\scriptsize\bfseries] at (-261:1.5) {5~\%};
% Légende
\begin{scope}[shift={(3.0,-1.6)}]
\draw[popline] (0,0) rectangle (4.2,3.2);
\node[font=\bfseries\small, anchor=west] at (0.15,2.9) {Légende};
\fill[popblue] (0.2,2.4) rectangle (0.45,2.6); \node[anchor=west, font=\scriptsize] at (0.6,2.5) {Barrages hydroélectriques};
\fill[poporange] (0.2,1.8) rectangle (0.45,2.0); \node[anchor=west, font=\scriptsize] at (0.6,1.9) {Centrales thermiques};
\fill[popgold] (0.2,1.2) rectangle (0.45,1.4); \node[anchor=west, font=\scriptsize] at (0.6,1.3) {Centrales à gaz (Kribi, Logbaba)};
\fill[popgreen] (0.2,0.6) rectangle (0.45,0.8); \node[anchor=west, font=\scriptsize] at (0.6,0.7) {Solaire (Guider, Maroua)};
\end{scope}
\node[font=\bfseries, popdark] at (0,2.8) {Puissance électrique installée au Cameroun par source};
\end{tikzpicture}
\end{center}
\legende{Diagramme circulaire de la puissance installée (total : 2000 MW, ordres de grandeur). Angles : hydro 198°, thermique 81°, gaz 63°, solaire 18°. Source : ENEO/ARSEL, valeurs arrondies.}
\end{popfigure}

\textbf{5. Commentaire.} L'hydroélectricité domine nettement avec 55 \% de la puissance installée (barrages d'Edéa, Songloulou, Nachtigal sur la Sanaga, Lagdo sur la Bénoué, Memve'ele sur le Ntem). Les centrales thermiques (22,5 \%) et à gaz (17,5 \%, notamment Kribi et Logbaba) complètent le dispositif, surtout en saison sèche. Le solaire reste marginal (5 \%) malgré un ensoleillement exceptionnel dans le Nord.

\textbf{Conclusion.} Le Cameroun dispose d'un mix énergétique dominé par l'hydraulique, reflet d'un énorme potentiel encore sous-exploité (le potentiel hydroélectrique national est estimé à environ 20 GW, dont moins de 10 \% est mobilisé).
\end{exoresolu}

\courssub{Méthode 5 : Commenter un document sur la crise industrielle}

\begin{methode}[Rédiger un commentaire de document (texte, tableau ou graphique)]
\begin{enumerate}
\item \textbf{Introduction} : présente le document (nature, source, date, thème) et annonce la problématique (ex. : comment la crise des années 1980-1990 a-t-elle brisé l'élan industriel camerounais ?).
\item \textbf{Analyse} : dégage 2 ou 3 idées principales, chacune appuyée par un chiffre ou un fait du document (fermetures d'unités, chute de la production, licenciements).
\item \textbf{Explication} : mobilise tes connaissances pour expliquer (chute des cours des matières premières, endettement, programmes d'ajustement structurel du FMI et de la Banque mondiale, privatisations, concurrence des importations).
\item \textbf{Conclusion} : bilan nuancé (reconversion, relance récente avec la cimenterie, les grands projets) et ouverture (industrialisation par la transformation locale).
\end{enumerate}
\textbf{Réflexe} : alterne toujours « ce que dit le document » et « ce que je sais » ; ne recopie jamais le texte.
\end{methode}

\begin{exoresolu}[La crise industrielle des années 1980-1990]
\textbf{Énoncé (type Bac)} : Document : « Entre 1986 et 1995, le Cameroun connaît la fermeture de nombreuses unités industrielles : filatures, huileries, sociétés parapubliques (SODECOTON affaiblie, CAMSUCO en faillite, CELLUCAM en démantèlement). La part de l'industrie manufacturière dans le PIB, qui atteignait environ 15 \% au milieu des années 1980, tombe sous les 10 \% au milieu des années 1990 (ordres de grandeur, Banque mondiale). » Commentez ce document en expliquant les causes et les conséquences de la crise industrielle. (10 points)
\tcblower
\textbf{Solution détaillée}

\textbf{Introduction.} Ce texte statistique, de source Banque mondiale, porte sur la crise industrielle camerounaise de 1986-1995. Problématique : comment un élan industriel prometteur, né des plans quinquennaux de l'indépendance, a-t-il été brisé en une décennie ?

\textbf{Analyse.} Deux constats majeurs :
\begin{itemize}
\item Des \textbf{fermetures massives} : le document cite CAMSUCO (sucre), CELLUCAM (cellulose) et des filatures et huileries ; le tissu parapublic, pilier de l'industrialisation, s'effondre.
\item Un \textbf{recul économique mesurable} : la part de l'industrie manufacturière dans le PIB passe d'environ 15 \% à moins de 10 \%, soit une baisse relative de
\[ \frac{10 - 15}{15} \times 100 \approx -33\,\text{\%} : \]
le poids de l'industrie a chuté d'environ un tiers en dix ans.
\end{itemize}

\textbf{Explication (connaissances).} Les causes sont externes (chute des cours du pétrole, du café et du cacao après 1985, surendettement) et internes (gestion déficitaire des sociétés d'État, coûts de production élevés). Les programmes d'ajustement structurel imposés par le FMI et la Banque mondiale (dévaluation du FCFA en 1994, privatisations, licenciements) achèvent les unités non rentables, tandis que la libéralisation expose les survivantes à la concurrence des produits importés.

\textbf{Conséquences.} Perte de milliers d'emplois industriels, désindustrialisation des villes moyennes, retour à l'exportation de matières premières brutes, essor du secteur informel.

\textbf{Conclusion.} La crise a brisé l'élan industriel camerounais et durablement fragilisé le tissu productif. La relance actuelle (cimenteries DANGOTE et CIMENCAM, projet de transformation locale du cacao et du bois) montre que la réindustrialisation passe par la valorisation des ressources nationales, au service de l'intégration nationale.
\end{exoresolu}

\courssub{Méthode 6 : Construire et lire un diagramme en barres comparatif}

\begin{methode}[Diagramme en barres : ressources minières ou production industrielle]
\begin{enumerate}
\item \textbf{Choisir l'échelle} : repère la valeur maximale, puis fixe une échelle simple (ex. 1 cm pour 5 unités) qui tient dans la feuille.
\item \textbf{Tracer les axes} : axe horizontal = catégories (minerais, régions, années) ; axe vertical = grandeur avec son \textbf{unité} obligatoire.
\item \textbf{Tracer des barres de largeur égale}, régulièrement espacées, de hauteur proportionnelle aux valeurs.
\item \textbf{Annoter} chaque barre de sa valeur, titrer le graphique et indiquer la source.
\item \textbf{Lire} : identifie le maximum, le minimum, les rapports (ex. « la bauxite représente X fois le fer »).
\end{enumerate}
\textbf{Réflexe} : un diagramme en barres compare des catégories ; pour montrer une évolution dans le temps, une courbe est souvent préférable.
\end{methode}

\begin{exoresolu}[Les réserves minières du Cameroun]
\textbf{Énoncé (type Bac)} : Les réserves estimées de minerais au Cameroun sont les suivantes (ordres de grandeur, sources ministérielles) : bauxite (Minim-Martap, Ngaoundal) : environ 1000 millions de tonnes ; fer (Mbalam, Kribi) : environ 600 millions de tonnes ; rutile (Akonolinga) : environ 3 millions de tonnes. Représente ces données par un diagramme en barres et tire-en deux enseignements. (10 points)
\tcblower
\textbf{Solution détaillée}

\textbf{1. Échelle.} Maximum : 1000 Mt. On choisit 1 cm pour 200 Mt, soit des hauteurs de 5 cm, 3 cm et 0,015 cm. La barre du rutile étant invisible à cette échelle, on l'indique par une valeur annotée (réflexe de rigueur : signaler toute valeur non représentable).

\textbf{2. Tracé.}

\begin{popfigure}
\begin{center}
\begin{tikzpicture}[scale=1]
\draw[-{Stealth}, thick] (0,0) -- (0,5.6) node[above, font=\scriptsize]{Réserves (Mt)};
\draw[-{Stealth}, thick] (0,0) -- (7.2,0);
% Barres : échelle 1 cm = 200 Mt pour lisibilité
\fill[popblue] (0.8,0) rectangle (2.0,5.0);
\fill[poporange] (3.0,0) rectangle (4.2,3.0);
\fill[popgreen] (5.2,0) rectangle (6.4,0.015);
% Valeurs
\node[above, font=\scriptsize\bfseries, popblue] at (1.4,5.0) {1000};
\node[above, font=\scriptsize\bfseries, poporange] at (3.6,3.0) {600};
\node[above, font=\scriptsize\bfseries, popgreen] at (5.8,0.05) {3};
% Catégories
\node[below, font=\scriptsize] at (1.4,0) {Bauxite};
\node[below, font=\scriptsize] at (3.6,0) {Fer};
\node[below, font=\scriptsize] at (5.8,0) {Rutile};
% Graduations
\foreach \y in {1,2,3,4,5} {\draw (-0.08,\y) -- (0.08,\y); \node[left, font=\scriptsize] at (-0.1,\y) {\pgfmathparse{int(\y*200)}\pgfmathresult};}
\node[font=\bfseries, popdark] at (3.6,6.2) {Principales réserves minières du Cameroun};
\end{tikzpicture}
\end{center}
\legende{Diagramme en barres des réserves estimées (millions de tonnes, ordres de grandeur). Échelle : 1 cm = 200 Mt. La barre du rutile (3 Mt) est trop faible pour être visible ; sa valeur est annotée.}
\end{popfigure}

\textbf{3. Lecture et enseignements.}
\begin{itemize}
\item \textbf{Hiérarchie marquée} : la bauxite domine (1000 Mt), devant le fer (600 Mt) ; le rapport est de $\dfrac{1000}{600} \approx 1{,}7$ : les réserves de bauxite sont environ 1,7 fois celles de fer.
\item \textbf{Contraste d'exploitation} : malgré ces réserves considérables (auxquelles s'ajoutent les terres rares d'Akonolinga), aucun de ces gisements n'est encore exploité à grande échelle ; seul l'aluminium est transformé à Edéa (ALUCAM) à partir d'alumine \emph{importée}, paradoxe qui illustre la faiblesse de la valorisation locale.
\end{itemize}

\textbf{Conclusion.} Le sous-sol camerounais recèle d'importantes réserves minières dont l'exploitation effective (projets de Minim-Martap et de Mbalam) constitue un enjeu majeur de la relance industrielle et de l'intégration nationale (chemin de fer minier, port de Kribi).
\end{exoresolu}

\begin{attention}[Les 5 erreurs qui coûtent des points au Bac]
\begin{enumerate}
\item \textbf{Camembert sans vérification} : ne pas contrôler que la somme des angles fait $360^{\circ}$. Une erreur d'arrondi non signalée (ex. $198 + 81 + 63 + 17 = 359^{\circ}$) doit être corrigée en ajustant le plus grand secteur. Retiens : $\text{angle} = \text{part} \times 3{,}6$.
\item \textbf{Croquis sans titre, sans nord ni légende} : c'est la cause première de perte de points en travaux pratiques. La légende doit reprendre \emph{tous} les symboles, dans l'ordre.
\item \textbf{Confondre puissance installée (MW) et production (MWh)} : ce sont deux grandeurs différentes. Cite toujours l'unité exacte du document et recopie-la sans la transformer.
\item \textbf{Décrire au lieu d'expliquer (ou l'inverse)} : à la consigne « décrire », reste factuel ; à « expliquez », mobilise les facteurs (port, énergie, main-d'œuvre, marché). Mélanger les deux consignes dans la même phrase affaiblit la copie.
\item \textbf{Erreurs de calcul de pourcentage} : oublier le $\times 100$, ou diviser le total par la partie au lieu de la partie par le total. Vérifie toujours que la somme des parts fait 100 \%, et qu'un taux de variation négatif est bien présenté comme une \textbf{baisse} (ex. $-33\,\text{\%}$ = baisse d'un tiers, et non « augmentation de $-33\,\text{\%}$ »).
\end{enumerate}
\end{attention}

\newpage

\coursec{Exercices}

\courssub{Je vérifie mes connaissances}

\begin{exercice}[QCM : l'industrie camerounaise]
Pour chaque question, choisis la (ou les) bonne(s) réponse(s) en justifiant brièvement.
\begin{enumerate}
\item Le premier barrage hydroélectrique du Cameroun, construit sur le fleuve Sanaga, est :
\begin{enumerate}
\item Song Loulou \item Lagdo \item Edéa \item Memve'ele
\end{enumerate}
\item La bauxite camerounaise se trouve principalement :
\begin{enumerate}
\item à Minim-Martap et Ngaoundal (Adamaoua) \item à Kribi (Sud) \item à Batouri (Est) \item à Wum (Nord-Ouest)
\end{enumerate}
\item ALUCAM est une entreprise de :
\begin{enumerate}
\item transformation du cacao \item production d'aluminium \item cimenterie \item brasserie
\end{enumerate}
\item Le potentiel hydroélectrique du Cameroun est estimé à environ :
\begin{enumerate}
\item 2 GW \item 20 GW \item 200 GW \item 2\,000 GW
\end{enumerate}
\end{enumerate}
\end{exercice}

\begin{exercice}[Vrai ou faux ?]
Indique si chaque affirmation est vraie ou fausse, puis justifie ta réponse en une ou deux phrases.
\begin{enumerate}
\item Les industries lourdes sont majoritaires dans le tissu industriel camerounais.
\item La crise économique des années 1980-1990 a entraîné la fermeture de nombreuses unités industrielles au Cameroun.
\item La centrale à gaz de Kribi utilise une énergie fossile produite sur le territoire camerounais.
\item Le Cameroun exploite déjà industriellement l'ensemble de ses réserves de fer et de bauxite.
\end{enumerate}
\end{exercice}

\begin{exercice}[Définitions]
Définis précisément, en deux ou trois phrases chacune, les termes suivants : \cle{industrie lourde}, \cle{agro-industrie}, \cle{PME}, \cle{zone industrielle}. Pour chaque terme, donne un exemple camerounais.
\end{exercice}

\begin{exercice}[Calculs directs]
La production électrique installée du Cameroun est d'environ \SI{1500}{MW} (ordre de grandeur), répartie ainsi : hydroélectricité 60\,\%, centrales thermiques et à gaz 38\,\%, solaire et autres sources 2\,\%.
\begin{enumerate}
\item Calcule la puissance installée en MW pour chacune des trois sources.
\item Calcule l'angle, en degrés, du secteur correspondant à l'hydroélectricité dans un diagramme circulaire.
\item Si le potentiel hydroélectrique exploitable est estimé à \SI{20000}{MW}, quel pourcentage de ce potentiel la puissance hydroélectrique installée représente-t-elle ?
\end{enumerate}
\end{exercice}

\courssub{Je m'entraîne}

\begin{exercice}[Classer les entreprises]
Classer les entreprises suivantes en trois catégories (industries lourdes, agro-industries, PME) en justifiant chaque choix par la nature de la production et la taille de l'unité : ALUCAM, SOSUCAM, CIMENCAM, une menuiserie de Bafoussam, SODECOTON, DANGOTE Cement Cameroon, une brasserie artisanale de Dschang.
\end{exercice}

\begin{exercice}[Localisation sur carte muette]
Sur une carte muette du Cameroun, localise et légende à l'aide de symboles adaptés : deux barrages hydroélectriques (Edéa et Song Loulou, ou Lagdo), la centrale à gaz de Kribi, le gisement de bauxite de Minim-Martap, le gisement de fer de Mbalam, et deux zones industrielles (Douala-Bonabéri et Yaoundé). N'oublie ni le titre, ni l'orientation, ni la légende numérotée.
\end{exercice}

\begin{exercice}[La crise et l'industrie]
Rédige un texte structuré de 15 à 20 lignes expliquant comment la crise économique des années 1980-1990 a brisé l'élan industriel du Cameroun. Tu citeras au moins deux causes de la crise et deux conséquences concrètes sur le tissu industriel (fermetures d'unités, endettement des entreprises publiques, privatisations).
\end{exercice}

\begin{exercice}[Diagramme en barres]
Le tableau ci-dessous donne la répartition approximative du nombre d'entreprises industrielles au Cameroun par type (ordres de grandeur).

\begin{center}
\begin{tabularx}{\linewidth}{|l|X|X|X|}
\hline
\rowcolor{popblueL} \textbf{Type d'industrie} & \textbf{Agro-industries} & \textbf{PME manufacturières} & \textbf{Industries lourdes} \\
\hline
Part estimée (\%) & 40 & 55 & 5 \\
\hline
\end{tabularx}
\end{center}

\begin{enumerate}
\item Construis un diagramme en barres légendé représentant cette répartition.
\item Dégage en quelques lignes deux enseignements sur la structure du tissu industriel camerounais.
\end{enumerate}
\end{exercice}

\courssub{Je me prépare au Bac}

\begin{exercice}[Étude de document : la production électrique]
\textbf{Document.} Production d'électricité au Cameroun par source (ordres de grandeur, début des années 2020 ; production totale : environ \SI{8000}{GWh} par an).

\begin{center}
\begin{tabularx}{\linewidth}{|l|X|X|}
\hline
\rowcolor{popblueL} \textbf{Source} & \textbf{Production (GWh)} & \textbf{Part (\%)} \\
\hline
Barrages hydroélectriques & 5\,600 & 70 \\
\hline
Centrales thermiques et à gaz & 2\,240 & 28 \\
\hline
Solaire et autres & 160 & 2 \\
\hline
\textbf{Total} & \textbf{8\,000} & \textbf{100} \\
\hline
\end{tabularx}
\end{center}

\begin{enumerate}
\item Construis un diagramme circulaire (camembert) légendé représentant la structure de la production électrique camerounaise. Tu calculeras au préalable l'angle de chaque secteur et tu indiqueras le matériel et la méthode utilisés.
\item Dégage deux enseignements sur la structure énergétique du Cameroun.
\item Explique pourquoi la forte dépendance à l'hydroélectricité constitue à la fois un atout et une faiblesse pour l'industrie camerounaise.
\item Cite deux projets ou réalisations récents visant à diversifier ou à renforcer la production d'énergie au Cameroun.
\end{enumerate}
\end{exercice}

\begin{exercice}[Question de synthèse]
Sujet : \emph{« Le Cameroun dispose d'atouts considérables, pourtant son industrie reste fragile. »}

À l'aide d'exemples précis (ALUCAM, DANGOTE, CIMENCAM, SODECOTON, SOSUCAM), rédige un devoir structuré (introduction avec problématique, développement en deux ou trois parties, conclusion) dans lequel tu :
\begin{enumerate}
\item présentes les atouts de l'industrie camerounaise (potentiel énergétique, réserves minières, matières premières agricoles) ;
\item décris les réalisations industrielles (types d'industries, zones industrielles) ;
\item analyses les limites qui expliquent la fragilité persistante de ce secteur.
\end{enumerate}
\end{exercice}

\begin{exercice}[Croquis : la zone industrielle de Douala]
\begin{enumerate}
\item Construis un croquis légendé de la zone industrielle de Douala (Bonabéri, Bassa, zone portuaire) en y faisant figurer : le port autonome de Douala, les principales usines, le fleuve Wouri et le pont, ainsi que les axes de communication (routes, voie ferrée vers Yaoundé et l'intérieur du pays). Ton croquis comportera un titre, une orientation et une légende numérotée.
\item À partir de ton croquis, montre en un paragraphe argumenté le rôle de cette zone industrielle dans l'intégration nationale (approvisionnement des autres régions, emplois, lien avec l'hinterland et les pays voisins enclavés comme le Tchad et la Centrafrique).
\item Propose deux mesures susceptibles de renforcer le rôle moteur de cette zone industrielle.
\end{enumerate}
\end{exercice}

\newpage

\coursec{Corrigés des exercices}

\begin{corrige}[Exercice 1 — QCM : l'industrie camerounaise]
\begin{enumerate}
\item \textbf{Réponse c : Edéa.} Le barrage d'Edéa, sur la Sanaga, est la première grande centrale hydroélectrique du Cameroun ; sa construction (années 1950) a été liée à l'usine d'aluminium ALUCAM, grande consommatrice d'électricité. Song Loulou est venu ensuite (années 1970-1980), Lagdo (1982) se trouve sur la Bénoué dans le Nord, et Memve'ele (sur la Ntem, Sud) n'a été mis en service qu'à la fin des années 2010.
\item \textbf{Réponse a : Minim-Martap et Ngaoundal (Adamaoua).} Les plus importantes réserves de bauxite camerounaise (estimées à plus d'un milliard de tonnes, ordre de grandeur) se situent sur les plateaux de l'Adamaoua. Elles ne sont pas encore exploitées industriellement.
\item \textbf{Réponse b : production d'aluminium.} ALUCAM (Compagnie camerounaise de l'aluminium), installée à Edéa, transforme l'alumine importée en aluminium grâce à l'électricité du barrage voisin. C'est l'archétype de l'industrie lourde au Cameroun.
\item \textbf{Réponse b : 20 GW.} Le potentiel hydroélectrique exploitable du Cameroun est estimé à environ \SI{20}{GW} (soit \SI{20000}{MW}), le deuxième d'Afrique après celui de la République démocratique du Congo. La puissance réellement installée n'en représente qu'une faible part.
\end{enumerate}
\textbf{Point de vigilance :} ne pas confondre \emph{potentiel} (ce qu'on pourrait produire) et \emph{puissance installée} (ce qui est réellement en service, environ \SI{1500}{MW}).
\end{corrige}

\begin{corrige}[Exercice 2 — Vrai ou faux ?]
\begin{enumerate}
\item \textbf{Faux.} Les industries lourdes (ALUCAM, CIMENCAM, DANGOTE Cement) sont \emph{minoritaires} : le tissu industriel camerounais est dominé par les agro-industries (SOSUCAM, SODECOTON, brasseries) et par une multitude de PME.
\item \textbf{Vrai.} La crise économique des années 1980-1990 (chute des cours des matières premières, endettement) a provoqué la fermeture de nombreuses unités industrielles, la faillite ou la privatisation d'entreprises publiques (ex. : CELLUCAM, CAMSUCO) et une forte désindustrialisation.
\item \textbf{Vrai.} La centrale à gaz de Kribi (environ \SI{216}{MW}) fonctionne au gaz naturel extrait du gisement offshore de Sanaga Sud, donc une énergie fossile produite sur le territoire national.
\item \textbf{Faux.} Les gisements de bauxite de Minim-Martap et de fer de Mbalam, bien que considérables, ne sont pas encore exploités industriellement : ils attendent des investissements massifs, notamment en infrastructures de transport (chemin de fer, port minéralier).
\end{enumerate}
\end{corrige}

\begin{corrige}[Exercice 3 — Définitions]
\begin{itemize}
\item \cle{Industrie lourde} : industrie qui utilise de grandes quantités de matières premières et d'énergie, exige des capitaux importants et produit des biens intermédiaires (métaux, ciment, produits chimiques). Elle emploie souvent peu de main-d'œuvre par rapport aux capitaux investis. \emph{Exemple camerounais :} ALUCAM à Edéa (aluminium) ou CIMENCAM (cimenterie).
\item \cle{Agro-industrie} : industrie qui transforme les produits agricoles, d'élevage ou de pêche en biens de consommation (sucre, huile, farine, boissons, coton égrené). Elle relie directement l'agriculture au marché. \emph{Exemple camerounais :} SOSUCAM (sucre), SODECOTON (coton), la SOCAPALM (huile de palme).
\item \cle{PME} (petite et moyenne entreprise) : entreprise de taille modeste (effectifs et chiffre d'affaires limités), souvent familiale, très nombreuse dans l'agroalimentaire, la menuiserie ou la confection. Elle crée beaucoup d'emplois mais manque de capitaux. \emph{Exemple camerounais :} une menuiserie de Bafoussam, une minoterie de quartier à Douala.
\item \cle{Zone industrielle} : espace aménagé et viabilisé (routes, eau, électricité) destiné à regrouper des usines afin de mutualiser les infrastructures et de limiter les nuisances en ville. \emph{Exemple camerounais :} la zone industrielle de Douala (Bonabéri, Bassa, zone portuaire) ou celle de Yaoundé.
\end{itemize}
\end{corrige}

\begin{corrige}[Exercice 4 — Calculs directs]
\begin{enumerate}
\item Puissance installée par source, pour un total de \SI{1500}{MW} (ordre de grandeur début 2020) :
\begin{align*}
\text{Hydroélectricité (70\,\%)} &= 1500 \times \frac{70}{100} = \SI{1050}{MW}\\
\text{Thermique et gaz (28\,\%)} &= 1500 \times \frac{28}{100} = \SI{420}{MW}\\
\text{Solaire et autres (2\,\%)} &= 1500 \times \frac{2}{100} = \SI{30}{MW}
\end{align*}
Vérification : $1050 + 420 + 30 = \SI{1500}{MW}$. Le total est correct.
\item Un cercle complet mesure $360^{\circ}$. Les angles des secteurs sont :
\begin{align*}
\text{Hydroélectricité} &: 360^{\circ} \times \frac{70}{100} = 252^{\circ}\\
\text{Thermique et gaz} &: 360^{\circ} \times \frac{28}{100} = 100{,}8^{\circ}\\
\text{Solaire et autres} &: 360^{\circ} \times \frac{2}{100} = 7{,}2^{\circ}
\end{align*}
Vérification : $252 + 100{,}8 + 7{,}2 = 360^{\circ}$.
\item Taux d'exploitation du potentiel hydroélectrique :
\[ \frac{1050}{20000} \times 100 = 5{,}25\,\%. \]
La puissance hydroélectrique installée ne représente qu'environ \textbf{5,3\,\%} du potentiel exploitable : la marge de progression est donc immense.
\end{enumerate}
\textbf{Points de vigilance :} bien convertir les pourcentages en fractions avant de multiplier ; vérifier toujours que la somme des parts redonne 100\,\% et que la somme des angles redonne $360^{\circ}$ ; utiliser \texttt{\textbackslash SI} pour les unités dans les formules.
\end{corrige}

\begin{corrige}[Exercice 5 — Classer les entreprises]
\begin{center}
\begin{tabularx}{\linewidth}{|l|X|X|}
\hline
\rowcolor{popblueL} \textbf{Catégorie} & \textbf{Entreprises} & \textbf{Justification} \\
\hline
Industries lourdes & ALUCAM, CIMENCAM, DANGOTE Cement Cameroon & Forte intensité capitalistique, grande consommation d'énergie, production de biens intermédiaires (aluminium, clinker et ciment). \\
\hline
Agro-industries & SOSUCAM, SODECOTON & Transformation de matières premières agricoles nationales : canne à sucre (SOSUCAM) et coton (SODECOTON) ; unités de grande taille adossées à des plantations. \\
\hline
PME & Menuiserie de Bafoussam, brasserie artisanale de Dschang & Petites unités, capitaux limités, main-d'œuvre locale, production destinée au marché local ou régional. \\
\hline
\end{tabularx}
\end{center}
\textbf{Point de vigilance :} une agro-industrie peut être une \emph{grande} entreprise (SOSUCAM emploie des milliers de salariés) : le critère « agro-industrie » porte sur la \emph{nature de la production} (transformation de produits agricoles), pas sur la taille. Le critère « PME » porte, lui, sur la \emph{taille}.
\end{corrige}

\begin{corrige}[Exercice 6 — Localisation sur carte muette]
\textbf{Méthode et éléments attendus sur la carte :}
\begin{itemize}
\item \textbf{Titre :} « Les principaux sites énergétiques, miniers et industriels du Cameroun » ; flèche du nord ; légende numérotée avec symboles distincts.
\item \textbf{Barrages hydroélectriques :} Edéa et Song Loulou sur la \emph{Sanaga} (région du Littoral/Centre, au sud-est de Douala) ; Lagdo sur la \emph{Bénoué} (région du Nord, près de Garoua).
\item \textbf{Centrale à gaz de Kribi :} sur la côte, région du Sud.
\item \textbf{Bauxite de Minim-Martap :} plateau de l'Adamaoua, au nord de Ngaoundéré.
\item \textbf{Fer de Mbalam :} région de l'Est, près de la frontière congolaise (non loin de Batouri).
\item \textbf{Zones industrielles :} Douala-Bonabéri (Littoral, sur le Wouri) et Yaoundé (Centre).
\end{itemize}
\textbf{Barème indicatif (sur 10) :} titre, nord, légende numérotée : 3 pts ; localisation correcte des 7 éléments : 7 pts (1 pt chacun).
\textbf{Point de vigilance :} Edéa et Song Loulou sont sur la Sanaga, \emph{pas} sur la Bénoué ; Lagdo est le seul grand barrage du Nord. Ne pas placer Mbalam dans l'Adamaoua (confusion fréquente avec Minim-Martap).
\end{corrige}

\begin{corrige}[Exercice 7 — La crise et l'industrie]
\textbf{Corrigé type (texte structuré) :}

Après l'indépendance, le Cameroun connaît un véritable élan industriel : l'État crée de nombreuses entreprises publiques (ALUCAM, SOSUCAM, SODECOTON, CELLUCAM) et le pétrole, découvert à la fin des années 1970, finance de grands projets. Mais à partir de 1985-1986, la crise économique brise nettement cette dynamique.

Les causes de la crise sont d'abord externes : l'effondrement des cours mondiaux des matières premières (cacao, café, pétrole) réduit brutalement les recettes d'exportation et le budget de l'État. Elles sont aussi internes : mauvaise gestion et endettement massif des entreprises publiques, coûts de production élevés, concurrence des produits importés.

Les conséquences sur le tissu industriel sont lourdes. De nombreuses unités ferment leurs portes : CELLUCAM (pâtes à papier) cesse toute activité, CAMSUCO disparaît, des filatures et des conserveries ferment. Les entreprises publiques survivantes, surendettées, sont restructurées ou privatisées dans le cadre des programmes d'ajustement structurel imposés par le FMI et la Banque mondiale : la SONEL et la SNEC sont cédées, ALUCAM passe sous contrôle étranger avant de péricliter. Les licenciements massifs appauvrissent les villes et le secteur informel explose.

Ainsi, en une décennie, le Cameroun perd une partie de son appareil productif : la part de l'industrie dans le PIB recule et le pays se réenferme dans l'exportation de matières premières brutes, une fragilité dont il peine encore à sortir.

\textbf{Barème indicatif (sur 10) :} introduction et problématisation : 2 pts ; deux causes citées et expliquées : 3 pts ; deux conséquences concrètes (fermetures, privatisations, exemples nommés) : 3 pts ; organisation et qualité de la langue : 2 pts.
\textbf{Points de vigilance :} dater la crise (années 1980-1990) ; citer des entreprises \emph{nommément} ; ne pas attribuer la crise à la seule mauvaise gestion (la chute des cours est le déclencheur).
\end{corrige}

\begin{corrige}[Exercice 8 — Diagramme en barres]
\begin{enumerate}
\item Diagramme en barres attendu (échelle : 1 cm pour 10\,\% par exemple) :

\begin{popfigure}
\begin{center}
\begin{tikzpicture}
\begin{axis}[
    ybar, bar width=1.1cm,
    width=0.85\linewidth, height=6.5cm,
    ymin=0, ymax=65,
    ylabel={Part estimée (\%)},
    symbolic x coords={Agro-industries, PME manufacturières, Industries lourdes},
    xtick=data,
    x tick label style={font=\small, align=center},
    nodes near coords,
    nodes near coords style={font=\small\bfseries, color=popdark},
    axis lines=left,
    ytick={0,10,20,30,40,50,60},
    grid=major, grid style={popline!40},
]
\addplot[fill=popblue, draw=popdark] coordinates {(Agro-industries,40) (PME manufacturières,55) (Industries lourdes,5)};
\end{axis}
\end{tikzpicture}
\end{center}
\legende{Répartition estimée du nombre d'entreprises industrielles du Cameroun par type (ordres de grandeur). Source : données du cours.}
\end{popfigure}

\item \textbf{Deux enseignements :}
\begin{itemize}
\item Le tissu industriel camerounais est \emph{dominé par les petites unités} : les PME manufacturières représentent à elles seules plus de la moitié des entreprises (55\,\%), ce qui traduit un secteur fragmenté, peu capitalistique mais créateur d'emplois.
\item Les \emph{industries lourdes sont marginales en nombre} (5\,\%) malgré leur poids économique et symbolique (ALUCAM, CIMENCAM, DANGOTE) ; les agro-industries (40\,\%) confirment l'ancrage de l'industrie camerounaise dans la transformation des produits agricoles.
\end{itemize}
\end{enumerate}
\textbf{Barème indicatif (sur 10) :} axes gradués et légendés, titre : 3 pts ; hauteurs exactes des trois barres : 3 pts ; deux enseignements pertinents et rédigés : 4 pts.
\textbf{Point de vigilance :} les barres doivent être de \emph{même largeur} et séparées ; ne pas oublier l'unité (\%) sur l'axe vertical.
\end{corrige}

\begin{corrige}[Exercice 9 — Étude de document : la production électrique]
\begin{enumerate}
\item \textbf{Calcul préalable des angles} (règle de trois : angle $=$ part $\times 360^{\circ}/100$) :
\begin{align*}
\text{Barrages hydroélectriques} &: 70 \times 3{,}6 = 252^{\circ}\\
\text{Centrales thermiques et à gaz} &: 28 \times 3{,}6 = 100{,}8^{\circ}\\
\text{Solaire et autres} &: 2 \times 3{,}6 = 7{,}2^{\circ}
\end{align*}
Vérification : $252 + 100{,}8 + 7{,}2 = 360^{\circ}$. Le calcul est exact.

\textbf{Matériel :} compas, rapporteur, règle, crayons de couleur. \textbf{Méthode :} tracer un cercle ; tracer un rayon de référence (vers le haut) ; reporter au rapporteur $252^{\circ}$ pour l'hydroélectricité, puis $100{,}8^{\circ}$ pour le thermique/gaz ; le secteur restant ($7{,}2^{\circ}$) correspond au solaire ; colorier, légender, titrer.

\begin{popfigure}
\begin{center}
\begin{tikzpicture}[scale=1.6]
% Secteur hydro : de 90° à 90-252 = -162° (sens horaire)
\fill[popblue] (0,0) -- (90:1.5) arc (90:-162:1.5) -- cycle;
% Secteur thermique : de -162° à -162-100.8 = -262.8°
\fill[poporange] (0,0) -- (-162:1.5) arc (-162:-262.8:1.5) -- cycle;
% Secteur solaire : de -262.8° à -270° (= 90-360)
\fill[popgold] (0,0) -- (-262.8:1.5) arc (-262.8:-270:1.5) -- cycle;
\draw[popdark, thick] (0,0) circle (1.5);
% Étiquettes
\node[font=\small\bfseries, text=white] at (-36:0.9) {70\,\%};
\node[font=\small\bfseries, text=white] at (145:0.9) {28\,\%};
\node[font=\scriptsize\bfseries, text=popdark] at (95:0.9) {2\,\%};
% Légende
\node[anchor=west, font=\small] at (2.0,0.6) {\textcolor{popblue}{\rule{0.35cm}{0.25cm}} Barrages hydroélectriques ($252^{\circ}$)};
\node[anchor=west, font=\small] at (2.0,0.0) {\textcolor{poporange}{\rule{0.35cm}{0.25cm}} Centrales thermiques et à gaz ($100{,}8^{\circ}$)};
\node[anchor=west, font=\small] at (2.0,-0.6) {\textcolor{popgold}{\rule{0.35cm}{0.25cm}} Solaire et autres ($7{,}2^{\circ}$)};
\end{tikzpicture}
\end{center}
\legende{Structure de la production électrique du Cameroun par source (ordres de grandeur, début des années 2020 ; total : environ \SI{8000}{GWh}/an). Source : données du cours.}
\end{popfigure}

\item \textbf{Deux enseignements :}
\begin{itemize}
\item La production électrique camerounaise repose \emph{très largement sur l'hydroélectricité} (70\,\%), grâce aux barrages de la Sanaga (Edéa, Song Loulou) et de la Bénoué (Lagdo) : c'est une énergie renouvelable et peu coûteuse à l'usage.
\item Les autres sources restent \emph{marginales} : le thermique et le gaz (28\,\%) servent surtout d'appoint, et le solaire (2\,\%) est embryonnaire malgré un ensoleillement exceptionnel, surtout dans le Nord.
\end{itemize}

\item \textbf{Atout et faiblesse de la dépendance hydraulique.} \emph{Atout :} l'hydroélectricité est une énergie renouvelable, relativement bon marché, qui a permis l'installation d'industries électro-intensives comme ALUCAM à Edéa et qui offre un énorme potentiel inexploité (environ \SI{20}{GW}). \emph{Faiblesse :} en saison sèche, la baisse du niveau des retenues (notamment Lagdo et Song Loulou) provoque des délestages qui pénalisent les usines ; de plus, la concentration des barrages sur un seul fleuve (la Sanaga) rend le système vulnérable à tout incident. D'où des coûts supplémentaires pour les industriels (groupes électrogènes) et une compétitivité réduite.

\item \textbf{Deux projets ou réalisations récents :}
\begin{itemize}
\item Le \textbf{barrage de Memve'ele} (environ \SI{200}{MW}) sur la Ntem (région du Sud), mis en service à la fin des années 2010 ;
\item la \textbf{centrale à gaz de Kribi} (environ \SI{216}{MW}, extensible), qui diversifie le mix énergétique ;
\item on peut aussi citer le \textbf{barrage de Nachtigal} (environ \SI{420}{MW}) sur la Sanaga, en construction, ou les centrales solaires de \textbf{Guider} et \textbf{Maroua} dans le Grand Nord.
\end{itemize}
\end{enumerate}
\textbf{Barème indicatif (sur 20) :} calculs d'angles exacts et vérifiés : 4 pts ; camembert soigné (cercle, rapporteur, couleurs, légende, titre) : 5 pts ; deux enseignements : 4 pts ; analyse atout/faiblesse argumentée : 4 pts ; deux projets correctement nommés : 3 pts.
\textbf{Points de vigilance :} exiger le \emph{calcul préalable} des angles (un camembert sans calculs n'est pas recevable) ; citer le matériel (compas, rapporteur) ; vérifier que la somme des angles fait $360^{\circ}$ ; pour la question 3, traiter \emph{les deux} aspects (atout ET faiblesse).
\end{corrige}

\begin{corrige}[Exercice 10 — Question de synthèse]
\textbf{Corrigé type (devoir structuré) :}

\emph{Introduction.} Le Cameroun est souvent présenté comme « l'Afrique en miniature » : son sous-sol recèle d'importantes réserves minières, son réseau hydrographique offre le deuxième potentiel hydroélectrique d'Afrique et son agriculture fournit d'abondantes matières premières. Pourtant, plus de soixante ans après l'indépendance, l'industrie camerounaise reste fragile et contribue modestement au PIB. On peut alors se demander : comment expliquer le décalage entre des atouts considérables et des réalisations industrielles limitées ? Nous présenterons d'abord les atouts, puis les réalisations, avant d'analyser les limites persistantes.

\emph{I. Des atouts considérables.} Le Cameroun dispose d'un énorme potentiel énergétique : potentiel hydroélectrique estimé à environ \SI{20}{GW} (deuxième d'Afrique), gaz naturel exploité à Sanaga Sud et Kribi, ensoleillement exceptionnel dans le Grand Nord. Ses réserves minières sont importantes : bauxite de Minim-Martap et Ngaoundal (plus d'un milliard de tonnes, ordre de grandeur), fer de Mbalam, rutile, terres rares à Akonolinga. S'y ajoutent des matières premières agricoles abondantes : coton, cacao, café, canne à sucre, palme, hévéa, qui alimentent les agro-industries.

\emph{II. Des réalisations réelles mais inégales.} L'appareil industriel repose sur quelques industries lourdes symboliques : ALUCAM (aluminium à Edéa), CIMENCAM et DANGOTE Cement Cameroon (cimenteries à Douala et Nomayos). Mais il est dominé par les agro-industries — SOSUCAM (sucre), SODECOTON (coton), SOCAPALM (huile de palme), les brasseries — et par une myriade de PME. Ces unités se concentrent dans les zones industrielles de Douala (Bonabéri, Bassa, zone portuaire) et de Yaoundé, ce qui traduit une forte polarisation littorale. La production électrique (environ \SI{8000}{GWh}/an) repose à 70\,\% sur les barrages d'Edéa, Song Loulou et Lagdo.

\emph{III. Une fragilité persistante.} Plusieurs limites expliquent cette fragilité. D'abord, l'héritage de la crise des années 1980-1990 : fermetures d'unités (CELLUCAM, CAMSUCO), privatisations douloureuses, désindustrialisation. Ensuite, les contraintes structurelles : énergie insuffisante et délestages, infrastructures de transport défaillantes (routes dégradées, chemin de fer vieillissant), qui bloquent l'exploitation des gisements de Minim-Martap et de Mbalam. Enfin, la concurrence des importations asiatiques, le faible accès au financement des PME et la transformation limitée des matières premières sur place (cacao et coton exportés largement bruts) maintiennent le pays dans une position de fournisseur de produits primaires.

\emph{Conclusion.} Le Cameroun possède donc les cartes de l'industrialisation — énergie, minerais, agriculture — mais ne les a pas encore toutes jouées : les projets récents (Memve'ele, Nachtigal, centrale à gaz de Kribi, cimenteries nouvelles) montrent une relance, qui suppose des infrastructures, des capitaux et une volonté de transformation locale pour devenir, comme l'y invite la vision « Cameroun émergent 2035 », un véritable pays industriel.

\textbf{Barème indicatif (sur 20) :} introduction avec problématique : 3 pts ; atouts illustrés d'exemples : 4 pts ; réalisations (types d'industries, zones industrielles, entreprises nommées) : 5 pts ; analyse des limites : 5 pts ; conclusion ouverte, cohérence et langue : 3 pts.
\textbf{Points de vigilance :} annoncer le plan dans l'introduction ; nommer \emph{au moins} les cinq entreprises du sujet ; éviter le catalogue : chaque exemple doit être \emph{exploité} (que prouve-t-il ?) ; ne pas confondre atouts (potentiels) et réalisations (existant).
\end{corrige}

\begin{corrige}[Exercice 11 — Croquis : la zone industrielle de Douala]
\begin{enumerate}
\item \textbf{Croquis attendu} (modèle de rédaction ; l'élève le reproduit à main levée sur sa copie) :

\begin{popfigure}
\begin{center}
\includegraphics[width=\linewidth,height=11cm,keepaspectratio]{N08/douala_map.jpg}
\end{center}
\legende{Plan de Douala (fond OpenStreetMap, nord en haut) servant de modèle pour le croquis : le Wouri, le pont de la route N3 vers Bonabéri, le port au bord du fleuve côté centre-ville, la zone industrielle de Bassa à l'est, la voie ferrée avec la gare centrale, la route N3 vers Édéa et Yaoundé, l'aéroport. Le croquis attendu en simplifie ces éléments : 1. port autonome ; 2. axe routier vers Yaoundé et l'intérieur ; 3. voie ferrée ; usines à figurer par des points. \\ Source : Wikimedia Commons, OpenStreetMap contributors, CC BY-SA 3.0.}
\end{popfigure}

\textbf{Conseils de construction :} tracer d'abord le fleuve Wouri et le pont (éléments structurants) ; placer le port sur la rive gauche, au pied du centre-ville ; situer la zone de Bassa au sud-est et Bonabéri sur la rive droite, au-delà du pont ; figurer les usines par des points ou des carrés ; terminer par le titre, la flèche du nord et la légende numérotée.

\item \textbf{Paragraphe argumenté (rôle dans l'intégration nationale).} La zone industrielle de Douala est le poumon économique du Cameroun : adossée au premier port du pays, qui traite l'essentiel du commerce extérieur national, elle concentre cimenteries (CIMENCAM, DANGOTE), brûleries de cacao, brasseries et agro-industries. Par la route et le chemin de fer (ligne Douala--Yaoundé--Ngaoundéré), ses produits — ciment, boissons, huiles, savons — irriguent toutes les régions, tandis qu'elle offre des dizaines de milliers d'emplois qui attirent les migrants de l'intérieur. Elle sert enfin de débouché maritime aux pays voisins enclavés : le Tchad et la Centrafrique transitent par le corridor de Douala. Elle tisse ainsi un lien matériel entre les régions camerounaises et entre le Cameroun et son hinterland continental.

\item \textbf{Deux mesures pour renforcer ce rôle moteur :}
\begin{itemize}
\item \emph{Améliorer les infrastructures :} moderniser et draguer le port, réhabiliter la voie ferrée et les routes vers Yaoundé et le Grand Nord, garantir une fourniture électrique stable pour supprimer les délestages qui freinent les usines ;
\item \emph{Déconcentrer et sécuriser l'investissement :} créer des zones économiques spéciales avec des incitations fiscales, développer la transformation locale des matières premières (cacao, coton, bois) et faciliter le financement des PME pour densifier le tissu industriel.
\end{itemize}
\end{enumerate}
\textbf{Barème indicatif (sur 20) :} croquis (titre, nord, légende numérotée, éléments correctement placés) : 8 pts ; paragraphe argumenté (au moins trois fonctions : approvisionnement, emplois, hinterland) : 8 pts ; deux mesures réalistes et expliquées : 4 pts.
\textbf{Points de vigilance :} le port est sur la \emph{rive gauche} du Wouri, Bonabéri sur la rive droite (au-delà du pont) ; ne pas oublier la légende \emph{numérotée} ; dans le paragraphe, citer explicitement le Tchad et la Centrafrique comme pays enclavés desservis par le corridor de Douala.
\end{corrige}

\newpage

\coursec{Fiche bilan}

\begin{popfigure}
\begin{tikzpicture}[mindmap, grow cyclic, every node/.style={concept}, concept color=popblue!70, text=white, root concept/.append style={font=\bfseries\small}, level 1/.append style={level distance=3.4cm, sibling angle=51.4, font=\scriptsize}, level 2/.append style={level distance=1.9cm, sibling angle=40, font=\tiny}]
\node[root concept]{L'INDUSTRIE AU CAMEROUN}
  child[concept color=poporange!85]{ node{Définitions}
    child{ node{Industrie, PME} }
    child{ node{Zone industrielle} }
  }
  child[concept color=popgold!90]{ node{Potentiel énergétique}
    child{ node{Barrages hydro} }
    child{ node{Gaz, thermique, solaire} }
  }
  child[concept color=popgreen!80]{ node{Réserves minières}
    child{ node{Bauxite, fer} }
    child{ node{Terres rares} }
  }
  child[concept color=poppink!85]{ node{Crise industrielle}
    child{ node{Fermetures d'unités} }
    child{ node{Élan brisé (années 1980-90)} }
  }
  child[concept color=poppurple!85]{ node{Types d'industries}
    child{ node{Lourdes : ALUCAM, CIMENCAM, DANGOTE} }
    child{ node{Agro-industries et PME} }
  }
  child[concept color=popblue!60]{ node{Zones industrielles}
    child{ node{Douala-Bonabéri} }
    child{ node{Yaoundé, Kribi} }
  }
  child[concept color=popmuted!90]{ node{TP 3 : camembert}
    child{ node{Angle = part $\times$ 3,6} }
  };
\end{tikzpicture}
\legende{Carte mentale de la leçon : les sept branches essentielles de l'industrie camerounaise (définitions, énergie, mines, crise, types d'industries, zones industrielles, TP 3).}
\end{popfigure}

\begin{aretenir}[L'essentiel de la leçon]
\textbf{Définitions clés}
\begin{itemize}
  \item \cle{Industrie} : ensemble des activités économiques qui transforment les matières premières (agricoles, minières, énergétiques) en biens de consommation ou en biens d'équipement.
  \item \cle{Industrie lourde} : industrie de base utilisant d'importants capitaux et beaucoup d'énergie (métallurgie, cimenterie, chimie) ; au Cameroun : \textbf{ALUCAM} (aluminium, Edéa), \textbf{CIMENCAM} (ciment, Douala-Bonabéri et Nomayos) et \textbf{DANGOTE Cement} (ciment, Douala).
  \item \cle{Agro-industrie} : industrie qui transforme les produits agricoles (huileries, savonneries, sucreries, brasseries) ; elle domine le tissu industriel camerounais avec les \cle{PME} (petites et moyennes entreprises).
  \item \cle{Zone industrielle} : espace aménagé et viabilisé (eau, électricité, routes) regroupant plusieurs unités industrielles (ex. : Bonabéri à Douala, Yaoundé, zone portuaire de Kribi).
\end{itemize}

\textbf{Données et repères à connaître par cœur}
\begin{center}
\begin{tabularx}{\linewidth}{|l|X|}\hline
\rowcolor{popblueL} \textbf{Repère} & \textbf{Contenu chiffré ou nominal} \\ \hline
Potentiel hydroélectrique & Environ \SI{20000}{MW} exploitables (ordre de grandeur ; certaines estimations du potentiel total vont bien au-delà), l'un des tout premiers d'Afrique ; barrages de Songloulou et Edéa (Sanaga), Lagdo (Bénoué), Memve'ele (environ \SI{211}{MW}) et Nachtigal (environ \SI{420}{MW}, récents). \\ \hline
Énergie gazière et thermique & Centrales à gaz de Kribi (environ \SI{216}{MW}) et Logbaba à Douala, alimentées par le gaz de Sanaga Sud ; centrales thermiques de secours ; solaire encore marginal mais en progression (Nord). \\ \hline
Réserves minières & Bauxite (Minim-Martap et Ngaoundal, Adamaoua), fer (Mbalam, à l'Est, et gisements de la région de Kribi), nickel et cobalt (Lomié), rutile (Akonolinga), terres rares. \\ \hline
Crise industrielle & À partir de 1986-1987 : baisse des cours des matières premières, ajustements structurels, fermeture de nombreuses unités (textile, huileries, montage), désindustrialisation partielle. \\ \hline
Structure industrielle & Industries lourdes minoritaires ; prédominance des agro-industries et des PME ; Douala concentre la majorité des unités industrielles du pays. \\ \hline
\end{tabularx}
\end{center}

\textbf{Méthodes express}
\begin{itemize}
  \item \cle{Camembert (TP 3)} : angle de chaque secteur $= \dfrac{\text{part en \%}}{100}\times 360$ (soit part $\times$ 3,6) ; reporter au rapporteur, colorier, légender, donner un titre et la source.
  \item \cle{Croquis industriel} : contour simplifié, villes (points nommés), barrages (symbole), zones industrielles (surfaces), flux (flèches), nord et échelle indicative.
  \item \cle{Commentaire de tableau} : dégager les grandes masses, comparer (plus/moins), expliquer par les facteurs (énergie, matières premières, marché, crise).
\end{itemize}
\end{aretenir}

\begin{attention}[Précautions sur les chiffres]
Les données industrielles et énergétiques évoluent vite (mise en service de Nachtigal, montée en puissance de Memve'ele, projets miniers de Mbalam et de Minim-Martap encore en attente d'exploitation effective). Au Bac, retiens les \emph{ordres de grandeur} et les \emph{localisations} plutôt que des chiffres au MW près, et signale toujours la source et l'année des données utilisées dans un document.
\end{attention}

\begin{aretenir}[Checklist avant le Bac]
\begin{itemize}
  \item $\square$ Je sais définir industrie, industrie lourde, agro-industrie, PME et zone industrielle.
  \item $\square$ Je connais l'ordre de grandeur du potentiel hydroélectrique camerounais ($\approx$ \SI{20000}{MW} exploitables).
  \item $\square$ Je cite au moins trois barrages : Songloulou, Edéa, Lagdo (et Memve'ele, Nachtigal).
  \item $\square$ Je cite les centrales à gaz de Kribi et de Logbaba (Douala).
  \item $\square$ Je localise les principales réserves : bauxite de Minim-Martap et Ngaoundal, fer de Mbalam.
  \item $\square$ J'explique la crise industrielle des années 1980-1990 et ses conséquences (fermetures d'unités, désindustrialisation partielle).
  \item $\square$ Je cite les trois industries lourdes : ALUCAM, CIMENCAM, DANGOTE Cement.
  \item $\square$ Je sais que les agro-industries et les PME dominent le tissu industriel.
  \item $\square$ Je localise les grandes zones industrielles : Douala-Bonabéri, Yaoundé, Kribi.
  \item $\square$ Je maîtrise la formule du camembert : angle $=$ part en \% $\times$ 3,6.
  \item $\square$ Je sais légender un croquis (titre, nord, échelle, légende numérotée).
  \item $\square$ Je relie l'industrie à l'intégration nationale (emplois, désenclavement, transformation locale des matières premières).
\end{itemize}
\end{aretenir}

\findecours

\end{document}
